% Generated mechanically from the frozen formal-spaces.tex target.
% Do not edit this generated file; edit the bound locale source instead.

% OK, start here.
%

\setcounter{chapter}{86}
\chapter{形式代数空间}



\phantomsection
\label{formal-spaces-section-phantom}


\section{引言}
\label{formal-spaces-section-introduction}

\noindent
形式概形由 \cite{EGA} 引入。\cite{McQuillan} 引入了形式概形的一个更一般的版本，
\cite{Yasuda} 又引入了另一个版本。形式代数空间由 \cite{Kn} 引入。
相关材料以及更多内容可见 \cite{Abbes} 和 \cite{Fujiwara-Kato}。
本章介绍我们将采用的形式代数空间概念。我们的定义足够一般，
可以把文献中大多数类别的形式概形/空间作为全子范畴包括进来。

\medskip\noindent
虽然我们确实讨论了其中若干替代理论与本理论的比较，
但在理论的逻辑发展并不需要时，我们并不总是给出全部细节。

\medskip\noindent
除引入形式代数空间外，我们还证明若干非常基本的性质，并讨论几类态射。









\section{EGA 意义下的形式概形}
\label{formal-spaces-section-formal-schemes-EGA}

\noindent
本节回顾 \cite{EGA} 中形式概形的构造。这个概念虽然在代数几何中非常有用，
却未必总是应当考虑的正确概念。或许更准确的说法是：建立该理论时作出了若干选择，
而对于不同目的，其他选择可能更合适。事实上，文献中可以找到许多彼此密切相关、
并针对作者所考虑问题而调整的不同理论。不过，这里所概述的理论有一个主要优点：
我们处理的是确定的几何对象。

\medskip\noindent
开始之前，我们应指出拓扑环层，或者更一般地拓扑空间层的层条件所涉及的一个问题。
具体而言，下列大范畴
\begin{enumerate}
\item 拓扑空间范畴，
\item 拓扑群范畴，
\item 拓扑环范畴，
\item 给定拓扑环上的拓扑模范畴，
\end{enumerate}
连同它们到 $\textit{Sets}$ 的自然遗忘函子，并不是
《层》之第 \ref{sheaves-section-algebraic-structures} 节所定义的代数结构类型的实例。
因此，我们不能贸然对它们应用该章建立的机制。另一方面，上列每个范畴都有极限和等化子，
并且到集合、群、环、模的遗忘函子与这些构造交换
（见《拓扑》之引理 \ref{topology-lemma-limits}、
\ref{topology-lemma-topological-group-limits}、
\ref{topology-lemma-topological-ring-limits} 和
\ref{topology-lemma-topological-module-limits}）。
因此，我们可以像《层》之定义
\ref{sheaves-definition-sheaf-values-in-category} 那样定义层的概念，
而其底层的集合、群、环或模预层确实是层。关键区别在于：对于开覆盖
$U = \bigcup_{i \in I} U_i$，图表
$$
\xymatrix{
\mathcal{F}(U) \ar[r]
&
\prod\nolimits_{i\in I}
\mathcal{F}(U_i)
\ar@<1ex>[r] \ar@<-1ex>[r]
&
\prod\nolimits_{(i_0, i_1) \in I \times I}
\mathcal{F}(U_{i_0} \cap U_{i_1})
}
$$
必须是拓扑空间、拓扑群、拓扑环或拓扑模范畴中的等化子图表；
也就是说，第一个映射把 $\mathcal{F}(U)$ 识别为
$\prod_{i \in I} \mathcal{F}(U_i)$ 的一个子空间，
而后者带有乘积拓扑。

\medskip\noindent
拓扑空间、拓扑群、拓扑环或拓扑模层 $\mathcal{F}$ 在点 $x \in X$ 处的茎
$\mathcal{F}_x$，定义为相应范畴中对开邻域所取的余极限
$$
\mathcal{F}_x = \colim_{x\in U} \mathcal{F}(U)
$$
这等同于在集合、群、环或模层面取余极限
（见《拓扑》之引理 \ref{topology-lemma-colimits}、
\ref{topology-lemma-topological-group-colimits}、
\ref{topology-lemma-topological-ring-colimits} 和
\ref{topology-lemma-topological-module-colimits}），
但它还带有一个拓扑。警告：所得拓扑取决于所处的范畴，见
《例》之第 \ref{examples-section-colimit-topology} 节。
我们可以对拓扑空间、拓扑群、拓扑环或拓扑模的预层作层化，
而取茎与此操作交换，见注 \ref{formal-spaces-remark-sheafification-of-presheaves-in-top}。

\medskip\noindent
设 $f : X \to Y$ 是拓扑空间的连续映射。从拓扑空间、拓扑群、拓扑环或拓扑模层的范畴，
到 $Y$ 上相应层范畴有一个函子 $f_*$，照常由
$f_*\mathcal{F}(V) = \mathcal{F}(f^{-1}V)$ 定义。
（这个情形中的拉回留待后文讨论。）我们定义 $Y$ 上拓扑空间层 $\mathcal{G}$
与 $X$ 上拓扑空间层 $\mathcal{F}$ 之间的 $f$-映射
$\xi : \mathcal{G} \to \mathcal{F}$，其方式与《层》之定义
\ref{sheaves-definition-f-map} 完全相同，但另加条件：
对每个开集 $V \subset Y$，
$\xi_V : \mathcal{G}(V) \to \mathcal{F}(f^{-1}V)$ 都连续。与
《层》之引理 \ref{sheaves-lemma-f-map} 一样，有
$$
\{\mathcal{G}\text{ 到 }\mathcal{F}\text{ 的 }f\text{-映射}\} =
\Mor_{\Sh(Y, \textit{Top})}(\mathcal{G}, f_*\mathcal{F})
$$
拓扑群、拓扑环和拓扑模层的情形也相同。最后，设
$\xi : \mathcal{G} \to \mathcal{F}$ 是上述 $f$-映射。
给定 $x \in X$，令其像为 $y = f(x)$，则有茎之间的连续映射
$$
\xi_x : \mathcal{G}_y \longrightarrow \mathcal{F}_x
$$
其定义方式与《层》之定义
\ref{sheaves-definition-composition-f-maps} 之后的讨论完全相同。

\medskip\noindent
利用上述讨论，我们可以定义“局部拓扑赋环空间”的范畴 $LTRS$。
其对象是一个二元组 $(X, \mathcal{O}_X)$，其中 $X$ 是拓扑空间，
$\mathcal{O}_X$ 是拓扑环层，并且各茎 $\mathcal{O}_{X, x}$
（忘掉拓扑后）都是局部环。$LTRS$ 中从
$(X, \mathcal{O}_X)$ 到 $(Y, \mathcal{O}_Y)$ 的态射是二元组
$(f, f^\sharp)$，其中 $f : X \to Y$ 是拓扑空间的连续映射，
$f^\sharp : \mathcal{O}_Y \to \mathcal{O}_X$ 是一个 $f$-映射，
并且对每个 $x \in X$，诱导映射
$$
f^\sharp_x : \mathcal{O}_{Y, f(x)} \longrightarrow \mathcal{O}_{X, x}
$$
（忘掉拓扑后）是局部环的局部同态。复合的定义方式与局部赋环空间态射的复合完全相同。

\medskip\noindent
现在假设拓扑空间 $X$ 有一个由拟紧开集组成的基。给定集合、群、环或某环上模的层
$\mathcal{F}$，可以赋予 $\mathcal{F}$ 以拓扑空间、拓扑群、拓扑环或拓扑模层的结构。
具体而言，若 $U \subset X$ 是拟紧开集，就给 $\mathcal{F}(U)$ 赋离散拓扑。
若 $U \subset X$ 任意，则选择由拟紧开集组成的开覆盖
$U = \bigcup_{i \in I} U_i$，并给 $\mathcal{F}(U)$ 赋予从
$\prod_{i \in I} \mathcal{F}(U_i)$ 诱导的拓扑
（按照上述讨论，本来就应如此）。读者可验证（此处略去）按这种方式确实得到拓扑空间、
拓扑群、拓扑环或拓扑模层。若对每个拟紧开集 $U \subset X$，
$\mathcal{F}(U)$ 上的拓扑都是离散的，就称一个拓扑空间、拓扑群、拓扑环或拓扑模层
为{\it 伪离散的}。上述构造是遗忘函子的一个伴随，并诱导集合层范畴与伪离散拓扑空间层范畴之间的等价
（群、环、模的情形同样如此）。

\medskip\noindent
Grothendieck 和 Dieudonn\'e 首先定义了仿射形式概形。
它们对应于可容许拓扑环 $A$，见《代数续篇》之定义
\ref{more-algebra-definition-topological-ring}。具体而言，给定 $A$，考虑环 $A$
的一族定义理想基本系 $I_\lambda$。（在任意可容许拓扑环中，所有定义理想组成的族都是基本系。）
对每个 $\lambda$，可以考虑概形 $\Spec(A/I_\lambda)$。若
$I_\lambda \subset I_\mu$，则诱导态射
$$
\Spec(A/I_\mu) \to \Spec(A/I_\lambda)
$$
是一个增厚，因为对某个 $n$ 有 $I_\mu^n \subset I_\lambda$。
另一种看法是注意到每个映射
$$
\Spec(A/I_\lambda) \to \Spec(A)
$$
的像都同胚于 $A$ 的开素理想集合。这促使我们定义
$$
\text{Spf}(A) = \{\text{开素理想 }\mathfrak p \subset A\}
$$
并赋予它来自 $\Spec(A)$ 的拓扑。对每个 $\lambda$，可以把结构层
$\mathcal{O}_{\Spec(A/I_\lambda)}$ 看作 $\text{Spf}(A)$ 上的层。
令 $\mathcal{O}_\lambda$ 为上文所述相应的伪离散拓扑环层。然后令
$$
\mathcal{O}_{\text{Spf}(A)} = \lim \mathcal{O}_\lambda
$$
其中极限在拓扑环层范畴中取。二元组
$(\text{Spf}(A), \mathcal{O}_{\text{Spf}(A)})$ 称为 $A$ 的{\it 形式谱}。

\medskip\noindent
此时应当核验几件事。第一，茎 $\mathcal{O}_{\text{Spf}(A), x}$
（忘掉拓扑后）是局部环。第二，给定 $f \in A$，对相应开集
$D(f) \cap \text{Spf}(A)$，有拓扑环的等式
$$
\Gamma(D(f) \cap \text{Spf}(A), \mathcal{O}_{\text{Spf}(A)})
= A_{\{f\}} = \lim (A/I_\lambda)_f
$$
其中 $I_\lambda$ 是上述定义理想基本系。此外，环 $A_{\{f\}}$ 也是可容许的，并且
% P10-E0290
$(\text{Spf}(A_f), \mathcal{O}_{\text{Spf}(A_{\{f\}})})$
同构于
$(D(f) \cap \text{Spf}(A),
\mathcal{O}_{\text{Spf}(A)}|_{D(f) \cap \text{Spf}(A)})$。
最后，给定一对可容许拓扑环 $A,B$，有
\begin{equation}
\label{formal-spaces-equation-morphisms-affine-formal-schemes}
\Mor_{LTRS}((\text{Spf}(B), \mathcal{O}_{\text{Spf}(B)}),
(\text{Spf}(A), \mathcal{O}_{\text{Spf}(A)}))
= \Hom_{cont}(A, B)
\end{equation}
其中 $LTRS$ 是上文定义的“局部拓扑赋环空间”范畴。

\medskip\noindent
在此基础上，\cite{EGA} 把{\it 形式概形}定义为二元组
$(\mathfrak X, \mathcal{O}_\mathfrak X)$，其中 $\mathfrak X$ 是拓扑空间，
$\mathcal{O}_\mathfrak X$ 是拓扑环层，并且每个点都有一个在 $LTRS$ 中
同构于某个仿射形式概形的开邻域。形式概形的{\it 态射}
$f : (\mathfrak X, \mathcal{O}_\mathfrak X) \to
(\mathfrak Y, \mathcal{O}_\mathfrak Y)$ 是范畴 $LTRS$ 中的态射。

\medskip\noindent
设 $A$ 是带离散拓扑的环。则 $A$ 是可容许的，并且形式概形
$\text{Spf}(A)$ 等于 $\Spec(A)$。结构层
$\mathcal{O}_{\text{Spf}(A)}$ 是与 $\mathcal{O}_{\Spec(A)}$ 相伴的伪离散拓扑环层；
换言之，其底层环层等于 $\mathcal{O}_{\Spec(A)}$，并且在拟紧开集 $U$ 上，
环 $\mathcal{O}_{\text{Spf}(A)}(U) = \mathcal{O}_{\Spec(A)}(U)$
带离散拓扑，但一般的开集上并非如此。因此，可以给每个仿射概形关联一个仿射形式概形。
完全类似地，可以从一般概形 $(X, \mathcal{O}_X)$ 出发，把它关联到
$(X, \mathcal{O}'_X)$，其中 $\mathcal{O}'_X$ 是底层环层为
$\mathcal{O}_X$ 的伪离散拓扑环层。此构造与态射相容，并定义一个函子
\begin{equation}
\label{formal-spaces-equation-compare-schemes-formal-schemes}
\textit{概形} \longrightarrow \textit{形式概形}
\end{equation}
由式 (\ref{formal-spaces-equation-morphisms-affine-formal-schemes}) 立即可知，这个函子全忠实。

\medskip\noindent
设 $\mathfrak X$ 是形式概形。我们用公式
$\text{size}(\mathfrak X) = \max(\aleph_0, \kappa_1, \kappa_2)$
定义形式概形的{\it 大小}，其中 $\kappa_1$ 是 $\mathfrak X$ 的仿射形式开集的基数，
$\kappa_2$ 是所有 $\mathcal{O}_\mathfrak X(\mathfrak U)$ 的基数的上确界，
这里 $\mathfrak U \subset \mathfrak X$ 取遍这类仿射形式开集。

\begin{lemma}
\label{formal-spaces-lemma-fully-faithful}
如《集合》之引理 \ref{sets-lemma-construct-category} 选取一个概形范畴
$\Sch_\alpha$。给定形式概形 $\mathfrak X$，令
$$
h_\mathfrak X : (\Sch_\alpha)^{opp} \longrightarrow \textit{Sets},\quad
h_\mathfrak X(S) = \Mor_{\textit{形式概形}}(S, \mathfrak X)
$$
为其点函子。若 $\mathfrak X$ 的大小不过大，则
$$
\Mor_{\textit{形式概形}}(\mathfrak X, \mathfrak Y) =
\Mor_{\textit{PSh}(\Sch_\alpha)}(h_\mathfrak X, h_\mathfrak Y)
$$
\end{lemma}

\begin{proof}
首先注意，对任意形式概形 $\mathfrak X$，$h_\mathfrak X$ 都满足 Zariski 拓扑的层条件
（见《概形》之定义 \ref{schemes-definition-representable-by-open-immersions}）。
这是形式概形范畴中态射的局部性所致。此外，对形式概形的开浸入
$\mathfrak V \to \mathfrak W$，相应的函子变换
$h_\mathfrak V \to h_\mathfrak W$ 是单射且可由开浸入表示
（见《概形》之定义 \ref{schemes-definition-representable-by-open-immersions}）。
选取形式概形由仿射形式概形组成的开覆盖
$\mathfrak X = \bigcup \mathfrak U_i$。则函子族 $h_{\mathfrak U_i}$ 覆盖
$h_\mathfrak X$（见《概形》之定义
\ref{schemes-definition-representable-by-open-immersions}）。最后注意
$$
h_{\mathfrak U_i} \times_{h_\mathfrak X} h_{\mathfrak U_j} =
h_{\mathfrak U_i \cap \mathfrak U_j}
$$
所以，给出一个映射 $h_\mathfrak X \to h_\mathfrak Y$，等价于给出一族在交叠处一致的映射
$h_{\mathfrak U_i} \to h_\mathfrak Y$。因此，可以把引理中映射的双射性（分别地，单射性）
归约为各对 $(\mathfrak U_i, \mathfrak Y)$ 的双射性（分别地，单射性），
以及各对 $(\mathfrak U_i \cap \mathfrak U_j, \mathfrak Y)$ 的单射性（分别地，无需条件）。
这样便归约到 $\mathfrak X$ 是仿射形式概形的情形。设
$\mathfrak X = \text{Spf}(A)$，其中 $A$ 是某个可容许拓扑环；
再选取 $A$ 的一族定义理想基本系 $I_\lambda \subset A$。

\medskip\noindent
我们也可以在 $\mathfrak Y$ 上局部化。事实上，设
$\mathfrak V \subset \mathfrak Y$ 是开形式子概形，并设
$\varphi : h_\mathfrak X \to h_\mathfrak Y$。则
$$
h_\mathfrak V \times_{h_\mathfrak Y, \varphi} h_\mathfrak X
\to h_\mathfrak X
$$
可由开浸入表示。对所有 $\lambda$ 拉回到 $\Spec(A/I_\lambda)$，
得到开子概形 $U_\lambda \subset \Spec(A/I_\lambda)$。然而，若
$I_\lambda \subset I_\mu$，则态射
% P10-E0291
$\Spec(A/I_\lambda) \to \Spec(A/I_\mu)$ 把 $U_\mu$ 拉回到 $U_\lambda$。
因而它们黏合成一个开形式子概形 $\mathfrak U \subset \mathfrak X$。
一个直接的论证（略去）表明
$$
h_\mathfrak U = h_\mathfrak V \times_{h_\mathfrak Y} h_\mathfrak X
$$
由此可见，给定开覆盖 $\mathfrak Y = \bigcup \mathfrak V_j$ 和函子变换
$\varphi : h_\mathfrak X \to h_\mathfrak Y$，可得到 $\mathfrak X$ 的相应开覆盖。
由于 $\mathfrak X$ 是仿射的，可以把该覆盖加细为由仿射形式子概形组成的有限开覆盖
$\mathfrak X = \mathfrak U_1 \cup \ldots \cup \mathfrak U_n$。
换言之，对每个 $i$，存在某个 $j$ 和映射
$\varphi_i : h_{\mathfrak U_i} \to h_{\mathfrak V_j}$，使图表
$$
\xymatrix{
h_{\mathfrak U_i} \ar[r]_{\varphi_i} \ar[d] & h_{\mathfrak V_j} \ar[d] \\
h_{\mathfrak X} \ar[r]^\varphi & h_\mathfrak Y
}
$$
交换。再加上几个论证（略去），这说明只须在
$\mathfrak X$ 和 $\mathfrak Y$ 都是仿射形式概形时证明引理中的双射性。

\medskip\noindent
假设 $\mathfrak X$ 和 $\mathfrak Y$ 是仿射形式概形。写成
$\mathfrak X = \text{Spf}(A)$ 和 $\mathfrak Y = \text{Spf}(B)$。
设 $\varphi : h_\mathfrak X \to h_\mathfrak Y$ 是函子变换，
$I_\lambda \subset A$ 是一族定义理想基本系。典范包含态射
$i_\lambda : \Spec(A/I_\lambda) \to \mathfrak X$ 被映到态射
$\varphi(i_\lambda) : \Spec(A/I_\lambda) \to \mathfrak Y$。
由式 (\ref{formal-spaces-equation-morphisms-affine-formal-schemes})，它对应于连续映射
$\chi_\lambda : B \to A/I_\lambda$。由于 $\varphi$ 是函子变换，
若 $I_\lambda \subset I_\mu$，则复合
$B \to A/I_\lambda \to A/I_\mu$ 等于 $\chi_\mu$。换言之，得到环映射
$$
\chi = \lim \chi_\lambda : B \longrightarrow \lim A/I_\lambda = A
$$
这是连续同态，因为对每个 $\lambda$，$I_\lambda$ 的逆像都是开集
（$A/I_\lambda$ 带离散拓扑，而 $\chi_\lambda$ 连续）。因此，由式
(\ref{formal-spaces-equation-morphisms-affine-formal-schemes}) 得到态射
$\text{Spf}(\chi) : \mathfrak X \to \mathfrak Y$。
我们略去核验：在此情形中，这个构造是引理中映射的逆。

\medskip\noindent
集合论说明。为了使上述论证在给定的概形范畴 $\Sch_\alpha$ 上成立，只须保证证明中使用的所有概形
都同构于 $\Sch_\alpha$ 的对象。事实上，仔细分析可知，只须上述出现的概形
$\Spec(A/I_\lambda)$ 同构于 $\Sch_\alpha$ 的对象。为此，假设
$\mathfrak X$ 的大小不超过 $\Sch_\alpha$ 中某个概形的大小当然就足够了。
\end{proof}


\begin{lemma}
\label{formal-spaces-lemma-formal-scheme-sheaf-fppf}
\begin{slogan}
形式概形是 fpqc 层
\end{slogan}
设 $\mathfrak X$ 是形式概形。其点函子 $h_\mathfrak X$
（见引理 \ref{formal-spaces-lemma-fully-faithful}）满足 fpqc 覆盖的层条件。
\end{lemma}

\begin{proof}
《拓扑》之引理 \ref{topologies-lemma-sheaf-property-fpqc}
把问题归约为 Zariski 覆盖以及 $R \to S$ 忠实平坦时的覆盖
$\{\Spec(S) \to \Spec(R)\}$。我们在引理 \ref{formal-spaces-lemma-fully-faithful}
的证明中已观察到，$h_\mathfrak X$ 满足 Zariski 覆盖的层条件。

\medskip\noindent
设 $R \to S$ 是忠实平坦环映射。用
$\pi : \Spec(S) \to \Spec(R)$ 表示相应的概形态射；它是满且平坦的。
设 $f : \Spec(S) \to \mathfrak X$ 是态射，并且作为从
$\Spec(S \otimes_R S)$ 到 $\mathfrak X$ 的映射，有
$f \circ \text{pr}_1 = f \circ \text{pr}_2$。
由《下降》之引理 \ref{descent-lemma-equiv-fibre-product}，作为底层集合间的映射，
$f$ 具有形式 $f = g \circ \pi$，其中
$g : \Spec(R) \to \mathfrak X$ 是某个（集合论意义下的）映射。
由《态射》之引理 \ref{morphisms-lemma-fpqc-quotient-topology}
以及 $f$ 连续可知 $g$ 连续。

\medskip\noindent
取 $y \in \Spec(R)$。选取包含 $g(y)$ 的仿射形式开子概形
$\mathfrak U \subset \mathfrak X$。设
$\mathfrak U = \text{Spf}(A)$，其中 $A$ 是某个可容许拓扑环。
由上述结论，可选取 $r \in R$，使得
$y \in D(r) \subset g^{-1}(\mathfrak U)$。
$f$ 在 $\pi^{-1}(D(r))$ 上到 $\mathfrak U$ 的限制，由式
(\ref{formal-spaces-equation-morphisms-affine-formal-schemes}) 对应于连续环映射
$A \to S_r$。按对 $f$ 的假设，两个诱导环映射
$A \to S_r \otimes_{R_r} S_r = (S \otimes_R S)_r$ 相等。
注意 $R_r \to S_r$ 忠实平坦。由《下降》之引理
\ref{descent-lemma-ff-exact}，两个箭头
$S_r \to S_r \otimes_{R_r} S_r$ 的等化子是 $R_r$。
于是 $A \to S_r$ 唯一地分解经过一个映射 $A \to R_r$；
这个映射也连续，因为它与映射 $A \to S_r$ 有相同的（开）核。
该映射再由式 (\ref{formal-spaces-equation-morphisms-affine-formal-schemes})
给出态射 $D(r) \to \mathfrak U$。

\medskip\noindent
至此证明了什么？我们证明了：对任意 $y \in \Spec(R)$，存在标准仿射开集
$y \in D(r) \subset \Spec(R)$，使得态射
$f|_{\pi^{-1}(D(r))} : \pi^{-1}(D(r)) \to \mathfrak X$
唯一地分解经过某个态射 $D(r) \to \mathfrak X$。
我们略去核验这些态射黏合成所需态射 $\Spec(R) \to \mathfrak X$。
\end{proof}

\begin{remark}[McQuillan 的变体]
\label{formal-spaces-remark-mcquillan}
McQuillan 给出了形式概形构造的一个变体，见 \cite{McQuillan}。
他建议稍微减弱可容许性条件。回忆，可容许拓扑环是一个完备的
（按我们的约定也分离的）线性拓扑环 $A$，且存在一个定义理想：
开理想 $I$，使得 $0$ 的任意邻域都包含某个 $I^n$，其中 $n \geq 1$。
McQuillan 使用我们将称为{\it 弱可容许}的拓扑环。
弱可容许拓扑环 $A$ 是一个完备的（按我们的约定也分离的）线性拓扑环，
且存在一个{\it 弱定义理想}：开理想 $I$，使得对所有 $f \in I$，
当 $n \to \infty$ 时有 $f^n \to 0$。与可容许情形类似，
若 $I$ 是弱定义理想而 $J \subset A$ 是开理想，则 $I \cap J$ 是弱定义理想。
因此，弱定义理想组成 $0$ 的开邻域基本系，可以沿着与上文大致相同的路线，
以此概念为基础定义一个更大的形式概形范畴。引理 \ref{formal-spaces-lemma-fully-faithful}
和 \ref{formal-spaces-lemma-formal-scheme-sheaf-fppf} 的对应结论在这个背景下仍然成立
（证明相同）。
\end{remark}

\begin{remark}[拓扑空间预层的层化]
\label{formal-spaces-remark-sheafification-of-presheaves-in-top}
\begin{reference}
\cite{Gray}
\end{reference}
本注简要讨论拓扑空间预层的层化。完全相同的论证适用于拓扑阿贝尔群、拓扑环以及
（给定拓扑环上的）拓扑模的预层。为了在正确的一般性下进行讨论，设我们工作在站点
$\mathcal{C}$ 上。只关心拓扑空间 $X$ 上（预）层情形的读者，应把
$\mathcal{C}$ 的对象理解为 $X$ 的开集，把 $\mathcal{C}$ 的态射理解为开集的包含，
并把 $\mathcal{C}$ 中的覆盖理解为 $X$ 中的覆盖，见《站点》之例
\ref{sites-example-site-topological}。以
$\Sh(\mathcal{C}, \textit{Top})$ 表示 $\mathcal{C}$ 上拓扑空间层的范畴，
以 $\textit{PSh}(\mathcal{C}, \textit{Top})$ 表示 $\mathcal{C}$ 上拓扑空间预层的范畴。
设 $\mathcal{F}$ 是 $\mathcal{C}$ 上的拓扑空间预层。其层化
$\mathcal{F}^\#$ 应函子性地对
$\Sh(\mathcal{C}, \textit{Top})$ 中的 $\mathcal{G}$ 满足公式
$$
\Mor_{\textit{PSh}(\mathcal{C}, \textit{Top})}(\mathcal{F}, \mathcal{G})
=
\Mor_{\Sh(\mathcal{C}, \textit{Top})}(\mathcal{F}^\#, \mathcal{G})
$$
换言之，我们要构造包含函子
$\Sh(\mathcal{C}, \textit{Top}) \to
\textit{PSh}(\mathcal{C}, \textit{Top})$ 的左伴随。

我们首先断言 $\Sh(\mathcal{C}, \textit{Top})$ 有极限，且包含函子与极限交换。
事实上，给定范畴 $\mathcal{I}$ 以及到
$\Sh(\mathcal{C}, \textit{Top})$ 的函子 $i \mapsto \mathcal{G}_i$，
我们只须定义
$$
(\lim \mathcal{G}_i)(U) = \lim \mathcal{G}_i(U)
$$
其中右侧极限在拓扑空间范畴中取
（《拓扑》之引理 \ref{topology-lemma-limits}）。这确实定义一个层，
因为极限与极限交换（《范畴》之引理
\ref{categories-lemma-colimits-commute}），特别是与层公理所使用的乘积和等化子交换。
最后，从 $\mathcal{F}$ 到 $\lim \mathcal{G}_i$ 的预层态射显然等同于一族相容的态射
$\mathcal{F} \to \mathcal{G}_i$。换言之，对象 $\lim \mathcal{G}_i$
是拓扑空间预层范畴中的极限，从而当然也是拓扑空间层范畴中的极限。

第二个断言是：任意预层态射 $\mathcal{F} \to \mathcal{G}$，其中
$\mathcal{G}$ 是 $\Sh(\mathcal{C}, \textit{Top})$ 的对象，都分解经过某个大小有界的子层
$\mathcal{G}' \subset \mathcal{G}$。这里，把拓扑空间层 $\mathcal{H}$ 的{\it 大小}
$|\mathcal{H}|$ 定义为基数
$\sup_{U \in \Ob(\mathcal{C})} |\mathcal{H}(U)|$。
为证明该断言，令
$$
\mathcal{G}'(U) =
\left\{
\quad
s \in \mathcal{G}(U)
\quad \middle| \quad
\begin{matrix}
\text{存在覆盖 }\{U_i \to U\}_{i \in I} \\
\text{使得 }s|_{U_i} \in \Im(\mathcal{F}(U_i) \to \mathcal{G}(U_i))
\end{matrix}
\quad
\right\}
$$
给 $\mathcal{G}'(U)$ 赋诱导拓扑。则 $\mathcal{G}'$ 是拓扑空间层（细节略），
$\mathcal{G}' \to \mathcal{G}$ 是一个态射，且给定映射
$\mathcal{F} \to \mathcal{G}$ 分解经过它。此外，$\mathcal{G}'$ 的大小由某个只依赖于
$\mathcal{C}$ 和预层 $\mathcal{F}$ 的基数 $\kappa$ 控制
（提示：按我们的约定，$\mathcal{C}$ 中的覆盖组成一个集合）。综合以上结果，
《范畴》之定理 \ref{categories-theorem-adjoint-functor} 的假设得到满足，
从而得到层化，它是从层到预层的包含函子的左伴随。

最后，设 $p$ 是站点 $\mathcal{C}$ 的一个点，由函子
$u : \mathcal{C} \to \textit{Sets}$ 给出，见《站点》之定义
\ref{sites-definition-point}。对拓扑空间 $M$，规则
$$
U \mapsto \text{Map}(u(U), M) = \prod\nolimits_{x \in u(U)} M
$$
定义的预层在带乘积拓扑后是拓扑空间层。因此，与《站点》之引理
\ref{sites-lemma-point-pushforward-sheaf} 的证明完全相同的论证表明
$\mathcal{F}_p = \mathcal{F}^\#_p$；换言之，层化与在一点取茎交换。
\end{remark}





\section{约定和记号}
\label{formal-spaces-section-conventions}

\noindent
从现在起采用的约定与《空间的性质》之第
\ref{spaces-properties-section-conventions} 节中的约定相似。
因此，从现在起始终假设所有概形都包含在一个大 fppf 站点 $\Sch_{fppf}$ 中。
所考虑的每个环 $A$ 也都满足：$\Spec(A)$（同构于）这个大站点的某个对象。
对拓扑环 $A$，只假设所有离散商都具有这个性质
（不过通常会作更强的假设，比较注 \ref{formal-spaces-remark-set-theoretic}）。

\medskip\noindent
设 $S$ 是概形，$X$ 是 $S$ 上的“空间”，即 $(\Sch/S)_{fppf}$ 上的层。
本章把 $X$ 在 $(\Sch/S)_{fppf}$ 上层范畴中与自身的乘积写作
$X \times_S X$，而不是 $X \times X$。此外，若 $X$ 和 $Y$ 是“空间”，
则我们说“设 $f : X \to Y$ 是态射”，意指 $f$ 是函子的自然变换，
也就是 $(\Sch/S)_{fppf}$ 上层的映射。类似地，若 $U$ 是 $S$ 上的概形，
$X$ 是 $S$ 上的“空间”，则我们说“设 $f : U \to X$ 是态射”或
“设 $g : X \to U$ 是态射”，意指 $f$ 或 $g$ 是层映射
$h_U \to X$ 或 $X \to h_U$，其中 $h_U$ 如《范畴》之例
\ref{categories-example-hom-functor} 所定义。






\section{拓扑环与拓扑模}
\label{formal-spaces-section-topological-module}

\noindent
本节是《代数续篇》之第 \ref{more-algebra-section-topological-ring} 节的继续。
设 $R$ 是拓扑环，$M$ 是线性拓扑化的 $R$-模。当我们说“{\it 设 $M_\lambda$
是一族开子模基本系}”时，意思是每个 $M_\lambda$ 都是开子模，且 $0$ 的任意邻域
都包含某个 $M_\lambda$。换言之，这表示 $M_\lambda$ 是 $M$ 中由子模组成的
$0$ 的邻域基本系。类似地，若 $R$ 是线性拓扑环，则说“{\it 设 $I_\lambda$
是一族开理想基本系}”，意思是 $I_\lambda$ 是 $R$ 中由理想组成的 $0$ 的邻域基本系。

\begin{example}
\label{formal-spaces-example-what-does-it-mean}
设 $R$ 是线性拓扑环，$M$ 是线性拓扑化的 $R$-模。设 $I_\lambda$
是 $R$ 中的一族开理想基本系，$M_\mu$ 是 $M$ 中的一族开子模基本系。
$+ : M \times M \to M$ 的连续性是自动的，而 $R \times M \to M$ 的连续性表示
$$
\forall f, x, \mu\ \exists \lambda, \nu,\ (f + I_\lambda)(x + M_\nu)
\subset fx + M_\mu
$$
当 $M_\nu \subset M_\mu$ 时，有
$fM_\nu + I_\lambda M_\nu \subset M_\mu$，所以该条件等价于
$$
\forall x, \mu\ \exists \lambda\ I_\lambda x \subset M_\mu
$$
然而，给定 $\mu$，未必存在 $\lambda$ 使
$I_\lambda M \subset M_\mu$。例如，考虑带 $t$-进拓扑的
$R = k[[t]]$，以及 $M = \bigoplus_{n \in \mathbf{N}} R$，其开子模基本系为
$$
M_m = \bigoplus\nolimits_{n \in \mathbf{N}} t^{nm}R
$$
由于每个 $x \in M$ 只有有限多个非零坐标，可见给定 $m$ 和 $x$，
存在 $k$ 使 $t^k x \in M_m$。因此，$M$ 是线性拓扑化的 $R$-模，
但给定 $m$ 时，未必存在 $k$ 使 $t^kM \subset M_m$。
另一方面，若 $R \to S$ 是线性拓扑环的连续映射，则相应陈述确实成立；
也就是说，对每个开理想 $J \subset S$，存在开理想 $I \subset R$，使得
$IS \subset J$（读者可容易地从映射 $R \to S$ 的连续性推出这一点）。
\end{example}

\begin{lemma}
\label{formal-spaces-lemma-closed}
设 $R$ 是拓扑环。设 $M$ 是线性拓扑化的 $R$-模，并设
$M_\lambda$（$\lambda \in \Lambda$）是一族开子模基本系。
设 $N \subset M$ 是子模。则 $N$ 的闭包为
$\bigcap_{\lambda \in \Lambda} (N + M_\lambda)$。
\end{lemma}

\begin{proof}
因为每个 $N + M_\lambda$ 都是开集，所以也都是闭集；故其交是闭集。
若 $x \in M$ 不在 $N$ 的闭包中，则对某个 $\lambda$ 有
% P10-E0292
$(x + M_\lambda) \cap N = 0$。所以 $x \not \in N + M_\lambda$。
这就证明了引理。
\end{proof}

\noindent
除非另有说明，我们总给子模和商模赋诱导拓扑。设 $M$ 是拓扑环 $R$ 上的线性拓扑化模，
$0 \to N \to M \to Q \to 0$ 是 $R$-模的短正合列。
若 $M_\lambda$ 是 $M$ 的一族开子模基本系，则
$N \cap M_\lambda$ 是 $N$ 的一族开子模基本系。若
$\pi : M \to Q$ 是商映射，则 $\pi(M_\lambda)$ 是 $Q$ 的一族开子模基本系。
特别地，这些诱导拓扑都是线性拓扑。

\begin{lemma}
\label{formal-spaces-lemma-closure}
设 $R$ 是拓扑环。设 $M$ 是线性拓扑化的 $R$-模，$N \subset M$ 是子模。则
\begin{enumerate}
\item $0 \to N^\wedge \to M^\wedge \to (M/N)^\wedge$ 正合；
\item $N^\wedge$ 是 $N \to M^\wedge$ 的像的闭包。
\end{enumerate}
\end{lemma}

\begin{proof}
设 $M_\lambda$（$\lambda \in \Lambda$）是一族开子模基本系。
则 $N \cap M_\lambda$ 是 $N$ 的一族开子模基本系，并且
% P10-E0293
$M_\lambda + N/N$ 是 $M/N$ 的一族开子模基本系。因此，(1) 来自正合列
$$
0 \to N/N \cap M_\lambda \to M/M_\lambda
\to M/(M_\lambda + N) \to 0
$$
的正合性以及取极限与极限交换这一事实。第二个陈述由此推出，因为
$N \to N^\wedge$ 的像稠密，而 $M^\wedge \to (M/N)^\wedge$ 的核是闭的。
\end{proof}

\begin{lemma}
\label{formal-spaces-lemma-quotient-by-closed}
设 $R$ 是拓扑环。设 $M$ 是完备的线性拓扑化 $R$-模，
$N \subset M$ 是闭子模。若 $M$ 有可数的 $0$ 邻域基本系，则
$M/N$ 完备，且映射 $M \to M/N$ 是开映射。
\end{lemma}

\begin{proof}
设 $M_n$（$n \in \mathbf{N}$）是 $M$ 的一族开子模基本系。
可以假设对所有 $n$ 都有 $M_{n + 1} \subset M_n$。
系统 $(M_n + N)/N$ 是 $M/N$ 中的一个基本系。因此，只须证明
$M/N = \lim M/(M_n + N)$。考虑短正合列
$$
0 \to N/N \cap M_n \to M/M_n \to M/(M_n + N) \to 0
$$
由于系统 $\{N/N\cap M_n\}$ 的转移映射都是满射，
《代数》之引理 \ref{algebra-lemma-ML-exact-sequence} 表明
$M = \lim M/M_n$（利用 $M$ 的完备性）满射到
$\lim M/(M_n + N)$。由于 $N$ 闭，映射
$M \to \lim M/(M_n + N)$ 的核为 $N$
（见引理 \ref{formal-spaces-lemma-closed}）。最后，按商拓扑的定义，$M \to M/N$ 是开映射。
\end{proof}

\begin{lemma}
\label{formal-spaces-lemma-ses}
\begin{reference}
\cite[Theorem 8.1]{Ma}
\end{reference}
设 $R$ 是拓扑环。设 $M$ 是线性拓扑化的 $R$-模，$N \subset M$ 是子模。
假设 $M$ 有可数的 $0$ 邻域基本系。则
\begin{enumerate}
\item $0 \to N^\wedge \to M^\wedge \to (M/N)^\wedge \to 0$ 正合；
\item $N^\wedge$ 是 $N \to M^\wedge$ 的像的闭包；
\item $M^\wedge \to (M/N)^\wedge$ 是开映射。
\end{enumerate}
\end{lemma}

\begin{proof}
由引理 \ref{formal-spaces-lemma-closure}，列
$0 \to N^\wedge \to M^\wedge \to (M/N)^\wedge$ 正合，且 (2) 成立。
这给出典范映射 $c : M^\wedge/N^\wedge \to (M/N)^\wedge$。
由引理 \ref{formal-spaces-lemma-quotient-by-closed}，模 $M^\wedge/N^\wedge$ 完备，
且 $M^\wedge \to M^\wedge/N^\wedge$ 是开映射。
由完备化的泛性质，得到典范映射
$b : (M/N)^\wedge \to M^\wedge/N^\wedge$。
由于 $b$ 和 $c$ 在一个稠密子集上互逆，它们彼此互逆。
\end{proof}

\begin{lemma}
\label{formal-spaces-lemma-completion-adic-star}
设 $R$ 是拓扑环，$M$ 是拓扑 $R$-模。设 $I \subset R$ 是有限生成理想。
假设 $M$ 有一个拓扑为 $I$-进拓扑的开子模。则 $M^\wedge$
也有一个拓扑为 $I$-进拓扑的开子模，并且对所有 $n \geq 1$ 有
$M^\wedge/I^n M^\wedge = M/I^nM$。
\end{lemma}

\begin{proof}
设 $M' \subset M$ 是一个拓扑为 $I$-进拓扑的开子模。
则 $\{I^nM'\}_{n \geq 1}$ 是 $M$ 的一族开子模基本系。因此
$M^\wedge = \lim M/I^nM'$ 包含
$(M')^\wedge = \lim M'/I^nM'$ 作为开子模，并且由《代数》之引理
\ref{algebra-lemma-hathat-finitely-generated}，$(M')^\wedge$ 上的拓扑是 $I$-进拓扑。
因为 $I$ 有限生成，$I^n$ 也有限生成，设其生成元为 $f_1, \ldots, f_r$。
注意，满射
$(f_1, \ldots, f_r) : M^{\oplus r} \to I^n M$
按上述对 $M$ 上拓扑的描述是连续且开的。把引理 \ref{formal-spaces-lemma-ses}
应用于这个满射以及短正合列
$0 \to I^nM \to M \to M/I^nM \to 0$，可知
$$
(f_1, \ldots, f_r) :
(M^\wedge)^{\oplus r} \longrightarrow M^\wedge
$$
满射到满射 $M^\wedge \to M/I^nM$ 的核。由于
$f_1, \ldots, f_r$ 生成 $I^n$，结论成立。
\end{proof}

\begin{definition}
\label{formal-spaces-definition-toplogy-tensor-product}
\begin{reference}
比较 \cite[0, Section 7.7]{EGA}
\end{reference}
设 $R$ 是拓扑环，$M$ 和 $N$ 是线性拓扑化的 $R$-模。
$M$ 与 $N$ 的{\it 张量积}是（通常的）张量积 $M \otimes_R N$，
并赋予如下线性拓扑：规定
$$
\Im(M_\mu \otimes_R N + M \otimes_R N_\nu \longrightarrow M \otimes_R N)
$$
组成一族开子模基本系，其中 $M_\mu \subset M$ 和 $N_\nu \subset N$
分别遍历 $M$ 和 $N$ 的开子模基本系。{\it 完备张量积}
$$
M \widehat{\otimes}_R N =
\lim M \otimes_R N/(M_\mu \otimes_R N + M \otimes_R N_\nu) =
\lim M/M_\mu \otimes_R N/N_\nu
$$
是该张量积的完备化。
\end{definition}

\noindent
注意，$R$ 上的拓扑对张量积或完备张量积的构造并无影响。
若 $R \to A$ 和 $R \to B$ 是线性拓扑环的连续映射，
则上述构造给出张量积 $A \otimes_R B$ 和完备张量积
$A \widehat{\otimes}_R B$。

\medskip\noindent
这里记录注 \ref{formal-spaces-remark-mcquillan} 中引入的概念。

\begin{definition}
\label{formal-spaces-definition-weakly-admissible}
设 $A$ 是线性拓扑环。
\begin{enumerate}
\item 若当 $n \to \infty$ 时 $f^n \to 0$，则称元素 $f \in A$
为{\it 拓扑幂零的}。
\item $A$ 的一个{\it 弱定义理想}是完全由拓扑幂零元素组成的开理想
$I \subset A$。
\item 若 $A$ 有弱定义理想，则称 $A$ {\it 弱预可容许}。
\item 若 $A$ 弱预可容许且完备\footnote{按我们的约定，这包括分离性。}，
则称 $A$ {\it 弱可容许}。
\end{enumerate}
\end{definition}

\noindent
给定线性拓扑环 $A$ 中的弱定义理想 $I$ 和开理想 $J$，
交 $I \cap J$ 是弱定义理想。因此，一旦存在一个弱定义理想，
就存在一族由弱定义理想组成的开理想基本系。特别地，给定弱可容许拓扑环 $A$，有
$A = \lim A/I_\lambda$，其中 $\{I_\lambda\}$ 是一族弱定义理想基本系。

\begin{lemma}
\label{formal-spaces-lemma-weakly-admissible-henselian}
设 $A$ 是弱可容许拓扑环，$I \subset A$ 是弱定义理想。
则 $(A, I)$ 是 Hensel 对。
\end{lemma}

\begin{proof}
设 $A \to A'$ 是 étale 环映射，$\sigma : A' \to A/I$ 是 $A$-代数映射。
由《代数续篇》之引理 \ref{more-algebra-lemma-characterize-henselian-pair}，
只须把 $\sigma$ 提升为 $A$-代数映射 $A' \to A$。
为此，由于 $A$ 完备，只须对每个开理想 $J \subset I$，找到唯一的
$A$-代数映射 $A' \to A/J$ 来提升 $\sigma$。由于 $I$ 是弱定义理想，
理想 $I/J$ 局部幂零。结论遂由《代数续篇》之引理
\ref{more-algebra-lemma-locally-nilpotent-henselian} 得出。
\end{proof}

\begin{lemma}
\label{formal-spaces-lemma-topologically-nilpotent}
设 $B$ 是线性拓扑环。$B$ 中拓扑幂零元素的集合是 $B$ 的闭根式理想。
设 $\varphi : A \to B$ 是线性拓扑环的连续映射。
\begin{enumerate}
\item 若 $f \in A$ 拓扑幂零，则 $\varphi(f)$ 拓扑幂零。
\item 若 $I \subset A$ 完全由拓扑幂零元素组成，则
$\varphi(I)B$ 的闭包完全由拓扑幂零元素组成。
\end{enumerate}
\end{lemma}

\begin{proof}
令 $\mathfrak b \subset B$ 为拓扑幂零元素的集合。
我们略去 $\mathfrak b$ 是根式理想的证明（这是练习定义的好题目）。
设 $g$ 是 $\mathfrak b$ 闭包中的元素。目标是证明 $g$ 拓扑幂零。
设 $J \subset B$ 是开理想。须证明对某个 $e \geq 1$ 有 $g^e \in J$。
由引理 \ref{formal-spaces-lemma-closed}，有 $g \in \mathfrak b + J$。
因此，对某个 $f \in \mathfrak b$ 和 $h \in J$ 有 $g = f + h$。
选取 $m \geq 1$ 使 $f^m \in J$，则如所需有 $g^{m + 1} \in J$。

\medskip\noindent
设 $\varphi : A \to B$ 如引理陈述。(1) 显然，(2) 由此以及
$\mathfrak b$ 是闭理想这一事实推出。
\end{proof}

\begin{lemma}
\label{formal-spaces-lemma-closure-image-ideal}
设 $A \to B$ 是线性拓扑环的连续映射，$I \subset A$ 是理想。
则 $IB$ 的闭包是 $B \to B \widehat{\otimes}_A A/I$ 的核。
\end{lemma}

\begin{proof}
设 $J_\mu$ 是 $B$ 的一族开理想基本系。由引理 \ref{formal-spaces-lemma-closed}，
% P10-E0294
$IB$ 的闭包是 $\bigcap (IB + J_\lambda)$。设 $I_\mu$ 是 $A$
中的一族开理想基本系。则
$$
B \widehat{\otimes}_A A/I =
\lim (B/J_\lambda \otimes_A A/(I_\mu + I)) =
\lim B/(J_\lambda + I_\mu B + I B)
$$
由于 $A \to B$ 连续，对每个 $\lambda$ 都存在 $\mu$ 使
$I_\mu B \subset J_\lambda$，见例 \ref{formal-spaces-example-what-does-it-mean} 中的讨论。
因此，该极限可写成 $\lim B/(J_\lambda + IB)$，结论显然。
\end{proof}

\begin{lemma}
\label{formal-spaces-lemma-completed-tensor-product}
设 $B \to A$ 和 $B \to C$ 是线性拓扑环的连续同态。
\begin{enumerate}
\item 若 $A$ 和 $C$ 弱预可容许，则 $A \widehat{\otimes}_B C$ 弱可容许。
\item 若 $A$ 和 $C$ 预可容许，则 $A \widehat{\otimes}_B C$ 可容许。
\item 若 $A$ 和 $C$ 有可数的开理想基本系，则
$A \widehat{\otimes}_B C$ 有可数的开理想基本系。
\item 若 $A$ 和 $C$ 是预-adic 的，并有有限生成的定义理想，
则 $A \widehat{\otimes}_B C$ 是 adic 的，并有有限生成的定义理想。
\item 若 $A$ 和 $C$ 是预-adic Noether 环，并且
$B/\mathfrak b \to A/\mathfrak a$ 是有限型的，其中
$\mathfrak a \subset A$ 和 $\mathfrak b \subset B$ 是拓扑幂零元素的理想，
则 $A \widehat{\otimes}_B C$ 是 adic Noether 环。
\end{enumerate}
\end{lemma}

\begin{proof}
设 $I_\lambda \subset A$（$\lambda \in \Lambda$）和
$J_\mu \subset C$（$\mu \in M$）分别是开理想基本系。按定义
$$
A \widehat{\otimes}_B C =
\lim_{\lambda, \mu} A/I_\lambda \otimes_B C/J_\mu
$$
并带极限拓扑。因此，一族开理想基本系由映射
$A \widehat{\otimes}_B C \to A/I_\lambda \otimes_B C/J_\mu$
的核 $K_{\lambda, \mu}$ 给出。注意，$K_{\lambda, \mu}$ 是理想
$I_\lambda(A \widehat{\otimes}_B C) +
J_\mu(A \widehat{\otimes}_B C)$ 的闭包。最后，有一个像稠密的环同态
$\tau : A \otimes_B C \to A \widehat{\otimes}_B C$。

\medskip\noindent
(1) 的证明。若 $I_\lambda$ 和 $J_\mu$ 完全由拓扑幂零元素组成，
则由引理 \ref{formal-spaces-lemma-topologically-nilpotent}，$K_{\lambda, \mu}$ 也如此。
所以按定义，$A \widehat{\otimes}_B C$ 弱可容许。

\medskip\noindent
(2) 的证明。假设对某个 $\lambda_0$ 和 $\mu_0$，理想
$I = I_{\lambda_0} \subset A$ 和 $J_{\mu_0} \subset C$ 是定义理想。
所以，对每个 $\lambda$，存在 $n$ 使 $I^n \subset I_\lambda$；
对每个 $\mu$，存在 $m$ 使 $J^m \subset J_\mu$。于是
$$
\left(I(A \widehat{\otimes}_B C) + J(A \widehat{\otimes}_B C)\right)^{n + m}
\subset
I_\lambda(A \widehat{\otimes}_B C) + J_\mu(A \widehat{\otimes}_B C)
$$
由此，开理想 $K = K_{\lambda_0, \mu_0}$ 满足
$K^{n + m} \subset K_{\lambda, \mu}$。所以 $K$ 是
$A \widehat{\otimes}_B C$ 的定义理想，而 $A \widehat{\otimes}_B C$ 按定义可容许。

\medskip\noindent
(3) 的证明。若 $\Lambda$ 和 $M$ 可数，则 $\Lambda \times M$ 也可数。

\medskip\noindent
(4) 的证明。假设 $\Lambda = \mathbf{N}$、$M = \mathbf{N}$，
并且有有限生成理想 $I \subset A$ 和 $J \subset C$，使
$I_n = I^n$ 且 $J_n = J^n$。则
$$
I(A \widehat{\otimes}_B C) + J(A \widehat{\otimes}_B C)
$$
是有限生成理想，并且容易看出 $A \widehat{\otimes}_B C$ 是
$A \otimes_B C$ 关于该理想的完备化。因此，(4) 由《代数》之引理
\ref{algebra-lemma-hathat-finitely-generated} 得出。

\medskip\noindent
(5) 的证明。令 $\mathfrak c \subset C$ 为拓扑幂零元素的理想。
由于 $A$ 和 $C$ 是 adic Noether 环，可知 $\mathfrak a$ 和
$\mathfrak c$ 是定义理想（细节略）。由 (4) 已知
$A \widehat{\otimes}_B C$ 是 adic 的，并且
$\mathfrak a(A \widehat{\otimes}_B C) +
\mathfrak c(A \widehat{\otimes}_B C)$ 是有限生成的定义理想。由于
$$
A \widehat{\otimes}_B C /
\left(\mathfrak a(A \widehat{\otimes}_B C) +
\mathfrak c(A \widehat{\otimes}_B C)\right)
=
A/\mathfrak a \otimes_{B/\mathfrak b} C/\mathfrak c
$$
是 Noether 环 $C/\mathfrak c$ 上的有限型代数，结论由《代数》之引理
\ref{algebra-lemma-completion-Noetherian} 得出。
\end{proof}









\section{taut 环映射}
\label{formal-spaces-section-taut-ring-maps}

\noindent
为线性拓扑环之间的连续映射所具有的下述性质取一个名称是方便的。

\begin{definition}
\label{formal-spaces-definition-taut}
设 $\varphi : A \to B$ 是线性拓扑环的连续映射。若对每个开理想
$I \subset A$，理想 $\varphi(I)B$ 的闭包都是开的，并且这些闭包组成一族开理想基本系，
则称 $\varphi$ 是 {\it taut} 的\footnote{这是非标准记号。该定义可推广到模：
若对每个开理想 $I \subset A$，$IM$ 在 $M$ 中的闭包都是开的，
且这些闭包组成 $M$ 中 $0$ 的一族邻域基本系，则称线性拓扑化的 $A$-模
$M$ 是 $A$-taut 的。}。
\end{definition}

\noindent
若 $\varphi : A \to B$ 是线性拓扑环的连续映射，而 $I_\lambda$
是 $A$ 的一族开理想基本系，则 $\varphi$ 是 taut 的，当且仅当
% P10-E0295
$I_\lambda B$ 的闭包都是开的，并组成 $A$ 中的一族开理想基本系。

\begin{lemma}
\label{formal-spaces-lemma-taut-weakly-admissible}
设 $\varphi : A \to B$ 是弱可容许拓扑环的连续映射。下列条件等价：
\begin{enumerate}
\item $\varphi$ 是 taut 的；
\item 对 $A$ 的每个弱定义理想 $I \subset A$，$\varphi(I)B$ 的闭包
都是 $B$ 的弱定义理想，并且这些闭包组成 $B$ 的一族弱定义理想基本系。
\end{enumerate}
\end{lemma}

\begin{proof}
定义 \ref{formal-spaces-definition-taut} 后的说明表明 (2) 推出 (1)。反过来，
假设 $\varphi$ 是 taut 的。若 $I \subset A$ 是弱定义理想，则按 taut 性的定义，
$\varphi(I)B$ 的闭包是开的；由引理 \ref{formal-spaces-lemma-topologically-nilpotent}，
它完全由拓扑幂零元素组成。因此 $\varphi(I)B$ 的闭包是弱定义理想。
此外，按 taut 性的定义，这些理想组成一族开理想基本系，所以 (2) 成立。
\end{proof}

\begin{lemma}
\label{formal-spaces-lemma-completion-taut}
设 $A$ 是线性拓扑环。从 $A$ 到其完备化的映射
$A \to A^\wedge$ 是 taut 的。
\end{lemma}

\begin{proof}
设 $I_\lambda$ 是 $A$ 的一族开理想基本系。回忆
$A^\wedge = \lim A/I_\lambda$ 带极限拓扑；这表示各核
$J_\lambda = \Ker(A^\wedge \to A/I_\lambda)$
组成 $A^\wedge$ 的一族开理想基本系。由于 $J_\lambda$ 是
$I_\lambda A^\wedge$ 的闭包（比较引理 \ref{formal-spaces-lemma-closure-image-ideal}），结论成立。
\end{proof}

\begin{lemma}
\label{formal-spaces-lemma-composition-taut}
设 $A \to B$ 和 $B \to C$ 是线性拓扑环的连续同态。
若 $A \to B$ 和 $B \to C$ 都是 taut 的，则 $A \to C$ 是 taut 的。
\end{lemma}

\begin{proof}
略。提示：若 $I \subset A$ 是理想，$J$ 是 $IB$ 的闭包，
则 $JC$ 的闭包等于 $IC$ 的闭包。
\end{proof}

\begin{lemma}
\label{formal-spaces-lemma-permanence-taut}
设 $A \to B$ 和 $B \to C$ 是线性拓扑环的连续同态。
若 $A \to C$ 是 taut 的，则 $B \to C$ 是 taut 的。
\end{lemma}

\begin{proof}
设 $J \subset B$ 是开理想，其逆像为 $I \subset A$。
则 $JC$ 的闭包包含 $IC$ 的闭包。由于 $A \to C$ 是 taut 的，这个闭包是开的。
设 $I_\lambda$ 是 $A$ 的一族开理想基本系，$K_\lambda$ 是
$I_\lambda C$ 的闭包。由于 $A \to C$ 是 taut 的，它们组成
$C$ 的一族开理想基本系。以 $J_\lambda \subset B$ 表示 $K_\lambda$ 的逆像。
则 $J_\lambda C$ 的闭包是 $K_\lambda$。因此可见，当 $J$ 遍历 $B$
的开理想时，理想 $JC$ 的闭包组成 $C$ 的一族开理想基本系。
\end{proof}

\begin{lemma}
\label{formal-spaces-lemma-base-change-taut}
设 $A \to B$ 和 $A \to C$ 是线性拓扑环的连续同态。
若 $A \to B$ 是 taut 的，则
$C \to B \widehat{\otimes}_A C$ 是 taut 的。
\end{lemma}

\begin{proof}
设 $K \subset C$ 是开理想。选取任意开理想 $I \subset A$，使其在 $C$ 中的像包含在
% P10-E0296
$J$ 中。按假设，$IB$ 的闭包 $J$ 是开的。由于 $A \to B$ 是 taut 的，
可见 $B \widehat{\otimes}_A C$ 是环
$B/J \otimes_{A/I} C/K$ 对所有 $K$ 和 $I$ 的选择所取的极限；也就是说，理想
$J(B \widehat{\otimes}_A C) + K(B \widehat{\otimes}_A C)$
组成一族开理想基本系。又因为
$B \to B \widehat{\otimes}_A C$ 连续，并且 $I$ 映入 $K$，
可知 $J$ 映入 $K(B \widehat{\otimes}_A C)$ 的闭包。
因此该闭包等于
$J(B \widehat{\otimes}_A C) + K(B \widehat{\otimes}_A C)$，证明完成。
\end{proof}

\begin{lemma}
\label{formal-spaces-lemma-taut-ascent-countable}
设 $\varphi : A \to B$ 是线性拓扑环的连续同态。
若 $\varphi$ 是 taut 的，且 $A$ 有可数的开理想基本系，
则 $B$ 有可数的开理想基本系。
\end{lemma}

\begin{proof}
由定义立即得到。
\end{proof}

\begin{lemma}
\label{formal-spaces-lemma-taut-ascent-weakly-admissible}
设 $\varphi : A \to B$ 是线性拓扑环的连续同态。
若 $\varphi$ 是 taut 的，且 $A$ 弱预可容许，则 $B$ 弱预可容许。
\end{lemma}

\begin{proof}
设 $I \subset A$ 是弱定义理想。则 $IB$ 的闭包 $J$ 是开的，
并且由引理 \ref{formal-spaces-lemma-topologically-nilpotent}，它完全由拓扑幂零元素组成。
因此 $J$ 是 $B$ 的弱定义理想。
\end{proof}

\begin{lemma}
\label{formal-spaces-lemma-taut-ascent-admissible}
设 $\varphi : A \to B$ 是线性拓扑环的连续同态。
若 $\varphi$ 是 taut 的，且 $A$ 预可容许，则 $B$ 预可容许。
\end{lemma}

\begin{proof}
设 $I \subset A$ 是定义理想，$I_\lambda \subset A$ 是一族开理想基本系。
则 $IB$ 的闭包 $J$ 是开的，而 $I_\lambda B$ 的闭包 $J_\lambda$
也是开的，并组成 $B$ 的一族开理想基本系。对每个 $\lambda$，
都存在 $n$ 使 $I^n \subset I_\lambda$。注意，$J^n$ 包含在 $I^nB$ 的闭包中。
所以 $J^n \subset J_\lambda$，从而 $J$ 是定义理想。
\end{proof}

\begin{lemma}
\label{formal-spaces-lemma-dense-image-surjective}
设 $\varphi : A \to B$ 是线性拓扑环的连续同态。假设
\begin{enumerate}
\item $\varphi$ 是 taut 的，且像稠密；
\item $A$ 完备，并有可数的开理想基本系；
\item $B$ 分离。
\end{enumerate}
则 $\varphi$ 是满的开映射，$B$ 完备，并且对某个闭理想
$K \subset A$ 有 $B = A/K$。
\end{lemma}

\begin{proof}
由开映射引理（《代数续篇》之引理
\ref{more-algebra-lemma-open-mapping}）和 $\varphi$ 的 taut 性，
可知映射 $\varphi$ 是开的。由于 $\varphi$ 的像稠密，$\varphi$ 是满射。
$\varphi$ 的核 $K$ 是闭的，因为 $\varphi$ 连续。因此
$B = A/K$ 完备，例如可见引理 \ref{formal-spaces-lemma-quotient-by-closed}。
\end{proof}







\section{Adic 环映射}
\label{formal-spaces-section-adic-ring-maps}

\noindent
作如下定义。

\begin{definition}
\label{formal-spaces-definition-adic-homomorphism}
设 $A$ 和 $B$ 是预-adic 拓扑环。若存在 $A$ 的定义理想
$I \subset A$，使得 $B$ 上的拓扑为 $I$-进拓扑，则称环同态
$\varphi : A \to B$ 是 {\it adic} 的\footnote{这可能是非标准术语。}。
\end{definition}

\noindent
若 $\varphi : A \to B$ 是预-adic 环的 adic 同态，则 $\varphi$ 连续，
并且对 $A$ 的每个定义理想 $I$，$B$ 上的拓扑都是 $I$-进拓扑。

\begin{lemma}
\label{formal-spaces-lemma-composition-adic}
设 $A \to B$ 和 $B \to C$ 是预-adic 环的连续同态。
若 $A \to B$ 和 $B \to C$ 都是 adic 的，则 $A \to C$ 是 adic 的。
\end{lemma}

\begin{proof}
略。
\end{proof}

\begin{lemma}
\label{formal-spaces-lemma-permanence-adic}
设 $A \to B$ 和 $B \to C$ 是预-adic 环的连续同态。
若 $A \to C$ 是 adic 的，则 $B \to C$ 是 adic 的。
\end{lemma}

\begin{proof}
选取 $A$ 的一个定义理想 $I$。因为 $A \to C$ 是 adic 的，
$IC$ 是 $C$ 的定义理想。因为 $B \to C$ 连续，可以找到映入 $IC$
的定义理想 $J \subset B$。因为 $A \to B$ 连续，$J$ 在 $I$ 中的逆像
$I' \subset I$ 也是 $A$ 的定义理想。所以
$I'C \subset JC \subset IC$ 夹在两个定义理想之间，本身也是定义理想。
\end{proof}

\begin{lemma}
\label{formal-spaces-lemma-adic-taut}
设 $\varphi : A \to B$ 是预-adic 拓扑环之间的连续同态。
若 $\varphi$ 是 adic 的，则 $\varphi$ 是 taut 的。
\end{lemma}

\begin{proof}
由定义立即得到。
\end{proof}

\noindent
下一个引理说明两件事：
\begin{enumerate}
\item adic 性沿完备线性拓扑环的 taut 映射上升；
\item 对 adic 环之间的连续映射 $\varphi : A \to B$，
若 $A$ 有有限生成的定义理想，则“$\varphi$ 是 taut 的”与
“$\varphi$ 是 adic 的”等价。
\end{enumerate}
由于 (2)，可以说“taut 性”把“adic 性”推广到了任意线性拓扑环之间的连续环映射。
另见第 \ref{formal-spaces-section-adic} 节。

\begin{lemma}
\label{formal-spaces-lemma-taut-is-adic}
设 $\varphi : A \to B$ 是线性拓扑环的连续映射。
若 $\varphi$ 是 taut 的，$A$ 是预-adic 的且有有限生成的定义理想，
而 $B$ 完备，则 $B$ 是 adic 的且有有限生成的定义理想，并且环映射
$\varphi$ 是 adic 的。
\end{lemma}

\begin{proof}
选取 $A$ 的一个有限生成定义理想 $I$。令 $J_n$ 为
$\varphi(I^n)B$ 在 $B$ 中的闭包。由于 $B$ 完备，有
$B = \lim B/J_n$。令 $B' = \lim B/I^nB$ 为 $B$ 的 $I$-进完备化。
由《代数》之引理 \ref{algebra-lemma-hathat-finitely-generated}，
$B'$ 上的 $I$-进拓扑完备，并且 $B'/I^nB' = B/I^nB$。
所以环映射 $B' \to B$ 连续，并且像稠密，因为对所有 $n$，
$B' \to B/I^nB \to B/J_n$ 都是满射。最后，映射 $B' \to B$ 是 taut 的，
因为 $(I^nB')B = I^nB$，且 $A \to B$ 是 taut 的。
由引理 \ref{formal-spaces-lemma-dense-image-surjective}，$B' \to B$ 是满的开映射。
因此 $B$ 上的拓扑是 $I$-进拓扑，证明完成。
\end{proof}








\section{弱 adic 环}
\label{formal-spaces-section-weakly-adic}

\noindent
建议读者跳过本节。下面是 adic 环的一个自然推广。

\begin{definition}
\label{formal-spaces-definition-weakly-adic}
\begin{reference}
\cite[Definition 8.3.8]{Gabber-Ramero}
\end{reference}
设 $A$ 是线性拓扑环。
\begin{enumerate}
\item 若存在理想 $I \subset A$，使得对所有 $n \geq 0$，$I^n$ 的闭包都是开的，
且这些闭包组成一族开理想基本系，则称 $A$ {\it 弱预-adic}
\footnote{在 \cite{Gabber-Ramero} 中，作者称 $A$ 是 {\it $c$-adic} 的。}。
\item 若 $A$ 弱预-adic 且完备\footnote{按我们的约定，这包括分离性。}，
则称 $A$ {\it 弱 adic}。
\end{enumerate}
\end{definition}

\noindent
对完备线性拓扑环，有下列蕴含关系
$$
\xymatrix{
\text{adic 且 Noether} \ar@{=>}[d] \\
\text{adic 且有有限生成的定义理想} \ar@{=>}[d] \\
\text{adic} \ar@{=>}[d] \\
\text{弱 adic} \ar@{=>}[d] \\
\text{可容许且第一可数} \ar@{=>}[d] \ar@{=>}[r] &
\text{可容许} \ar@{=>}[d] \\
\text{弱可容许且第一可数} \ar@{=>}[r] &
\text{弱可容许}
}
$$
这里“第一可数”表示拓扑环有可数的开理想基本系。
对不完备的线性拓扑环也有类似的蕴含图表
（即使用预-adic、弱预-adic、预可容许和弱预可容许这些概念）。
与预-adic 环的情形相反，弱预-adic 环的完备化是弱 adic 的；
下面刻画弱预-adic 环的引理说明了这一点。

\begin{lemma}
\label{formal-spaces-lemma-weakly-adic}
设 $A$ 是线性拓扑环。下列条件等价：
\begin{enumerate}
\item $A$ 弱预-adic；
\item 存在 taut 连续环映射 $A' \to A$，其中 $A'$ 是预-adic 拓扑环；
\item $A$ 预可容许，且存在一个定义理想 $I$，使得对所有 $n \geq 1$，
$I^n$ 的闭包都是开的；
\item $A$ 预可容许，并且对每个定义理想 $I$ 以及所有 $n \geq 1$，
$I^n$ 的闭包都是开的。
\end{enumerate}
弱预-adic 环的完备化是弱 adic 的。若 $A$ 弱 adic，则 $A$ 可容许，
并有可数的开理想基本系。
\end{lemma}

\begin{proof}
假设 (1)。选取理想 $I$，使得对所有 $n$，$I^n$ 的闭包都是开的，
且这些闭包组成一族开理想基本系。以 $A'$ 表示赋予 $I$-进拓扑的环 $A$。
则 $A' \to A$ 按定义是 taut 的，故 (2) 成立。

\medskip\noindent
假设 (2)。设 $I' \subset A'$ 是定义理想，以 $I$ 表示 $I'A$ 的闭包。
$A' \to A$ 的 taut 性表示，$(I')^nA$ 的闭包 $I_n$ 都是开的，
并组成一族开理想基本系。因此 $I = I_1$ 是开的，而 $I^n$ 的闭包等于
$I_n$，从而是开的并组成一族开理想基本系。特别地，$I$ 是一个定义理想，
且所有 $I^n$ 的闭包都是开的。所以 (3) 成立。

\medskip\noindent
若 $I \subset A$ 如 (3) 所述，则 $I$ 是定义 \ref{formal-spaces-definition-weakly-adic}
中的理想，从而 (1) 成立。又若 $I' \subset A$ 是任意其他定义理想，
则 $I'$ 是开的（见《代数续篇》之定义
\ref{more-algebra-definition-topological-ring}），故它包含某个 $I^n$，其中
$n \geq 1$。因此对所有 $m \geq 1$，$(I')^m$ 包含 $I^{nm}$，
从而 $(I')^m$ 的闭包对所有 $m$ 都是开的。这样便看出 (3) 推出 (4)。
(4) $\Rightarrow$ (3) 是平凡的。

\medskip\noindent
设 $A$ 弱预-adic。选取如 (2) 中的 $A' \to A$。
由引理 \ref{formal-spaces-lemma-completion-taut} 和 \ref{formal-spaces-lemma-composition-taut}，
复合 $A' \to A^\wedge$ 是 taut 的。因此，由 (2) 与 (1) 的等价性，
$A^\wedge$ 弱预-adic。因为线性拓扑环 $A$ 的完备化是完备的
（《代数续篇》之第 \ref{more-algebra-section-topological-ring} 节），
可见 $A^\wedge$ 弱 adic。

\medskip\noindent
设 $A$ 弱 adic。则 $A$ 完备，并且
% P10-E0297
由 (1) $\Rightarrow$ (3)，$A$ 且预可容许，所以 $A$ 可容许。
当然，按定义 $A$ 有可数的开理想基本系。
\end{proof}

\noindent
给出两个判据，保证弱 adic 环是 adic 的，并有有限生成的定义理想。

\begin{lemma}
\label{formal-spaces-lemma-weakly-adic-finite-generation-bis}
设 $A$ 是完备线性拓扑环。设 $I \subset A$ 是有限生成理想，
并且对所有 $n \geq 0$，$I^n$ 的闭包都是开的，且这些闭包组成一族开理想基本系。
则 $A$ 是 adic 的，并有有限生成的定义理想。
\end{lemma}

\begin{proof}
以 $A'$ 表示赋予 $I$-进拓扑的环 $A$。这些假设说明 $A' \to A$ 是 taut 的。
结论由引理 \ref{formal-spaces-lemma-taut-is-adic} 得出（确切地说，该引理还告诉我们 $I$ 是定义理想）。
\end{proof}

\begin{lemma}
\label{formal-spaces-lemma-weakly-adic-finite-generation}
设 $A$ 是弱 adic 拓扑环。设 $I$ 是定义理想，并且 $I/I_2$ 是有限生成模，
其中 $I_2$ 是 $I^2$ 的闭包。则 $A$ 是 adic 的，并有有限生成的定义理想。
\end{lemma}

\begin{proof}
我们使用引理 \ref{formal-spaces-lemma-weakly-adic} 的刻画而不再说明。
选取 $f_1, \ldots, f_r \in I$，其像生成 $I/I_2$。令
$I' = (f_1, \ldots, f_r)$，则 $I' + I_2 = I$。
此时 $I_2$ 是
$I^2 = (I' + I_2)^2 \subset I' + I_3$ 的闭包，其中 $I_3$ 是 $I^3$ 的闭包。
所以 $I' + I_3 = I$。如此继续，可见对所有 $n \geq 2$，
$I' + I_n = I$，其中 $I_n$ 是 $I^n$ 的闭包。换言之，$I'$ 在 $A$ 中的闭包为 $I$。
所以 $(I')^n$ 的闭包为 $I_n$。因此，$(I')^n$ 的闭包组成 $A$ 的一族开理想基本系。
结论由引理 \ref{formal-spaces-lemma-weakly-adic-finite-generation-bis} 得出。
\end{proof}

\noindent
“弱预-adic”性质的一个关键特征是：它沿线性拓扑环的 taut 环同态上升。

\begin{lemma}
\label{formal-spaces-lemma-taut-ascent-weakly-adic}
设 $\varphi : A \to B$ 是线性拓扑环的连续同态。
若 $\varphi$ 是 taut 的，且 $A$ 弱预-adic，则 $B$ 弱预-adic。
\end{lemma}

\begin{proof}
设 $I \subset A$ 是理想，使 $I^n$ 的闭包 $I_n$ 都是开的，
且这些闭包定义一族开理想基本系。则 $I^nB$ 的闭包等于 $I_nB$ 的闭包。
由于 $\varphi$ 是 taut 的，这些闭包都是开的，并组成 $B$ 的一族开理想基本系。
故 $B$ 弱预-adic。
\end{proof}

\begin{lemma}
\label{formal-spaces-lemma-completed-tensor-product-weakly-adic}
设 $B \to A$ 和 $B \to C$ 是线性拓扑环的连续同态。
若 $A$ 和 $C$ 弱预-adic，则 $A \widehat{\otimes}_B C$ 弱 adic。
\end{lemma}

\begin{proof}
我们使用引理 \ref{formal-spaces-lemma-weakly-adic} 的刻画而不再说明。
由引理 \ref{formal-spaces-lemma-completed-tensor-product}，已知
$A \widehat{\otimes}_B C$ 可容许。此外，该引理的证明表明：
% P10-E0298
当 $I \subset A$ 和 $J \subset C$ 是定义理想时，
$I(A \widehat{\otimes}_B C) + J(A \widehat{\otimes}_B C)$ 的闭包
$K \subset A \widehat{\otimes}_B C$ 是定义理想。
于是只须证明，对所有 $n \geq 1$，$K^n$ 的闭包都是开的。
由于理想 $K^n$ 包含
$I^n(A \widehat{\otimes}_B C) + J^n(A \widehat{\otimes}_B C)$，
而 $I^n$ 在 $A$ 中的闭包和 $J^n$ 在 $C$ 中的闭包都是开的，
可见 $K^n$ 在 $A \widehat{\otimes}_B C$ 中的闭包是开的。
\end{proof}






\section{性质的下降}
\label{formal-spaces-section-taut-descent}

\noindent
本节考虑如下情形：
\begin{enumerate}
\item $\varphi : A \to B$ 是线性拓扑化拓扑环的连续映射；
\item $\varphi$ 是 taut 的；
\item 对每个开理想 $I \subset A$，若 $J \subset B$ 表示 $IB$ 的闭包，
则映射 $A/I \to B/J$ 忠实平坦。
\end{enumerate}
我们将证明，在此情形下，$B$ 的性质由 $A$ 继承。

\begin{lemma}
\label{formal-spaces-lemma-taut-descent-countable}
在上述情形中，若 $B$ 有可数的开理想基本系，则 $A$ 有可数的开理想基本系。
\end{lemma}

\begin{proof}
选取一族开理想基本系
$B \supset J_1 \supset J_2 \supset \ldots$。由 $\varphi$ 的 taut 性，
对每个 $n$，可以找到开理想 $I_n$，使 $J_n \supset I_nB$。
我们断言 $I_n$ 是 $A$ 的一族开理想基本系。事实上，设 $I \subset A$ 开。
由于 $\varphi$ 是 taut 的，$IB$ 的闭包是开的，因此对充分大的某个 $n$ 包含 $J_n$。
所以 $I_nB \subset IB$。令 $J$ 为 $IB$ 在 $B$ 中的闭包。
因为 $A/I \to B/J$ 忠实平坦，所以它是单射。又由于
$I_nB \subset IB \subset J$，映射 $I_n \to A/I \to B/J$ 为零，
故 $I_n \to A/I$ 为零。因此 $I_n \subset I$，结论成立。
\end{proof}

\begin{lemma}
\label{formal-spaces-lemma-taut-descent-weakly-admissible}
在上述情形中，若 $B$ 弱预可容许，则 $A$ 弱预可容许。
\end{lemma}

\begin{proof}
设 $J \subset B$ 是弱定义理想。取开理想 $I \subset A$ 使 $IB \subset J$。
为证明 $I$ 是弱定义理想，须证明每个 $f \in I$ 都拓扑幂零。
设 $I' \subset A$ 是开理想，以 $J' \subset B$ 表示 $I'B$ 的闭包。
则 $A/I' \to B/J'$ 忠实平坦，从而是单射。因此，为证明 $f^n \in I'$，
只须证明 $\varphi(f)^n \in J'$。当 $n \gg 0$ 时这成立，因为
$\varphi(f) \in J$，理想 $J$ 是 $B$ 的弱定义理想，而 $J'$ 在 $B$ 中开。
\end{proof}

\begin{lemma}
\label{formal-spaces-lemma-taut-descent-admissible}
在上述情形中，若 $B$ 预可容许，则 $A$ 预可容许。
\end{lemma}

\begin{proof}
% P10-E0299
设 $J \subset B$ 是弱定义理想。取开理想 $I \subset A$ 使 $IB \subset J$。
设 $I' \subset A$ 是开理想。为证明 $I$ 是定义理想，须证明当 $n \gg 0$ 时
$I^n \subset I'$。以 $J' \subset B$ 表示 $I'B$ 的闭包。
则 $A/I' \to B/J'$ 忠实平坦，从而是单射。因此，为证明
$I^n \subset I'$，只须证明 $\varphi(I)^n \subset J'$。
当 $n \gg 0$ 时这成立，因为 $\varphi(I) \subset J$，理想 $J$ 是 $B$ 的定义理想，
而 $J'$ 在 $B$ 中开。
\end{proof}

\begin{lemma}
\label{formal-spaces-lemma-taut-descent-weakly-adic}
在上述情形中，若 $B$ 弱预-adic，则 $A$ 弱预-adic。
\end{lemma}

\begin{proof}
我们使用引理 \ref{formal-spaces-lemma-weakly-adic} 给出的弱预-adic 环刻画而不再说明。
由引理 \ref{formal-spaces-lemma-taut-descent-admissible}，拓扑环 $A$ 预可容许。
设 $I \subset A$ 是定义理想。固定 $n \geq 1$。为证明引理，须证明
$I^n$ 的闭包是开的。设 $I_\lambda \subset A$ 是一族开理想基本系。
以 $J \subset B$（分别地，$J_\lambda \subset B$）表示
$IB$（分别地，$I_\lambda B$）的闭包。由于 $B$ 弱预-adic，$J^n$ 的闭包是开的。
所以存在 $\lambda$，使得
$$
J_\lambda \subset \bigcap\nolimits_\mu (J^n + J_\mu)
$$
因为由引理 \ref{formal-spaces-lemma-closed}，右侧是 $J^n$ 的闭包。
这表示 $J_\lambda$ 在 $B/J_\mu$ 中的像包含在 $J^n$ 在 $B/J_\mu$ 中的像内。
注意，$J^n$ 在 $B/J_\mu$ 中的像等于 $I^nB$ 在 $B/J_\mu$ 中的像
（因为 $J$ 的每个元素模 $J_\mu$ 都同余于 $IB$ 的一个元素）。
由于 $A/I_\mu \to B/J_\mu$ 忠实平坦，并且
$I_\lambda B \subset J_\lambda$，可知 $I_\lambda$ 在 $A/I_\mu$ 中的像
包含在 $I^n$ 的像内。因此 $I_\lambda$ 包含在 $I^n$ 的闭包中，证明完成。
\end{proof}

\begin{lemma}
\label{formal-spaces-lemma-taut-descent-adic-star}
在上述情形中，若 $B$ 是 adic 的且有有限生成的定义理想，而 $A$ 完备，
则 $A$ 是 adic 的且有有限生成的定义理想。
\end{lemma}

\begin{proof}
由引理 \ref{formal-spaces-lemma-taut-descent-weakly-adic}，已知 $A$ 弱 adic，
从而当然可容许（并用引理 \ref{formal-spaces-lemma-weakly-adic} 看出 adic 环是弱 adic 的）。
设 $I \subset A$ 是定义理想，$J \subset B$ 是有限生成的定义理想。
由于 $IB$ 的闭包是开的，可以找到 $n > 0$，使 $J^n$ 包含在 $IB$ 的闭包中。
所以用 $J^n$ 替换 $J$ 后，可以假设 $J$ 是包含在 $IB$ 闭包中的有限生成定义理想。
由引理 \ref{formal-spaces-lemma-closed}，这当然蕴含
$$
J \subset IB + J^2
$$
考虑有限生成 $A$-模 $M = (J + IB)/IB$。所显示的方程表明 $JM = M$。
由引理 \ref{formal-spaces-lemma-weakly-admissible-henselian}（例如）可见，$J$ 包含在
$B$ 的 Jacobson 根中。所以由 Nakayama 引理，更确切地说由《代数》之引理
\ref{algebra-lemma-NAK} 的 (2)，得到 $M = 0$。因此 $J \subset IB$。

\medskip\noindent
因为 $J$ 有限生成，可以找到有限生成理想 $I' \subset I$，使 $J \subset I'B$。
因为 $A \to B$ 连续、$J \subset B$ 开，并且 $I$ 是定义理想，
可以找到 $n > 0$ 使 $I^nB \subset J$。令 $J_{n + 1} \subset B$
为 $I^{n + 1}B$ 的闭包。我们有
$$
I^n \cdot (B/J_{n + 1}) \subset J \cdot (B/J_{n + 1}) \subset
I' \cdot (B/J_{n + 1})
$$
由于 $A/I^{n + 1} \to B/J_{n + 1}$ 忠实平坦，这蕴含
$I^n \cdot (A/I^{n + 1}) \subset I' \cdot (A/I^{n + 1})$，
进而意味着
$$
I^n \subset I' + I^{n + 1}
$$
这蕴含对所有 $k \geq 1$ 有 $I^n \subset I' + I^{n + k}$，
进而蕴含对所有 $k,m \geq 1$ 有
$I^{nm} \subset (I')^m + I^{nm + k}$。这说明 $(I')^m$ 的闭包包含 $I^{nm}$。
由于 $A$ 弱 adic，$I^{nm}$ 的闭包是开的，故
% P10-E0300
对所有 $m$，$(I')^m$ 的闭包都是开的。由于这些闭包组成 $A$ 的一族开理想基本系
（因为 $I^n$ 的闭包也具有此性质），结论由引理
\ref{formal-spaces-lemma-weakly-adic-finite-generation-bis} 得出。
\end{proof}








\section{仿射形式代数空间}
\label{formal-spaces-section-affine-formal-algebraic-spaces}

\noindent
本节引入仿射形式代数空间。它们事实上与 \cite{BVGD} 中所谓的仿射形式概形相同。
不过，为避免与 \cite{EGA} 所定义的仿射形式概形概念混淆，
我们将把它们称为仿射形式代数空间。

\medskip\noindent
回忆，概形的增厚是一个在底层拓扑空间上诱导满射的闭浸入，
见《态射续篇》之定义 \ref{more-morphisms-definition-thickening}。

\begin{definition}
\label{formal-spaces-definition-affine-formal-algebraic-space}
设 $S$ 是概形。若存在
\begin{enumerate}
\item 一个有向集 $\Lambda$；
\item $(\Sch/S)_{fppf}$ 中位于 $\Lambda$ 上的系统
$(X_\lambda, f_{\lambda \mu})$，其中
\begin{enumerate}
\item 每个 $X_\lambda$ 都仿射；
\item 每个 $f_{\lambda \mu} : X_\lambda \to X_\mu$ 都是增厚，
\end{enumerate}
\end{enumerate}
使得作为 fppf 层有
$$
X \cong \colim_{\lambda \in \Lambda} X_\lambda
$$
并且 $X$ 满足一个集合论条件（见注 \ref{formal-spaces-remark-set-theoretic}），
则称 $(\Sch/S)_{fppf}$ 上的层 $X$ 是{\it 仿射形式代数空间}。
$S$ 上仿射形式代数空间的{\it 态射}是层的映射。
\end{definition}

\noindent
注意，系统 $(X_\lambda, f_{\lambda \mu})$ 不是数据的一部分。
设 $U$ 是 $S$ 上的拟紧概形。由于转移映射都是单态射，由《站点》之引理
\ref{sites-lemma-directed-colimits-sections} 可见
$$
X(U) = \colim X_\lambda(U)
$$
因此，定义中的余极限所内含的 fppf 层化其实是 Zariski 层化，
而它对拟紧概形不起任何作用。

\begin{lemma}
\label{formal-spaces-lemma-diagonal-affine-formal-algebraic-space}
设 $S$ 是概形。若 $X$ 是 $S$ 上的仿射形式代数空间，则对角态射
$\Delta : X \to X \times_S X$ 可表示且为闭浸入。
\end{lemma}

\begin{proof}
设给定 $U \to X$ 和 $V \to X$，其中 $U,V$ 是 $S$ 上的概形。
来证明 $U \times_X V$ 可表示。按定义
\ref{formal-spaces-definition-affine-formal-algebraic-space} 写
$X = \colim X_\lambda$。上述讨论表明，在 $U$ 和 $V$ 上 Zariski 局部地，
% P10-E0301
这些态射分解经过某个 $X_\lambda$。此时
$U \times_X V = U \times_{X_\lambda} V$，它是概形。
所以由《空间》之引理 \ref{spaces-lemma-representable-diagonal}，对角可表示。
给定 $(a,b) : W \to X \times_S X$，其中 $W$ 是 $S$ 上的概形，考虑映射
$X \times_{\Delta, X \times_S X, (a, b)} W \to W$。
与前面一样，在 $W$ 上局部地，态射 $a$ 和 $b$ 映入某个 $\lambda$ 对应的仿射概形
$X_\lambda$，于是得到态射
$X_\lambda
\times_{\Delta_\lambda, X_\lambda \times_S X_\lambda, (a, b)} W \to W$。
这是 $\Delta_\lambda : X_\lambda \to X_\lambda \times_S X_\lambda$
的基变换，而后者是闭浸入，因为 $X_\lambda \to S$ 分离
（$X_\lambda$ 仿射）。所以 $X \to X \times_S X$ 是闭浸入。
\end{proof}

\noindent
概形态射 $X \to X'$ 是增厚，当且仅当它是闭浸入并在底层点集上诱导满射，
见（《态射续篇》之定义 \ref{more-morphisms-definition-thickening}）。
因此，增厚这一性质在任意基变换下保持，并且在目标上是 fpqc 局部的，
见《空间》之第 \ref{spaces-section-lists} 节。所以《空间》之定义
\ref{spaces-definition-relative-representable-property} 适用于“增厚”，
而我们知道 $(\Sch/S)_{fppf}$ 上预层的可表示变换 $F \to G$ 是增厚是什么意思。
我们指出，在 $F$ 和 $G$ 是代数空间时，这与《空间态射续篇》之定义
\ref{spaces-more-morphisms-definition-thickening} 中的增厚定义并不冲突。

\begin{lemma}
\label{formal-spaces-lemma-covering-by-thickenings}
设 $X_\lambda$（$\lambda \in \Lambda$）和 $X = \colim X_\lambda$
如定义 \ref{formal-spaces-definition-affine-formal-algebraic-space}。则
$X_\lambda \to X$ 可表示且为增厚。
\end{lemma}

\begin{proof}
由《空间》之第 \ref{spaces-section-representable} 节和第
\ref{spaces-section-representable-properties} 节的讨论，该陈述有意义。
由引理 \ref{formal-spaces-lemma-diagonal-affine-formal-algebraic-space}，态射
$X_\lambda \to X$ 可表示。给定 $U \to X$，其中 $U$ 是概形，
定义 \ref{formal-spaces-definition-affine-formal-algebraic-space} 后的讨论表明，
在 $U$ 上 Zariski 局部地，该态射分解经过某个满足
$\lambda \leq \mu$ 的 $X_\mu$。此时
$U \times_X X_\lambda = U \times_{X_\mu} X_\lambda$，
所以 $U \times_X X_\lambda \to U$ 是增厚
$X_\lambda \to X_\mu$ 的基变换。
\end{proof}

\begin{lemma}
\label{formal-spaces-lemma-factor-through-thickening}
设 $X_\lambda$（$\lambda \in \Lambda$）和 $X = \colim X_\lambda$
如定义 \ref{formal-spaces-definition-affine-formal-algebraic-space}。
若 $Y$ 是 $S$ 上的拟紧代数空间，则任意态射 $Y \to X$ 都分解经过某个
$X_\lambda$。
\end{lemma}

\begin{proof}
选取一个仿射概形 $V$ 和满的 étale 态射 $V \to Y$。
由定义 \ref{formal-spaces-definition-affine-formal-algebraic-space} 后的讨论，
复合 $V \to Y \to X$ 分解经过某个 $X_\lambda$。
由于 $V \to Y$ 是层的满射，结论成立。
\end{proof}

\begin{lemma}
\label{formal-spaces-lemma-characterize-affine-formal-algebraic-space}
设 $S$ 是概形，$X$ 是 $(\Sch/S)_{fppf}$ 上的层。
则 $X$ 是仿射形式代数空间，当且仅当下列条件成立：
\begin{enumerate}
\item 任意态射 $U \to X$，其中 $U$ 是 $S$ 上的仿射概形，
都分解经过一个可表示且为增厚的态射 $T \to X$，其中 $T$ 是 $S$ 上的仿射概形；
\item 满足如注 \ref{formal-spaces-remark-set-theoretic} 中的集合论条件。
\end{enumerate}
\end{lemma}

\begin{proof}
由引理 \ref{formal-spaces-lemma-covering-by-thickenings} 和
\ref{formal-spaces-lemma-factor-through-thickening}，仿射形式代数空间满足 (1) 和 (2)。
为证明逆命题，可以假设 $X$ 非空。令 $\Lambda$ 为如下范畴：
对象是可表示且为增厚的态射 $T \to X$，其中 $T$ 是 $S$ 上的仿射概形。
该范畴有向。因为 $X$ 非空，$\Lambda$ 至少含一个对象。
若 $T \to X$ 和 $T' \to X$ 属于 $\Lambda$，则可以把
$T \amalg T' \to X$ 分解经过 $\Lambda$ 中的 $T'' \to X$。
$\Lambda$ 的任意两个对象之间至多有唯一箭头。因此 $\Lambda$ 是有向集，且按假设
$X = \colim_{T \to X\text{ 属于 }\Lambda} T$。
为完成证明，须说明 $\Lambda$ 中任意箭头 $T \to T'$ 都是增厚。
这是因为 $T' \to X$ 是层的单态射，所以
$T = T \times_{T'} T' = T \times_X T'$；因而态射 $T \to T'$
等于投影 $T \times_X T' \to T'$，它是增厚，因为 $T \to X$ 是增厚。
\end{proof}

\noindent
一般的仿射形式代数空间 $X$ 不一定有足够多的函数来分离点（例如）。
见《例》之第 \ref{examples-section-affine-formal-algebraic-space} 节。
为刻画确实具有这种性质的空间，给出下列引理。

\begin{lemma}
\label{formal-spaces-lemma-mcquillan-affine-formal-algebraic-space}
设 $S$ 是概形。设 $X$ 是 $(\Sch/S)_{fppf}$ 上的 fppf 层，
满足注 \ref{formal-spaces-remark-set-theoretic} 的集合论条件。下列条件等价：
\begin{enumerate}
\item 存在 $S$ 上的弱可容许拓扑环 $A$（见注 \ref{formal-spaces-remark-mcquillan}），使得
$X = \colim_{I \subset A\text{ 为弱定义理想}} \Spec(A/I)$；
\item $X$ 是仿射形式代数空间，并且存在 $S$-代数 $A$ 和映射
$X \to \Spec(A)$，使得对任意闭浸入 $T \to X$，其中 $T$ 是仿射概形，
复合 $T \to \Spec(A)$ 是闭浸入；
\item $X$ 是仿射形式代数空间，并且存在 $S$-代数 $A$ 和映射
$X \to \Spec(A)$，使得对任意闭浸入 $T \to X$，其中 $T$ 是概形，
复合 $T \to \Spec(A)$ 是闭浸入；
\item $X$ 是仿射形式代数空间，并且对定义
\ref{formal-spaces-definition-affine-formal-algebraic-space} 中某个表示
$X = \colim X_\lambda$，各投影
$\lim \Gamma(X_\lambda, \mathcal{O}_{X_\lambda})
\to \Gamma(X_\lambda, \mathcal{O}_{X_\lambda})$ 都是满射；
\item $X$ 是仿射形式代数空间，并且对定义
\ref{formal-spaces-definition-affine-formal-algebraic-space} 中任意表示
$X = \colim X_\lambda$，各投影
$\lim \Gamma(X_\lambda, \mathcal{O}_{X_\lambda})
\to \Gamma(X_\lambda, \mathcal{O}_{X_\lambda})$ 都是满射。
\end{enumerate}
此外，该弱可容许拓扑环是
$A = \lim \Gamma(X_\lambda, \mathcal{O}_{X_\lambda})$，
赋予其极限拓扑；弱定义理想恰好分类那些可表示且为增厚的态射 $T \to X$。
\end{lemma}

\begin{proof}
(5) 推出 (4) 是显然的。

\medskip\noindent
对定义 \ref{formal-spaces-definition-affine-formal-algebraic-space} 中的
$X = \colim X_\lambda$ 假设 (4)。令
$A = \lim \Gamma(X_\lambda, \mathcal{O}_{X_\lambda})$。
设 $T \to X$ 是闭浸入，其中 $T$ 是概形
（注意，由引理 \ref{formal-spaces-lemma-diagonal-affine-formal-algebraic-space}，
$T \to X$ 可表示）。由于 $X_\lambda \to X$ 是增厚，
$X_\lambda \times_X T \to T$ 也是增厚。另一方面，
$X_\lambda \times_X T \to X_\lambda$ 是闭浸入，所以
$X_\lambda \times_X T$ 仿射。于是由《极限》之命题
\ref{limits-proposition-affine}，$T$ 仿射。再由引理
\ref{formal-spaces-lemma-factor-through-thickening}，$T \to X$ 分解经过某个
$X_\lambda$。因此
$A \to \Gamma(X_\lambda, \mathcal{O}) \to \Gamma(T, \mathcal{O})$
是满射。由此可见 (3) 成立。

\medskip\noindent
(3) 推出 (2) 是显然的。

\medskip\noindent
对 $A$ 和 $X \to \Spec(A)$ 假设 (2)。按定义
\ref{formal-spaces-definition-affine-formal-algebraic-space} 写
$X = \colim X_\lambda$。则按 (2)，
$A_\lambda = \Gamma(X_\lambda, \mathcal{O})$ 是 $A$ 的商。
所以 $A^\wedge = \lim A_\lambda$ 是完备拓扑环，见《代数续篇》之第
\ref{more-algebra-section-topological-ring} 节的讨论。因为
$A \to A_\lambda$ 是满射，各映射 $A^\wedge \to A_\lambda$ 也是满射。
我们断言，对任意 $\lambda$，$A^\wedge \to A_\lambda$ 的核
$I_\lambda \subset A^\wedge$ 都是弱定义理想。事实上，按极限拓扑的定义，它是开的。
若 $f \in I_\lambda$，则对任意 $\mu \in \Lambda$，$f$ 在
$A_\mu$ 中的像在 $X_\mu$ 各点的所有剩余域中都为零。因此它是
$A_\mu$ 的幂零元素，故某个幂 $f^n \in I_\mu$。所以
% P10-E0302
当 $n \to 0$ 时 $f^n \to 0$。因此 $A^\wedge$ 弱可容许。
最后，设 $I \subset A^\wedge$ 是弱定义理想。则
$I \subset A^\wedge$ 开，所以存在某个 $\lambda$ 使
$I \supset I_\lambda$。于是得到态射
$\Spec(A^\wedge/I) \to \Spec(A_\lambda) \to X$。
因而 $X = \colim \Spec(A^\wedge/I)$，这里余极限遍历所有弱定义理想。
所以 (1) 成立。

\medskip\noindent
假设 (1)。此时 $X$ 显然是仿射形式代数空间。任取定义
\ref{formal-spaces-definition-affine-formal-algebraic-space} 中的表示
$X = \colim X_\lambda$。对每个 $\lambda$，由引理
\ref{formal-spaces-lemma-factor-through-thickening}，可找到弱定义理想
$I \subset A$，使 $X_\lambda \to X$ 分解经过
$\Spec(A/I) \to X$。于是
$X_\lambda = \Spec(A/I_\lambda)$，其中 $I \subset I_\lambda$。
反过来，对任意弱定义理想 $I \subset A$，态射
$\Spec(A/I) \to X$ 分解经过某个 $X_\lambda$，即
$I_\lambda \subset I$。所以，每个 $I_\lambda$ 都是弱定义理想，
且它们组成弱定义理想集合的一个共尾子集。因此
$A = \lim A/I = \lim A/I_\lambda$，可见 (5) 成立，并且还有
$A = \lim \Gamma(X_\lambda, \mathcal{O}_{X_\lambda})$。
\end{proof}

\noindent
有了这个引理，可以作如下定义。

\begin{definition}
\label{formal-spaces-definition-types-affine-formal-algebraic-space}
设 $S$ 是概形，$X$ 是 $S$ 上的仿射形式代数空间。
若 $X$ 满足引理 \ref{formal-spaces-lemma-mcquillan-affine-formal-algebraic-space}
中的等价条件，则称 $X$ 是 {\it McQuillan} 的。
设 $A$ 是与 $X$ 相伴的弱可容许拓扑环。我们称
\begin{enumerate}
\item 若 $X$ 是 McQuillan 的且 $A$ 可容许
（《代数续篇》之定义 \ref{more-algebra-definition-topological-ring}），
则 $X$ 是{\it 经典的}；
\item 若 $X$ 是 McQuillan 的且 $A$ 弱 adic
（定义 \ref{formal-spaces-definition-weakly-adic}），则 $X$ {\it 弱 adic}；
\item 若 $X$ 是 McQuillan 的且 $A$ adic
（《代数续篇》之定义 \ref{more-algebra-definition-topological-ring}），
则 $X$ {\it adic}；
\item 若 $X$ 是 McQuillan 的、$A$ adic 且 $A$ 有有限生成的定义理想，
则 $X$ {\it adic*}；
\item 若 $X$ 是 McQuillan 的，且 $A$ 既 Noether 又 adic，
则 $X$ 是 {\it Noether} 的。
\end{enumerate}
\end{definition}

\noindent
在 \cite{Fujiwara-Kato} 中，adic 拓扑环 $A$ 包含有限生成定义理想这一性质
被称为“有限理想型”。给定仿射形式代数空间 $X$，定义中引入的概念之间有如下蕴含：
$$
\xymatrix{
X\text{ 是 Noether 的} \ar@{=>}[r] &
X\text{ 是 adic* 的} \ar@{=>}[r] &
X\text{ 是 adic 的} \ar@{=>}[lld] \\
X\text{ 是弱 adic 的} \ar@{=>}[r] &
X\text{ 是经典的} \ar@{=>}[r] &
X\text{ 是 McQuillan 的}
}
$$
见第 \ref{formal-spaces-section-weakly-adic} 节的讨论；精确陈述见引理
\ref{formal-spaces-lemma-implications-between-types}。

\begin{remark}
\label{formal-spaces-remark-compare-with-affine-formal-schemes}
经典仿射形式代数空间对应于 EGA（\cite{EGA}）所考虑的仿射形式概形。
为解释这一点，假设基概形是 $\Spec(\mathbf{Z})$。设
$\mathfrak X = \text{Spf}(A)$ 是仿射形式概形，并令
$h_\mathfrak X$ 为引理 \ref{formal-spaces-lemma-fully-faithful} 中的点函子。
则 $h_\mathfrak X = \colim h_{\Spec(A/I)}$，其中余极限遍历可容许拓扑环
$A$ 的定义理想集合。在仿射概形上取值时，这由式
(\ref{formal-spaces-equation-morphisms-affine-formal-schemes}) 得出；只须在仿射概形上核验，
因为两边都是 fppf 层，见引理 \ref{formal-spaces-lemma-formal-scheme-sheaf-fppf}。
因此 $h_\mathfrak X$ 是仿射形式代数空间。事实上，由定义
\ref{formal-spaces-definition-types-affine-formal-algebraic-space}，它是经典仿射形式代数空间。
所以引理 \ref{formal-spaces-lemma-fully-faithful} 告诉我们，仿射形式概形范畴
等价于经典仿射形式代数空间范畴。
\end{remark}

\noindent
作出上述与仿射形式概形的联系后，作如下定义显得自然。

\begin{definition}
\label{formal-spaces-definition-affine-formal-spectrum}
设 $S$ 是概形。设 $A$ 是 $S$ 上的弱可容许拓扑环，见定义
\ref{formal-spaces-definition-weakly-admissible}\footnote{经典情形见《代数续篇》之定义
\ref{more-algebra-definition-topological-ring}，差异的讨论见注
\ref{formal-spaces-remark-mcquillan}。}。$A$ 的{\it 形式谱}是仿射形式代数空间
$$
\text{Spf}(A) = \colim \Spec(A/I)
$$
其中余极限遍历 $A$ 的弱定义理想集合，并在范畴
$\Sh((\Sch/S)_{fppf})$ 中取。
\end{definition}

\noindent
这样的形式谱按构造是 McQuillan 的；反过来，每个 McQuillan 仿射形式代数空间
都同构于某个形式谱。需要指出的是，在我们的理论中，确实存在不能写成任何弱可容许拓扑环
的形式谱的仿射形式代数空间。依照 \cite{Yasuda}，可以把 $S$-pro-环定义为
$S$-代数范畴中的 pro-对象，见《范畴》之注
\ref{categories-remark-pro-category}。这样，每个 $S$ 上的仿射形式代数空间
都会是某个 $S$-pro-环的形式谱。我们不这样做，而将直接处理相应的仿射形式代数空间。

\medskip\noindent
形式谱的构造是函子性的。为解释这一点，设
% P10-E0303
$\varphi : B \to A$ 是 $S$ 上弱可容许拓扑环的连续映射。则
$$
\text{Spf}(\varphi) : \text{Spf}(B) \to \text{Spf}(A)
$$
是唯一的仿射形式代数空间态射，使得对所有满足
$\varphi(I) \subset J$ 的弱定义理想 $I \subset A$ 和 $J \subset B$，图表
$$
\xymatrix{
\Spec(B/J) \ar[d] \ar[r] & \Spec(A/I) \ar[d] \\
\text{Spf}(B) \ar[r] & \text{Spf}(A)
}
$$
交换。由于 $\varphi$ 的连续性蕴含：对每个弱定义理想 $J \subset B$，
存在具有所需性质的弱定义理想 $I \subset A$，可见所要求的交换性唯一决定并定义
$\text{Spf}(\varphi)$。

\begin{lemma}
\label{formal-spaces-lemma-morphism-between-formal-spectra}
设 $S$ 是概形，$A,B$ 是 $S$ 上的弱可容许拓扑环。
$S$ 上仿射形式代数空间的任意态射
$f : \text{Spf}(B) \to \text{Spf}(A)$ 都等于
$\text{Spf}(f^\sharp)$，其中
$f^\sharp : A \to B$ 是唯一的连续 $S$-代数映射。
\end{lemma}

\begin{proof}
设 $f : \text{Spf}(B) \to \text{Spf}(A)$ 如引理所述。
设 $J \subset B$ 是弱定义理想。由引理
\ref{formal-spaces-lemma-factor-through-thickening}，存在弱定义理想 $I \subset A$，使复合
$\Spec(B/J) \to \text{Spf}(B) \to \text{Spf}(A)$
分解经过 $\Spec(A/I)$。由《概形》之引理
\ref{schemes-lemma-morphism-into-affine}，得到 $S$-代数映射
$A/I \to B/J$。这些映射随 $J$ 变化相容，并定义映射
$f^\sharp : A \to B$。该映射连续，因为对每个弱定义理想
$J \subset B$，都存在弱定义理想 $I \subset A$ 使
$f^\sharp(I) \subset J$。等式 $f = \text{Spf}(f^\sharp)$
由我们对组成 $f^\sharp$ 的环映射 $A/I \to B/J$ 的选择成立。
\end{proof}

\begin{lemma}
\label{formal-spaces-lemma-presentation-representable}
设 $S$ 是概形，$f : X \to Y$ 是 $(\Sch/S)_{fppf}$ 上预层的映射。
若 $X$ 是仿射形式代数空间，而 $f$ 可由代数空间表示且局部拟有限，
则 $f$ 可表示（由概形表示）。
\end{lemma}

\begin{proof}
设 $T$ 是 $S$ 上的概形，且给定映射 $T \to Y$。须证明代数空间
$X \times_Y T$ 是概形。按定义
\ref{formal-spaces-definition-affine-formal-algebraic-space} 写
$X = \colim X_\lambda$。设 $W \subset X \times_Y T$ 是拟紧开子空间。
投影 $X \times_Y T \to X$ 在 $W$ 上的限制分解经过某个 $X_\lambda$。于是
$$
W \to X_\lambda \times_S T
$$
是单态射（从而分离）且局部拟有限（因为按对 $X \to Y$ 的假设，
$W \to X \times_Y T \to T$ 局部拟有限；见《空间的态射》之引理
\ref{spaces-morphisms-lemma-permanence-quasi-finite}）。
所以由《空间的态射》之命题
\ref{spaces-morphisms-proposition-locally-quasi-finite-separated-over-scheme}，
$W$ 是概形。因此由《空间的性质》之引理
\ref{spaces-properties-lemma-subscheme}，$X \times_Y T$ 是概形。
\end{proof}






\section{可数指标的仿射形式代数空间}
\label{formal-spaces-section-countably-indexed}

\noindent
这类仿射形式代数空间由下列引理描述。

\begin{lemma}
\label{formal-spaces-lemma-countable-affine-formal-algebraic-space}
设 $S$ 是概形，$X$ 是 $S$ 上的仿射形式代数空间。下列条件等价：
\begin{enumerate}
\item 存在 $S$ 上仿射概形的一列增厚
$X_1 \to X_2 \to X_3 \to \ldots$，使得 $X = \colim X_n$；
\item 存在定义 \ref{formal-spaces-definition-affine-formal-algebraic-space}
中的一种选择 $X = \colim X_\lambda$，使得 $\Lambda$ 可数。
\end{enumerate}
\end{lemma}

\begin{proof}
这来自如下观察：可数有向集有一个同构于 $(\mathbf{N}, \geq)$ 的共尾子集。
见《代数》之引理 \ref{algebra-lemma-ML-limit-nonempty} 的证明。
\end{proof}

\begin{definition}
\label{formal-spaces-definition-countable}
设 $S$ 是概形，$X$ 是 $S$ 上的仿射形式代数空间。
若引理 \ref{formal-spaces-lemma-countable-affine-formal-algebraic-space} 的等价条件成立，
则称 $X$ 是{\it 可数指标的}。
\end{definition}

\noindent
用 \cite{BVGD} 的语言，这表述为 $X$ 是 $\aleph_0$-ind 概形。

\begin{lemma}
\label{formal-spaces-lemma-implications-between-types}
设 $X$ 是概形 $S$ 上的仿射形式代数空间。
\begin{enumerate}
\item 若 $X$ 是 Noether 的，则 $X$ adic*。
\item 若 $X$ adic*，则 $X$ adic。
\item 若 $X$ adic，则 $X$ 弱 adic。
\item 若 $X$ 弱 adic，则 $X$ 经典。
\item 若 $X$ 弱 adic，则 $X$ 是可数指标的。
\item 若 $X$ 是可数指标的，则 $X$ 是 McQuillan 的。
\end{enumerate}
\end{lemma}

\begin{proof}
% P10-E0304
把 $X$ 写成 $X = \text{Spf}(A)$，其中 $A$ 是弱可容许的（从而完备的）
线性拓扑环。利用第 \ref{formal-spaces-section-weakly-adic} 节讨论的这类环各种类型之间的蕴含，
即可得到 (1)、(2)、(3) 和 (4)。

\medskip\noindent
(5) 的证明。按定义，存在弱 adic 拓扑环 $A$，使
$X = \colim \Spec(A/I)$，其中余极限遍历 $A$ 的定义理想。
由于 $A$ 弱 adic，特别地，
% P10-E0305
存在一族可数的开理想基本系 $I_\lambda$，见定义
\ref{formal-spaces-definition-weakly-adic}。于是按 $\text{Spf}(A)$ 的定义，
% P10-E0306
$X = \colim \Spec(A/I_n)$。所以 $X$ 是可数指标的。

\medskip\noindent
(6) 的证明。把 $X = \colim X_n$ 写成 $S$ 上仿射概形的一列增厚
$X_1 \to X_2 \to X_3 \to \ldots$ 的余极限。则
$$
A = \lim \Gamma(X_n, \mathcal{O}_{X_n})
$$
满射到每个 $\Gamma(X_n, \mathcal{O}_{X_n})$，因为态射
$X_n \to X_{n + 1}$ 是闭浸入，所以转移映射是满射。因此 $X$ 是 McQuillan 的。
\end{proof}

\begin{lemma}
\label{formal-spaces-lemma-countably-indexed}
设 $S$ 是概形，$X$ 是 $(\Sch/S)_{fppf}$ 上的预层。下列条件等价：
\begin{enumerate}
\item $X$ 是可数指标的仿射形式代数空间；
\item $X = \text{Spf}(A)$，其中 $A$ 是有可数 $0$ 邻域基本系的弱可容许拓扑
$S$-代数；
\item $X = \text{Spf}(A)$，其中 $A$ 是弱可容许拓扑 $S$-代数，
且有一族弱定义理想基本系
$A \supset I_1 \supset I_2 \supset I_3 \supset \ldots$；
\item $X = \text{Spf}(A)$，其中 $A$ 是完备拓扑 $S$-代数，
它的一族 $0$ 的开邻域基本系由理想的可数序列
$A \supset I_1 \supset I_2 \supset I_3 \supset \ldots$ 给出，
并且 $I_n/I_{n + 1}$ 局部幂零；
\item $X = \text{Spf}(A)$，其中 $A = \lim B/J_n$ 带极限拓扑，
而 $B \supset J_1 \supset J_2 \supset J_3 \supset \ldots$
是 $S$-代数 $B$ 中的一列理想，且 $J_n/J_{n + 1}$ 局部幂零。
\end{enumerate}
\end{lemma}

\begin{proof}
假设 (1)。由引理 \ref{formal-spaces-lemma-implications-between-types}，
可写 $X = \text{Spf}(A)$，其中 $A$ 是弱可容许拓扑 $S$-代数。
对引理 \ref{formal-spaces-lemma-countable-affine-formal-algebraic-space} (1)
中的任意表示 $X = \colim X_n$，可见
$A = \lim A_n$，其中 $X_n = \Spec(A_n)$，并且对某个弱定义理想
$I_n \subset A$ 有 $A_n = A/I_n$。这来自引理
\ref{formal-spaces-lemma-mcquillan-affine-formal-algebraic-space} 的最后一个陈述，
该陈述还蕴含 $\{I_n\}$ 是一族 $0$ 的开邻域基本系。
因此有弱定义理想的序列
$$
A \supset I_1 \supset I_2 \supset I_3 \supset \ldots
$$
且 $A = \lim A/I_n$。由此可见，条件 (1) 推出条件 (2)--(5) 中的每一个。

\medskip\noindent
假设 (5)。首先注意，$A = \lim B/J_n$ 上的极限拓扑是完备的线性拓扑，
见《代数续篇》之第 \ref{more-algebra-section-topological-ring} 节。
若 $f \in A$ 在 $B/J_1$ 中映为零，则由于其在 $J_1/J_2$ 中的像幂零，
某个幂在 $B/J_2$ 中映为零；再取一个幂，它在 $J_2/J_3$ 中映为零，等等。
由此可见，开理想 $\Ker(A \to B/J_1)$ 是弱定义理想。所以 $A$ 弱可容许。
由此可见 (5) 推出 (2)。

\medskip\noindent
取 $B = A$，显然 (4) 是 (5) 的特殊情形。显然 (3) 是 (2) 的特殊情形。

\medskip\noindent
假设 $A$ 如 (2) 所述。设 $E_n$ 是 $A$ 中一族可数的 $0$ 邻域基本系。
由于 $A$ 是弱可容许拓扑环，可以找到开理想 $I_n \subset E_n$。
还可以选取弱定义理想 $J \subset A$。则 $J \cap I_n$ 是 $A$
的一族弱定义理想基本系，并得到
$X = \text{Spf}(A) = \colim \Spec(A/(J \cap I_n))$，
这表明 $X$ 是可数指标的仿射形式代数空间。
\end{proof}

\begin{lemma}
\label{formal-spaces-lemma-characterize-noetherian-affine}
设 $S$ 是概形，$X$ 是仿射形式代数空间。下列条件等价：
\begin{enumerate}
\item $X$ 是 Noether 的；
\item $X$ adic*，并且对每个闭浸入 $T \to X$，其中 $T$ 是概形，
$T$ 都是 Noether 的；
\item $X$ adic*，并且对定义
\ref{formal-spaces-definition-affine-formal-algebraic-space} 中某种选择
$X = \colim X_\lambda$，概形 $X_\lambda$ 都是 Noether 的；
\item $X$ 弱 adic，并且对定义
\ref{formal-spaces-definition-affine-formal-algebraic-space} 中某种选择
$X = \colim X_\lambda$，概形 $X_\lambda$ 都是 Noether 的。
\end{enumerate}
\end{lemma}

\begin{proof}
假设 $X$ 是 Noether 的。则 $X = \text{Spf}(A)$，其中 $A$ 是 Noether adic 环。
设 $T \to X$ 是闭浸入，其中 $T$ 是概形。由引理
\ref{formal-spaces-lemma-mcquillan-affine-formal-algebraic-space}，$T$ 仿射，且
$T \to \Spec(A)$ 是闭浸入。由于 $A$ Noether，可见 $T$ Noether。
由此可见 (1) $\Rightarrow$ (2)。

\medskip\noindent
蕴含 (2) $\Rightarrow$ (3) 和 (2) $\Rightarrow$ (4) 是立即的
（见引理 \ref{formal-spaces-lemma-implications-between-types}）。

\medskip\noindent
为证明 (3) $\Rightarrow$ (1)，把 $X = \text{Spf}(A)$ 写成某个
adic 环 $A$ 的形式谱，其中 $A$ 有有限生成的定义理想 $I$。
还给定：对某个开理想基本系 $I_\lambda$，环 $A/I_\lambda$ 都 Noether。
由于 $I$ 开，可以找到 $\lambda$ 使 $I_\lambda \subset I$。
则 $A/I$ Noether，由《代数》之引理
\ref{algebra-lemma-completion-Noetherian} 可知 $A$ Noether。

\medskip\noindent
为证明 (4) $\Rightarrow$ (3)，把 $X = \text{Spf}(A)$ 写成某个弱 adic 环
$A$ 的形式谱。则 $A$ 可容许，并有一个定义理想 $I$；而且 $I^2$ 的闭包
$I_2$ 是开的，见引理 \ref{formal-spaces-lemma-weakly-adic}。
还给定：对某个开理想基本系 $I_\lambda$，环 $A/I_\lambda$ 都 Noether。
选取 $\lambda$ 使 $I_\lambda \subset I_2$。则 $A/I_2$ 是
$A/I_\lambda$ 的商，所以 Noether。故 $I/I_2$ 是有限 $A$-模。
由引理 \ref{formal-spaces-lemma-weakly-adic-finite-generation}，$A$ 是有有限生成定义理想的
adic 环。因此 $X$ adic*，即 (3) 成立。
\end{proof}





\section{形式代数空间}
\label{formal-spaces-section-formal-algebraic-spaces}

\noindent
我们暂时打破先对环、再对概形、最后对代数空间引入新概念的习惯，
不先引入形式概形便直接引入形式代数空间。一般思路是：形式代数空间是 fppf 拓扑中的层，
并且 étale 局部地是 \cite{BVGD} 意义下的仿射形式概形。
相关材料可见 \cite{Yasuda}。

\medskip\noindent
在形式代数空间的定义中，我们将借用《自举》之第
\ref{bootstrap-section-morphism-representable-by-spaces} 节和第
\ref{bootstrap-section-representable-by-spaces-properties} 节中的若干术语。

\begin{definition}
\label{formal-spaces-definition-formal-algebraic-space}
设 $S$ 是概形。若存在一族层映射
$\{X_i \to X\}_{i \in I}$，使得
\begin{enumerate}
\item $X_i$ 是仿射形式代数空间；
\item $X_i \to X$ 可由代数空间表示且为 étale；
\item $\coprod X_i \to X$ 作为层映射是满射，
\end{enumerate}
并且 $X$ 满足一个集合论条件（见注 \ref{formal-spaces-remark-set-theoretic}），
则称 $(\Sch/S)_{fppf}$ 上的层 $X$ 是{\it 形式代数空间}。
$S$ 上形式代数空间的{\it 态射}是层的映射。
\end{definition}

\noindent
讨论。合理性检查：仿射形式代数空间是形式代数空间。
在定义的情形中，态射 $X_i \to X$ 可表示（由概形表示），
见引理 \ref{formal-spaces-lemma-presentation-representable}。由《自举》之引理
\ref{bootstrap-lemma-surjective-flat-locally-finite-presentation}，
也可以不要求 $\coprod X_i \to X$ 作为层映射是满射，而只要求它满
（这是有意义的，因为它可表示）。

\medskip\noindent
我们的形式代数空间概念{\bf 非常一般}。事实上，即使上面定义的仿射形式代数空间
也是非常难处理的对象。

\begin{lemma}
\label{formal-spaces-lemma-diagonal-formal-algebraic-space}
设 $S$ 是概形。若 $X$ 是 $S$ 上的形式代数空间，则对角态射
$\Delta : X \to X \times_S X$ 可表示、为单态射、局部拟有限、局部有限型且分离。
\end{lemma}

\begin{proof}
设给定 $U \to X$ 和 $V \to X$，其中 $U,V$ 是 $S$ 上的概形。
则 $U \times_X V$ 是层。选取定义 \ref{formal-spaces-definition-formal-algebraic-space}
中的 $\{X_i \to X\}$。对每个 $i$，态射
$$
(U \times_X X_i) \times_{X_i} (V \times_X X_i)
= (U \times_X V) \times_X X_i \to U \times_X V
$$
作为 $X_i \to X$ 的基变换可表示且为 étale，而其源是概形
（使用引理 \ref{formal-spaces-lemma-diagonal-affine-formal-algebraic-space} 和
\ref{formal-spaces-lemma-presentation-representable}）。这些映射联合满，所以由《自举》之定理
\ref{bootstrap-theorem-final-bootstrap}，$U \times_X V$ 是代数空间。
态射 $U \times_X V \to U \times_S V$ 是单态射。它还局部拟有限，
因为预复合上面显示的态射后，得到复合
$$
(U \times_X X_i) \times_{X_i} (V \times_X X_i)
\to (U \times_X X_i) \times_S (V \times_X X_i)
\to U \times_S V
$$
它是闭浸入（引理 \ref{formal-spaces-lemma-diagonal-affine-formal-algebraic-space}）
与 étale 态射的复合，故局部拟有限，见《空间上的下降》之引理
\ref{spaces-descent-lemma-locally-quasi-finite-etale-local-source}。
因此由《空间的态射》之命题
\ref{spaces-morphisms-proposition-locally-quasi-finite-separated-over-scheme}，
$U \times_X V$ 是概形。所以由《空间》之引理
\ref{spaces-lemma-representable-diagonal}，$\Delta$ 可表示。

\medskip\noindent
事实上，上面已经证明概形态射
$U \times_X V \to U \times_S V$ 总是单态射且局部拟有限，
所以由《空间》之引理 \ref{spaces-lemma-transformation-diagonal-properties}，
$\Delta : X \to X \times_S X$ 是单态射且局部拟有限。
然后可用《空间》之引理
\ref{spaces-lemma-representable-transformations-property-implication} 的原理，
看出 $\Delta$ 分离且局部有限型。事实上，概形的单态射分离
（《概形》之引理 \ref{schemes-lemma-monomorphism-separated}），
而概形的局部拟有限态射局部有限型
（由《态射》之第 \ref{morphisms-section-quasi-finite} 节中的定义得到）。
\end{proof}

\begin{lemma}
\label{formal-spaces-lemma-space-to-formal-space}
设 $S$ 是概形，$f : X \to Y$ 是从 $S$ 上代数空间到 $S$ 上形式代数空间的态射。
则 $f$ 可由代数空间表示。
\end{lemma}

\begin{proof}
设 $Z \to Y$ 是态射，其中 $Z$ 是 $S$ 上的概形。须证明
$X \times_Y Z$ 是代数空间。选取概形 $U$ 以及满的 étale 态射 $U \to X$。
则 $U \times_Y Z \to X \times_Y Z$ 是可表示、满且 étale 的
（《空间》之引理
\ref{spaces-lemma-base-change-representable-transformations-property}），
并且由引理 \ref{formal-spaces-lemma-diagonal-formal-algebraic-space}，
$U \times_Y Z$ 是概形。结论遂由《自举》之定理
\ref{bootstrap-theorem-final-bootstrap} 得出。
\end{proof}

\begin{remark}
\label{formal-spaces-remark-compare-with-formal-schemes}
略去集合论问题，EGA 意义下的形式概形范畴
（见第 \ref{formal-spaces-section-formal-schemes-EGA} 节）等价于形式代数空间范畴的一个全子范畴。
为解释这一点，假设基概形是 $\Spec(\mathbf{Z})$。
由引理 \ref{formal-spaces-lemma-formal-scheme-sheaf-fppf}，与形式概形
$\mathfrak X$ 相伴的点函子 $h_\mathfrak X$ 是 fppf 拓扑中的层。
由引理 \ref{formal-spaces-lemma-fully-faithful}，赋值
$\mathfrak X \mapsto h_\mathfrak X$ 是从形式概形范畴到 fppf 层范畴的全忠实嵌入。
给定形式概形 $\mathfrak X$，选取由仿射形式概形 $\mathfrak X_i$ 组成的开覆盖
$\mathfrak X = \bigcup \mathfrak X_i$。由注
\ref{formal-spaces-remark-compare-with-affine-formal-schemes}，$h_{\mathfrak X_i}$
是仿射形式代数空间。态射
$h_{\mathfrak X_i} \to h_\mathfrak X$ 可表示且为开浸入。
所以 $\{h_{\mathfrak X_i} \to h_\mathfrak X\}$ 是定义
\ref{formal-spaces-definition-formal-algebraic-space} 中的一族映射，
可见 $h_\mathfrak X$ 是形式代数空间。
\end{remark}

\begin{remark}
\label{formal-spaces-remark-set-theoretic}
设 $S$ 是概形，并设 $(\Sch/S)_{fppf}$ 是《拓扑》之定义
\ref{topologies-definition-big-small-fppf} 中的大 fppf 站点。
在定义 \ref{formal-spaces-definition-affine-formal-algebraic-space} 和
\ref{formal-spaces-definition-formal-algebraic-space} 中，采用如下 $X$ 的集合论条件：
存在 $(\Sch/S)_{fppf}$ 的对象 $U,R$、作为 fppf 层满射的态射 $U \to X$，
以及作为 fppf 层满射的态射 $R \to U \times_X U$。
换言之，要求我们的层是两个可表示层之间两映射的余等化子。
下面是说明这一概念表现合理的若干观察：
\begin{enumerate}
\item 设 $X = \colim_{\lambda \in \Lambda} X_\lambda$，
并且该系统满足定义 \ref{formal-spaces-definition-affine-formal-algebraic-space} 的条件 (1) 和 (2)。
则 $U = \coprod_{\lambda \in \Lambda} X_\lambda \to X$ 是 fppf 层的满射。
此外，由引理 \ref{formal-spaces-lemma-diagonal-affine-formal-algebraic-space}，
$U \times_X U$ 是 $U \times_S U$ 的闭子概形。
所以，若 $U$ 可由 $(\Sch/S)_{fppf}$ 的对象表示，则 $U \times_S U$ 也可
（见《集合》之引理 \ref{sets-lemma-what-is-in-it}），集合论条件得到满足。
当 $\Lambda$ 可数时总是如此，见《集合》之引理
\ref{sets-lemma-what-is-in-it}。
\item 合理性检查。设 $\{X_i \to X\}_{i \in I}$ 如定义
\ref{formal-spaces-definition-formal-algebraic-space}（集合论条件按上文表述），
并假设每个 $X_i$ 实际上都是仿射概形。则 $X$ 是代数空间。
事实上，若选取更大的大 fppf 站点 $(\Sch'/S)_{fppf}$，使
$U' = \coprod X_i$ 和 $R' = \coprod X_i \times_X X_j$
都可由其中的对象表示，则 $X' = U'/R'$ 是这一选择下代数空间范畴的对象。
然后应用《空间》之引理 \ref{spaces-lemma-fully-faithful}，可知
$X$ 是 $(\Sch/S)_{fppf}$ 上的代数空间。
\item 设 $\{X_i \to X\}_{i \in I}$ 是满足定义
\ref{formal-spaces-definition-formal-algebraic-space} 中条件 (1)、(2)、(3) 的一族层映射。
对每个 $i$，可以选取 $U_i \in \Ob((\Sch/S)_{fppf})$ 以及层的满射
$U_i \to X_i$。所以，若 $I$ 不过大（例如可数），则
$U = \coprod U_i \to X$ 是层的满射，并且 $U$ 可由
$(\Sch/S)_{fppf}$ 的对象表示。要得到满射到 $U \times_X U$ 的
$R \in \Ob((\Sch/S)_{fppf})$，只须假设对角
$\Delta : X \to X \times_S X$ 不太难处理；例如，若 $X$ 的对角拟紧，
即 $X$ 拟分离，就总能做到。
\end{enumerate}
\end{remark}







\section{约化}
\label{formal-spaces-section-reduction}

\noindent
如下引理所示，每个形式代数空间都有一个底层既约代数空间。

\begin{lemma}
\label{formal-spaces-lemma-reduction-formal-algebraic-space}
设 $S$ 是概形，$X$ 是 $S$ 上的形式代数空间。
存在既约代数空间 $X_{red}$ 以及可表示且为增厚的态射
$X_{red} \to X$。若 $U$ 是既约代数空间，则态射 $U \to X$
唯一地分解经过 $X_{red}$。
\end{lemma}

\begin{proof}
首先假设 $X$ 是仿射形式代数空间。按定义
\ref{formal-spaces-definition-affine-formal-algebraic-space} 写
$X = \colim X_\lambda$。由于转移态射都是增厚，仿射概形
$X_\lambda$ 都有彼此同构的约化 $X_{red}$。
由引理 \ref{formal-spaces-lemma-covering-by-thickenings} 以及增厚的复合仍为增厚这一事实，
态射 $X_{red} \to X$ 可表示且为增厚。我们略去泛性质的核验
（使用《概形》之定义 \ref{schemes-definition-reduced-induced-scheme}、
《概形》之引理 \ref{schemes-lemma-map-into-reduction}、
《空间的性质》之定义 \ref{spaces-properties-definition-reduced-induced-space}
以及《空间的性质》之引理
\ref{spaces-properties-lemma-map-into-reduction}）。

\medskip\noindent
设 $X$ 和 $\{X_i \to X\}_{i \in I}$ 如定义
\ref{formal-spaces-definition-formal-algebraic-space}。对每个 $i$，令
$X_{i, red} \to X_i$ 为上面构造的约化。对 $i,j \in I$，投影
$X_{i, red} \times_X X_j \to X_{i, red}$ 是概形的 étale 态射
（étale 由假设，源是概形由引理 \ref{formal-spaces-lemma-presentation-representable}）。
所以 $X_{i, red} \times_X X_j$ 既约
（见《下降》之引理 \ref{descent-lemma-reduced-local-smooth}）。
因此，由泛性质，投影
$X_{i, red} \times_X X_j \to X_j$ 分解经过 $X_{j, red}$。
由于态射 $X_{i, red} \to X_i$ 是层的单射，可得
$$
R_{ij} = X_{i, red} \times_X X_j = X_{i, red} \times_X X_{j, red} =
X_i \times_X X_{j, red}
$$
令 $U = \coprod X_{i, red}$、$R = \coprod R_{ij}$，并以
$s,t : R \to U$ 表示两个投影。作为层，$R = U \times_X U$，
并且 $s$ 和 $t$ 都是 étale 的。于是
% P10-E0307
$(t,s) : R \to U$ 由上述观察定义一个 étale 等价关系。
所以由《空间》之定理 \ref{spaces-theorem-presentation}，
$X_{red} = U/R$ 是代数空间。按构造，图表
$$
\xymatrix{
\coprod X_{i, red} \ar[r] \ar[d] & \coprod X_i \ar[d] \\
X_{red} \ar[r] & X
}
$$
是笛卡尔的。由于右侧竖直箭头 étale 且满，顶部水平箭头可表示且为增厚，
由《自举》之引理 \ref{bootstrap-lemma-after-fppf-sep-lqf}，
$X_{red} \to X$ 可表示
（核验该引理的假设时，使用满的 étale 态射是满、平坦且局部有限表示的，
并使用增厚分离且局部拟有限）。然后可以用《空间》之引理
\ref{spaces-lemma-descent-representable-transformations-property}
推出 $X_{red} \to X$ 是增厚
（使用增厚等价于满的闭浸入这一事实）。

\medskip\noindent
最后，设 $U \to X$ 是态射，其中 $U$ 是 $S$ 上的既约代数空间。
则每个 $X_i \times_X U$ 都在 $U$ 上 étale，因而既约
（按《空间的性质》之第
\ref{spaces-properties-section-types-properties} 节中既约代数空间的定义）。
所以 $X_i \times_X U \to X_i$ 分解经过 $X_{i,red}$。
由于 $\{X_i \times_X U \to U\}$ 是 étale 覆盖，$U \to X$
分解经过 $X_{red}$。
\end{proof}

\begin{example}
\label{formal-spaces-example-reduction-affine-formal-spectrum}
设 $A$ 是弱可容许拓扑环。此时
$$
\text{Spf}(A)_{red} = \Spec(A/\mathfrak a)
$$
其中 $\mathfrak a \subset A$ 是拓扑幂零元素的理想。
事实上，$\mathfrak a$ 是根式理想
（引理 \ref{formal-spaces-lemma-topologically-nilpotent}），并且因为 $A$ 弱可容许，它是开的。
\end{example}

\begin{lemma}
\label{formal-spaces-lemma-reduction-smooth}
设 $S$ 是概形，$f : X \to Y$ 是 $S$ 上形式代数空间的态射，
它可由代数空间表示且光滑（例如 étale）。则
$X_{red} = X \times_Y Y_{red}$。
\end{lemma}

\begin{proof}
（étale 情形直接来自引理 \ref{formal-spaces-lemma-reduction-formal-algebraic-space}
证明中对底层既约代数空间的构造。）假设 $f$ 光滑。
注意，$X \times_Y Y_{red} \to Y_{red}$ 是代数空间的光滑态射。
所以由《空间上的下降》之引理
\ref{spaces-descent-lemma-reduced-local-smooth}，
$X \times_Y Y_{red}$ 是既约代数空间。于是约化的泛性质表明典范态射
$X_{red} \to X \times_Y Y_{red}$ 是同构。
\end{proof}

\begin{lemma}
\label{formal-spaces-lemma-reduction-surjective}
设 $S$ 是概形，$f : X \to Y$ 是 $S$ 上形式代数空间的态射，
它可由代数空间表示。则 $f$ 在《自举》之定义
\ref{bootstrap-definition-property-transformation} 的意义下为满射，
当且仅当 $f_{red} : X_{red} \to Y_{red}$ 是代数空间的满态射。
\end{lemma}

\begin{proof}
略。
\end{proof}






\section{沿增厚的代数空间余极限}
\label{formal-spaces-section-global-colimits}

\noindent
一种特殊的形式代数空间可以像下面引理那样，在整体上写成沿增厚的代数空间余滤余极限。
稍后（在第 \ref{formal-spaces-section-quasi-compact-quasi-separated} 节）将看到，
任意拟紧且拟分离的形式代数空间都是这样的整体余极限。

\begin{lemma}
\label{formal-spaces-lemma-colimit-is-formal}
设 $S$ 是概形。给定有向集 $\Lambda$ 以及位于 $\Lambda$ 上的代数空间系统
$(X_\lambda, f_{\lambda \mu})$，其中每个
$f_{\lambda \mu} : X_\lambda \to X_\mu$ 都是增厚。
则 $X = \colim_{\lambda \in \Lambda} X_\lambda$ 是 $S$ 上的形式代数空间。
\end{lemma}

\begin{proof}
因为在 fppf 层范畴中取余极限，$X$ 是层。选定 $\lambda \in \Lambda$。
选取 étale 覆盖 $\{X_{i,\lambda} \to X_\lambda\}$，其中
% P10-E0308
$X_i$ 是 $S$ 上的仿射概形，见《空间的性质》之引理
\ref{spaces-properties-lemma-cover-by-union-affines}。
对每个 $\mu \geq \lambda$，存在笛卡尔图表
$$
\xymatrix{
X_{i, \lambda} \ar[r] \ar[d] & X_{i, \mu} \ar[d] \\
X_\lambda \ar[r] & X_\mu
}
$$
其竖直箭头为 étale，见《空间态射续篇》之定理
\ref{spaces-more-morphisms-theorem-topological-invariance}
（这里还用到增厚是满足该定理条件的满闭浸入）。此外，这些图表在唯一同构意义下唯一，
所以当 $\mu' \geq \mu$ 时
$X_{i,\mu} = X_\mu \times_{X_{\mu'}} X_{i,\mu'}$。
% P10-E0309
态射 $X_{i,\mu} \to X_{i,\mu'}$ 是增厚，因为它是增厚的基变换。
由《空间的极限》之命题 \ref{spaces-limits-proposition-affine}
以及 $X_{i,\lambda}$ 仿射这一事实，每个 $X_{i,\mu}$ 都是仿射概形。
令 $X_i = \colim_{\mu \geq \lambda} X_{i,\mu}$。
则 $X_i$ 是仿射形式代数空间。态射 $X_i \to X$ 是 étale 的，
因为给定仿射概形 $U$，任意 $U \to X$ 都分解经过某个
$\mu \geq \lambda$ 对应的 $X_\mu$（细节略）。由此可见 $X$ 是形式代数空间。
\end{proof}

\noindent
设 $S$ 是概形，$X$ 是 $S$ 上的形式代数空间。如何证明或检验
$X$ 是引理 \ref{formal-spaces-lemma-colimit-is-formal} 中那样的整体余极限？
为此，寻找映射 $i : Z \to X$，其中 $Z$ 是 $S$ 上的代数空间，而
$i$ 是满的闭浸入；换言之，$i$ 是增厚。这是有意义的，因为
$i$ 可由代数空间表示（引理 \ref{formal-spaces-lemma-space-to-formal-space}），
可以像前面一样使用《自举》之定义
\ref{bootstrap-definition-property-transformation}。

\begin{example}
\label{formal-spaces-example-david-hansen}
设 $(A, \mathfrak m, \kappa)$ 是赋值环，对某个非零
$\pi \in \mathfrak m$，它关于 $(\pi)$-进拓扑完备。
再假设 $\mathfrak m$ 不是有限生成的。一个例子是
$A = \mathcal{O}_{\mathbf{C}_p}$、$\pi = p$，其中
$\mathcal{O}_{\mathbf{C}_p}$ 是 $p$-进复数域 $\mathbf{C}_p$
的整数环（$\mathbf{C}_p$ 是 $\mathbf{Q}_p$ 的代数闭包的完备化）。
另一个例子是
$$
A =
\left\{
\sum\nolimits_{\alpha \in \mathbf{Q},\ \alpha \geq 0} a_\alpha t^\alpha
\middle|
\begin{matrix}
a_\alpha \in \kappa\text{，并且对所有 }n\text{，只有有限多个} \\
\alpha \leq n\text{ 对应的 }a_\alpha\text{ 非零}
\end{matrix}
\right\}
$$
以及 $\pi = t$。则 $X = \text{Spf}(A)$ 是仿射形式代数空间，
$\Spec(\kappa) \to X$ 是增厚，对应于弱定义理想
$\mathfrak m \subset A$，但后者不是定义理想。
\end{example}

\begin{remark}[弱定义理想]
\label{formal-spaces-remark-weak-ideals-of-definition}
设 $\mathfrak X$ 是 McQuillan 意义下的形式概形，见注
\ref{formal-spaces-remark-mcquillan}。$\mathfrak X$ 的一个{\it 弱定义理想}
是理想层 $\mathcal{I} \subset \mathcal{O}_\mathfrak X$，使得对每个仿射形式开子概形
$\mathfrak U \subset \mathfrak X$，理想
$\mathcal{I}(\mathfrak U) \subset \mathcal{O}_\mathfrak X(\mathfrak U)$
都是弱可容许拓扑环 $\mathcal{O}_\mathfrak X(\mathfrak U)$ 的弱定义理想。
只须在某个仿射开覆盖的成员上检验该条件。有一一对应
$$
\{\mathfrak X\text{ 的弱定义理想}\}
\leftrightarrow
\{\text{如上所述的增厚 }i : Z \to h_\mathfrak X\}
$$
该对应把 $\mathcal{I}$ 映到概形
$Z = (\mathfrak X, \mathcal{O}_\mathfrak X/\mathcal{I})$
及其到 $\mathfrak X$ 的显然态射。{\it 弱定义理想基本系}是弱定义理想族
$\mathcal{I}_\lambda$，使得对每个仿射形式开子概形
$\mathfrak U \subset \mathfrak X$，理想
$$
I_\lambda = \mathcal{I}_\lambda(\mathfrak U) \subset
A = \Gamma(\mathfrak U, \mathcal{O}_\mathfrak X)
$$
组成弱可容许拓扑环 $A$ 的一族弱定义理想基本系。
只须在某个仿射开覆盖的成员上检验。由此得到：与 McQuillan 形式概形
$\mathfrak X$ 相伴的形式代数空间 $h_\mathfrak X$ 是引理
\ref{formal-spaces-lemma-colimit-is-formal} 中那样的概形余极限，当且仅当
$\mathfrak X$ 存在一族弱定义理想基本系。
\end{remark}

\begin{remark}[定义理想]
\label{formal-spaces-remark-ideals-of-definition}
设 $\mathfrak X$ 是 EGA 意义下的形式概形。$\mathfrak X$ 的一个{\it 定义理想}
是理想层 $\mathcal{I} \subset \mathcal{O}_\mathfrak X$，使得对每个仿射形式开子概形
$\mathfrak U \subset \mathfrak X$，理想
$\mathcal{I}(\mathfrak U) \subset \mathcal{O}_\mathfrak X(\mathfrak U)$
都是可容许拓扑环 $\mathcal{O}_\mathfrak X(\mathfrak U)$ 的定义理想。
只须在某个仿射开覆盖的成员上检验该条件。这里{\bf 并不}像注
\ref{formal-spaces-remark-weak-ideals-of-definition} 那样得到定义理想与增厚
$Z \to h_\mathfrak X$ 之间的同一对应；例
\ref{formal-spaces-example-david-hansen} 给出了一个例子。
{\it 定义理想基本系}是定义理想族 $\mathcal{I}_\lambda$，
使得对每个仿射形式开子概形 $\mathfrak U \subset \mathfrak X$，理想
$$
I_\lambda = \mathcal{I}_\lambda(\mathfrak U) \subset
A = \Gamma(\mathfrak U, \mathcal{O}_\mathfrak X)
$$
组成可容许拓扑环 $A$ 的一族定义理想基本系。
只须在某个仿射开覆盖的成员上检验。假设 $\mathfrak X$ 拟紧，
而 $\{\mathcal{I}_\lambda\}_{\lambda \in \Lambda}$ 是一族弱定义理想基本系。
若 $A$ 是可容许拓扑环，则所有充分小的开理想都是定义理想
（即包含在某个定义理想中的任意开理想都是定义理想）。
所以，因为只需在一个仿射开覆盖的有限多个成员上检验，
对充分大的 $\lambda$，$\mathcal{I}_\lambda$ 是定义理想。
利用注 \ref{formal-spaces-remark-weak-ideals-of-definition} 中的讨论，得到：
与 EGA 意义下拟紧形式概形 $\mathfrak X$ 相伴的形式代数空间
$h_\mathfrak X$ 是引理 \ref{formal-spaces-lemma-colimit-is-formal} 中那样的概形余极限，
当且仅当 $\mathfrak X$ 存在一族定义理想基本系。
\end{remark}






\section{沿闭子集的完备化}
\label{formal-spaces-section-completion}

\noindent
我们的形式代数空间概念很适合沿闭子集取完备化。

\begin{lemma}
\label{formal-spaces-lemma-completion-affine-is-affine-formal-algebraic-space}
设 $S$ 是概形，$X$ 是 $S$ 上的仿射概形，$T \subset |X|$ 是闭子集。
则函子
$$
(\Sch/S)_{fppf} \longrightarrow \textit{Sets},\quad
U \longmapsto \{f : U \to X \mid f(|U|) \subset T\}
$$
是 McQuillan 仿射形式代数空间。
\end{lemma}

\begin{proof}
设 $X = \Spec(A)$，而 $T$ 对应于根式理想 $I \subset A$。
设 $U = \Spec(B)$ 是 $S$ 上的仿射概形，$f : U \to X$ 是 $F(U)$
的元素。则 $f$ 对应于环映射 $\varphi : A \to B$，使得 $B$
的每个素理想都包含 $\varphi(I)B$。所以 $\varphi(I)$ 的每个元素在 $B$
中幂零，见《代数》之引理 \ref{algebra-lemma-Zariski-topology}。
% P10-E0310
令 $J = \Ker(\varphi)$，可知 $I/J$ 是 $A/J$ 中的局部幂零理想。
等价地，$V(J) = V(I) = T$。换言之，引理中的函子等于
$\colim \Spec(A/J)$，其中余极限遍历满足 $V(J) = T$ 的理想 $J$。
因此该函子是仿射形式代数空间。它是 McQuillan 的
（定义 \ref{formal-spaces-definition-types-affine-formal-algebraic-space}），
因为映射 $A \to A/J$ 是满射，从而
$A^\wedge = \lim A/J \to A/J$ 是满射，见引理
\ref{formal-spaces-lemma-mcquillan-affine-formal-algebraic-space}。
\end{proof}

\begin{lemma}
\label{formal-spaces-lemma-completion-is-formal-algebraic-space}
设 $S$ 是概形，$X$ 是 $S$ 上的代数空间，$T \subset |X|$ 是闭子集。
则函子
$$
(\Sch/S)_{fppf} \longrightarrow \textit{Sets},\quad
U \longmapsto \{f : U \to X \mid f(|U|) \subset T\}
$$
是形式代数空间。
\end{lemma}

\begin{proof}
以 $F$ 表示该函子。设 $\{U_i \to U\}$ 是 fppf 覆盖。
则 $\coprod |U_i| \to |U|$ 是满射。由于 $X$ 是 fppf 层，
可知 $F$ 是 fppf 层。

\medskip\noindent
选取 étale 覆盖 $\{g_i : X_i \to X\}$，使每个 $X_i$ 都仿射，且对所有 $i$，
见《空间的性质》之引理
\ref{spaces-properties-lemma-cover-by-union-affines}。
态射 $F \times_X X_i \to F$ 是 étale 的
（见《空间》之引理
\ref{spaces-lemma-base-change-representable-transformations-property}），
而映射 $\coprod F \times_X X_i \to F$ 是层的满射。
因此，只须证明 $F \times_X X_i$ 是仿射形式代数空间。
$F \times_X X_i$ 的一个 $U$-值点是态射 $U \to X_i$，
其像包含在闭子集 $g_i^{-1}(T) \subset |X_i|$ 中。
所以结论由引理
\ref{formal-spaces-lemma-completion-affine-is-affine-formal-algebraic-space} 得出。
\end{proof}

\begin{definition}
\label{formal-spaces-definition-completion}
设 $S$ 是概形，$X$ 是 $S$ 上的代数空间，$T \subset |X|$ 是闭子集。
引理 \ref{formal-spaces-lemma-completion-is-formal-algebraic-space} 中的形式代数空间
称为 $X$ {\it 沿 $T$ 的完备化}。
\end{definition}

\noindent
\cite[Chapter I, Section 10.8]{EGA} 用记号 $X_{/T}$ 表示该完备化，
我们有时也使用这一记号。设 $f : X \to X'$ 是概形 $S$ 上代数空间的态射。
假设 $T \subset |X|$ 和 $T' \subset |X'|$ 是闭子集，且
$|f|(T) \subset T'$。则显然，$f$ 定义完备化之间的形式代数空间态射
$$
X_{/T} \longrightarrow X'_{/T'}
$$


\begin{lemma}
\label{formal-spaces-lemma-map-completions-representable}
设 $S$ 是概形，$f : X' \to X$ 是 $S$ 上代数空间的态射。
设 $T \subset |X|$ 是闭子集，并令
$T' = |f|^{-1}(T) \subset |X'|$。则
$$
\xymatrix{
X'_{/T'} \ar[r] \ar[d] & X' \ar[d]^f \\
X_{/T} \ar[r] & X
}
$$
是层的笛卡尔图表。特别地，态射
$X'_{/T'} \to X_{/T}$ 可由代数空间表示。
\end{lemma}

\begin{proof}
事实上，设 $Y \to X$ 是从概形到 $X$ 的态射，使 $|Y|$ 映入 $T$。
则
% P10-E0311
$Y \times_X X' \to X$ 是代数空间的态射，使
$|Y \times_X X'|$ 映入 $T'$。所以函子
$Y \times_{X_{/T}} X'_{/T'}$ 由 $Y \times_X X'$ 表示，
从而引理成立。
\end{proof}

\begin{lemma}
\label{formal-spaces-lemma-reduction-completion}
设 $S$ 是概形，$X$ 是 $S$ 上的代数空间，$T \subset |X|$ 是闭子集。
$X$ 沿 $T$ 的完备化 $X_{/T}$ 的约化 $(X_{/T})_{red}$，
是 $X$ 中与 $T$ 对应的既约诱导闭子空间 $Z$。
\end{lemma}

\begin{proof}
由引理 \ref{formal-spaces-lemma-reduction-formal-algebraic-space}、《空间的性质》之定义
\ref{spaces-properties-definition-reduced-induced-space}
（它使用《空间的性质》之引理
\ref{spaces-properties-lemma-reduced-closed-subspace} 构造 $Z$）
以及 $X_{/T}$ 的定义可知，$Z$ 和 $(X_{/T})_{red}$ 是由同一映射性质
% P10-E0312
刻画既约代数空间：源为既约代数空间的态射 $g : Y \to X$
分解经过它们，当且仅当 $|Y|$ 映入 $T \subset |X|$。
\end{proof}

\begin{lemma}
\label{formal-spaces-lemma-affine-formal-completion-types}
设 $S$ 是概形，$X = \Spec(A)$ 是 $S$ 上的仿射概形，
$T \subset X$ 是闭子集，$X_{/T}$ 是 $X$ 沿 $T$ 的形式完备化。
\begin{enumerate}
\item 若 $X \setminus T$ 拟紧，即 $T$ 可构造，则 $X_{/T}$ adic*。
\item 若对某个有限生成理想 $I \subset A$ 有 $T = V(I)$，
则 $X_{/T} = \text{Spf}(A^\wedge)$，其中 $A^\wedge$ 是 $A$
的 $I$-进完备化。
\item 若 $X$ Noether，则 $X_{/T}$ Noether。
\end{enumerate}
\end{lemma}

\begin{proof}
由《代数》之引理 \ref{algebra-lemma-qc-open}，若 (1) 成立，
则可以找到 (2) 中的理想 $I \subset A$。若 (3) 成立，
也可以找到 (2) 中的理想 $I \subset A$。此外，由《代数》之引理
\ref{algebra-lemma-completion-Noetherian-Noetherian}，Noether 环的完备化仍 Noether。
综上，只须证明 (2)。

\medskip\noindent
(2) 的证明。设 $I = (f_1, \ldots, f_r) \subset A$ 切出 $T$。
若 $Z = \Spec(B)$ 是仿射概形，$g : Z \to X$ 是满足
$g(Z) \subset T$（集合论意义下）的态射，则对每个 $i$，
$g^\sharp(f_i)$ 在 $B$ 中幂零。所以对某个 $n$，$I^n$ 在 $B$ 中映为零。
因此可见
$X_{/T} = \colim \Spec(A/I^n) = \text{Spf}(A^\wedge)$。
\end{proof}

\noindent
下列引理由 Ofer Gabber 给出。

\begin{lemma}
\label{formal-spaces-lemma-completion-countably-indexed}
\begin{reference}
Ofer Gabber 于 2014 年 9 月 11 日的电子邮件。
\end{reference}
设 $S$ 是概形，$X = \Spec(A)$ 是 $S$ 上的仿射概形，
$T \subset X$ 是闭子概形。
\begin{enumerate}
\item 若形式完备化 $X_{/T}$ 是可数指标的，并且存在可数多个
$f_1,f_2,f_3,\ldots \in A$ 使
$T = V(f_1,f_2,f_3,\ldots)$，则 $X_{/T}$ adic*。
\item 若去掉 $T$ 可由 $X$ 中可数多个函数切出这一假设，(1) 的结论不成立。
\end{enumerate}
\end{lemma}

\begin{proof}
$X_{/T}$ 是可数指标的这一假设表示，存在理想序列
$$
A \supset J_1 \supset J_2 \supset J_3 \supset \ldots
$$
满足 $V(J_n) = T$，并且
% P10-E0313
每个满足 $V(J) = T$ 的理想 $J \subset A$ 都有某个 $n$ 使 $J \supset J_n$。

\medskip\noindent
为构造 (2) 的例子，令 $\omega_1$ 为第一个不可数序数。
设 $k$ 是域，$A$ 是由 $x_\alpha$（$\alpha \in \omega_1$）和
$y_{\alpha\beta}$（$\alpha \in \beta \in \omega_1$）生成的 $k$-代数，
满足关系 $x_\alpha = y_{\alpha\beta}x_\beta$。
令 $T = V(x_\alpha)$，$J_n = (x_\alpha^n)$。
若理想 $J \subset A$ 满足 $V(J) = T$，则对某个
$n_\alpha \geq 1$ 有 $x_\alpha^{n_\alpha} \in J$。
集合 $\{\alpha \mid n_\alpha = n\}$ 中至少有一个必须在 $\omega_1$ 中无界。
于是这些关系蕴含 $J_n \subset J$。

\medskip\noindent
为看出 (2) 成立，现在只须证明
$A^\wedge = \lim A/J_n$ 不是关于某个有限生成理想完备的环。
对 $\gamma \in \omega_1$，令 $A_\gamma$ 为 $A$ 关于如下理想的商：
该理想由 $x_\alpha$（$\alpha \in \gamma$）以及
$y_{\alpha\beta}$（$\alpha \in \gamma$）生成。
因为 $A/J_1$ 既约，$\lim A/J_n$ 的每个拓扑幂零元素 $f$ 都属于
$J_1^\wedge = \lim J_1/J_n$。这表示 $f$ 是只涉及可数多个生成元的无穷级数。
所以对某个 $\gamma$，$f$ 在
$A_\gamma^\wedge = \lim A_\gamma/J_nA_\gamma$ 中消失。
注意，由引理 \ref{formal-spaces-lemma-ses}，
$A^\wedge \to A_\gamma^\wedge$ 连续且开。
% P10-E0314
若 $A^\wedge$ 上的拓扑是某个有限生成理想
$I \subset A^\wedge$ 的 $I$-进拓扑，则 $I$ 会在某个
$A_\gamma^\wedge$ 中映为零。这将意味着 $A_\gamma^\wedge$ 离散，
但事实并非如此，因为有一个满的连续开映射（开性由引理
\ref{formal-spaces-lemma-ses}）
$A_\gamma^\wedge \to k[[t]]$，它对满足
$\gamma = \alpha$ 或 $\gamma \in \alpha$ 的指标由
$x_\alpha \mapsto t$、$y_{\alpha\beta} \mapsto 1$ 给出。

\medskip\noindent
证明 (1) 之前，先证明如下事实：若 $I \subset A^\wedge$ 是有限生成理想，
其闭包 $\bar I$ 是开的，则 $I = \bar I$。
由于 $V(J_n^2) = T$，存在 $m$ 使 $J_n^2 \supset J_m$。
所以，取子序列后可以假设对所有 $n$ 都有
$J_n^2 \supset J_{n+1}$。令
$J_n^\wedge = \lim_{k \geq n} J_n/J_k \subset A^\wedge$。
由于闭包 $\bar I = \bigcap (I + J_n^\wedge)$
（引理 \ref{formal-spaces-lemma-closed}）是开的，可见存在 $m$，使得对所有
$n \geq m$ 有 $I + J_n^\wedge \supset J_m^\wedge$。固定这样的 $m$。
对所有 $n \geq m+1$，有
$$
J_{n - 1}^\wedge I + J_{n + 1}^\wedge \supset
J_{n - 1}^\wedge (I + J_{n + 1}^\wedge) \supset
J_{n - 1}^\wedge J_m^\wedge
$$
事实上，第一个包含显然，第二个已在上面证明。
因为 $J_{n-1}J_m \supset J_{n-1}^2 \supset J_n$，
这些包含表明 $J_n$ 在 $A^\wedge$ 中的像包含在理想
$J_{n-1}^\wedge I + J_{n+1}^\wedge$ 中。因为该理想是开的，得到
$$
J_{n - 1}^\wedge I + J_{n + 1}^\wedge \supset J_n^\wedge.
$$
设 $I = (g_1,\ldots,g_t)$。取 $f \in J_{m+1}^\wedge$。
利用最后显示的包含（它对所有 $n \geq m+1$ 都成立），
可以对 $c \geq 0$ 归纳写成
$$
f = \sum f_{i, c} g_i \mod J_{m + 1+ c}^\wedge
$$
其中 $f_{i,c} \in J_m^\wedge$，且
$f_{i,c} \equiv f_{i,c-1} \bmod J_{m+c}^\wedge$。
由此 $IJ_m^\wedge \supset J_{m+1}^\wedge$。
再结合 $I + J_{m+1}^\wedge \supset J_m^\wedge$，得到 $I$ 开。

\medskip\noindent
(1) 的证明。假设 $T = V(f_1,f_2,f_3,\ldots)$。
令 $I_m \subset A^\wedge$ 为由 $f_1,\ldots,f_m$ 生成的理想。
分两种情形。

\medskip\noindent
情形 I：对某个 $m$，$I_m$ 的闭包是开的。
由前一段的结果，$I_m$ 开。按构造，对任意 $n$ 有
$(J_n)^2 \supset J_{n+1}$，所以 $(J_n^\wedge)^2$ 的闭包包含
$J_{n+1}^\wedge$，因而是开的。取充分大的 $n$，可知
$A^\wedge$ 中任意两个开理想的乘积的闭包是开的。
用关于 $k$ 的归纳证明，对 $k \ge 1$，$I_m^k$ 开。
$k=1$ 的情形是情形 I 中对 $m$ 的假设。若 $k>1$，并假设
$I_m^{k-1}$ 开，则 $I_m^k = I_m^{k-1}\cdot I_m$ 是两个开理想的乘积，
因而其闭包是开的。但 $I_m^k$ 有限生成，所以由前一段
（应用于 $I=I_m^k$），$I_m^k$ 是开的，可以继续关于 $k$ 的归纳。
由于 $I_m$ 的每个元素都拓扑幂零，可知 $I_m$ 是定义理想。
这证明 $A^\wedge$ 是有有限生成定义理想的 adic 环；即 $X_{/T}$ adic*。

\medskip\noindent
情形 II：对所有 $m$，$I_m$ 的闭包 $\bar I_m$ 都不开。
则 $A^\wedge/\bar I_m$ 上的拓扑不离散。这表示可以选取
$\phi(m) \geq m$，使
$$
\Im(J_{\phi(m)} \to A/(f_1, \ldots, f_m)) \not =
\Im(J_{\phi(m) + 1} \to A/(f_1, \ldots, f_m))
$$
这里使用了
$A^\wedge/(\bar I_m + J_n^\wedge)
= A/((f_1,\ldots,f_m)+J_n)$。
选取指数 $e_i>0$，使得当 $0<m<i$ 时
$f_i^{e_i} \in J_{\phi(m)+1}$。令
$J = (f_1^{e_1}, f_2^{e_2}, f_3^{e_3},\ldots)$。
则 $V(J)=T$。我们断言对所有 $n$ 都有 $J \not\supset J_n$，
这给出矛盾，证明情形 II 不会发生。事实上，$J$ 在
$A/(f_1,\ldots,f_m)$ 中的像包含在 $J_{\phi(m)+1}$ 的像中，
而后者真包含于 $J_m$ 的像。
\end{proof}















\section{纤维积}
\label{formal-spaces-section-fibre-products}

\noindent
这是关于形式代数空间纤维积的必备一节。

\begin{lemma}
\label{formal-spaces-lemma-etale-covering-by-formal-algebraic-spaces}
设 $S$ 是概形，$\{X_i \to X\}_{i \in I}$ 是
$(\Sch/S)_{fppf}$ 上的一族层映射。假设：
(a) $X_i$ 是 $S$ 上的形式代数空间；
(b) $X_i \to X$ 可由代数空间表示且为 étale；
(c) $\coprod X_i \to X$ 是层的满射。
则 $X$ 是 $S$ 上的形式代数空间。
\end{lemma}

\begin{proof}
对每个 $i$，选取定义 \ref{formal-spaces-definition-formal-algebraic-space}
中的 $\{X_{ij} \to X_i\}_{j \in J_i}$。则
$\{X_{ij} \to X\}_{i \in I, j \in J_i}$ 是定义
\ref{formal-spaces-definition-formal-algebraic-space} 中 $X$ 所需的一族映射。
\end{proof}

\begin{lemma}
\label{formal-spaces-lemma-fibre-products-general}
设 $S$ 是概形，$X,Y$ 是 $S$ 上的形式代数空间，$Z$ 是对角可由代数空间表示的层。
设 $X \to Z$ 和 $Y \to Z$ 是层映射。则 $X \times_Z Y$ 是形式代数空间。
\end{lemma}

\begin{proof}
选取定义 \ref{formal-spaces-definition-formal-algebraic-space} 中的
$\{X_i \to X\}$ 和 $\{Y_j \to Y\}$。则
$\{X_i \times_Z Y_j \to X \times_Z Y\}$ 是一族可由代数空间表示且为 étale 的映射。
所以引理 \ref{formal-spaces-lemma-etale-covering-by-formal-algebraic-spaces} 告诉我们，
只须在 $X$ 和 $Y$ 是仿射形式代数空间时证明
$X \times_Z Y$ 是形式代数空间。

\medskip\noindent
假设 $X$ 和 $Y$ 是仿射形式代数空间。按定义
\ref{formal-spaces-definition-affine-formal-algebraic-space} 写
$X = \colim X_\lambda$ 和 $Y = \colim Y_\mu$。则
$X \times_Z Y = \colim X_\lambda \times_Z Y_\mu$。
每个 $X_\lambda \times_Z Y_\mu$ 都是代数空间。
对 $\lambda \leq \lambda'$ 和 $\mu \leq \mu'$，态射
$$
X_\lambda \times_Z Y_\mu \to
X_\lambda \times_Z Y_{\mu'} \to
X_{\lambda'} \times_Z Y_{\mu'}
$$
是增厚基变换的复合，因而是增厚。因此应用引理
\ref{formal-spaces-lemma-colimit-is-formal} 即得结论。
\end{proof}

\begin{lemma}
\label{formal-spaces-lemma-fibre-products}
设 $S$ 是概形。$S$ 上形式代数空间的范畴有纤维积。
\end{lemma}

\begin{proof}
这是引理 \ref{formal-spaces-lemma-fibre-products-general} 的特殊情形，
因为形式代数空间的对角可表示，见引理
\ref{formal-spaces-lemma-diagonal-formal-algebraic-space}。
\end{proof}

\begin{lemma}
\label{formal-spaces-lemma-reduction-fibre-products}
设 $S$ 是概形，$X \to Z$ 和 $Y \to Z$ 是 $S$ 上形式代数空间的态射。则
$(X \times_Z Y)_{red} = (X_{red} \times_{Z_{red}} Y_{red})_{red}$。
\end{lemma}

\begin{proof}
这来自引理 \ref{formal-spaces-lemma-reduction-formal-algebraic-space} 中约化的泛性质。
\end{proof}

\noindent
我们已经证明了下列引理（当时尚不知道纤维积存在）。

\begin{lemma}
\label{formal-spaces-lemma-diagonal-morphism-formal-algebraic-spaces}
设 $S$ 是概形，$f : X \to Y$ 是 $S$ 上形式代数空间的态射。
对角态射 $\Delta : X \to X \times_Y X$ 可表示（由概形表示）、
为单态射、局部拟有限、局部有限型且分离。
\end{lemma}

\begin{proof}
设 $T$ 是概形，并给定态射 $T \to X \times_Y X$。则
$$
T \times_{(X \times_Y X)} X = T \times_{(X \times_S X)} X
$$
所以结论立即来自引理 \ref{formal-spaces-lemma-diagonal-formal-algebraic-space}。
\end{proof}






\section{形式代数空间的分离公理}
\label{formal-spaces-section-separation}

\noindent
本节讨论形式代数空间的“绝对”分离条件。形式代数空间态射的分离条件将在后文讨论。

\begin{lemma}
\label{formal-spaces-lemma-characterize-quasi-separated}
设 $S$ 是概形，$X$ 是 $S$ 上的形式代数空间。下列条件等价：
\begin{enumerate}
\item $X$ 的约化（引理 \ref{formal-spaces-lemma-reduction-formal-algebraic-space}）
是拟分离代数空间；
\item 对 $U \to X$、$V \to X$，若 $U,V$ 是拟紧概形，
则纤维积 $U \times_X V$ 拟紧；
\item 对 $U \to X$、$V \to X$，若 $U,V$ 仿射，
则纤维积 $U \times_X V$ 拟紧。
\end{enumerate}
\end{lemma}

\begin{proof}
由引理 \ref{formal-spaces-lemma-diagonal-formal-algebraic-space}，
$U \times_X V$ 是概形。令 $U_{red},V_{red},X_{red}$ 分别为 $U,V,X$ 的约化。
则
$$
U_{red} \times_{X_{red}} V_{red} = U_{red} \times_X V_{red}
\to U \times_X V
$$
是概形的增厚。牢记代数空间的对应引理，(1) 与 (2) 的等价性由此显然，
见《空间的性质》之引理
\ref{spaces-properties-lemma-characterize-quasi-separated}。
我们略去 (2) 与 (3) 等价性的证明。
\end{proof}

\begin{lemma}
\label{formal-spaces-lemma-characterize-separated}
设 $S$ 是概形，$X$ 是 $S$ 上的形式代数空间。下列条件等价：
\begin{enumerate}
\item $X$ 的约化（引理 \ref{formal-spaces-lemma-reduction-formal-algebraic-space}）
是分离代数空间；
\item 对 $U \to X$、$V \to X$，若 $U,V$ 仿射，则纤维积
$U \times_X V$ 仿射，并且
$$
\mathcal{O}(U) \otimes_\mathbf{Z} \mathcal{O}(V)
\longrightarrow
\mathcal{O}(U \times_X V)
$$
是满射。
\end{enumerate}
\end{lemma}

\begin{proof}
若 (2) 成立，则把《空间的性质》之引理
% P10-E0315
\ref{spaces-properties-lemma-characterize-quasi-separated}
应用于态射 $U \to X_{red}$ 和 $V \to X_{red}$，其中 $U,V$ 仿射，
并使用 $U \times_{X_{red}}V = U \times_X V$，可知
$X_{red}$ 是分离代数空间。

\medskip\noindent
假设 (1)。设 $U \to X$ 和 $V \to X$ 如 (2) 所述。
由引理 \ref{formal-spaces-lemma-diagonal-formal-algebraic-space}，
$U \times_X V$ 是概形。令 $U_{red},V_{red},X_{red}$ 分别为 $U,V,X$ 的约化。
则
$$
U_{red} \times_{X_{red}} V_{red} = U_{red} \times_X V_{red}
\to U \times_X V
$$
是概形的增厚。于是
$(U \times_X V)_{red} = (U_{red} \times_{X_{red}} V_{red})_{red}$。
特别地，$(U \times_X V)_{red}$ 是仿射概形，并且
$$
\mathcal{O}(U) \otimes_\mathbf{Z} \mathcal{O}(V)
\longrightarrow
\mathcal{O}((U \times_X V)_{red})
$$
是满射，见《空间的性质》之引理
% P10-E0315
\ref{spaces-properties-lemma-characterize-quasi-separated}。
于是由《空间的极限》之命题
\ref{spaces-limits-proposition-affine}，$U \times_X V$ 仿射。
另一方面，仿射概形的态射 $U \times_X V \to U \times V$ 是复合
$$
U \times_X V = X \times_{(X \times_S X)} (U \times_S V)
\to U \times_S V \to U \times V
$$
第一个态射是单态射且局部有限型
（引理 \ref{formal-spaces-lemma-diagonal-formal-algebraic-space}）。
第二个态射是浸入（《概形》之引理
\ref{schemes-lemma-fibre-product-after-map}）。
所以该复合是局部有限型的单态射。另一方面，它是整的，因为底层既约仿射概形上的映射
由上述结论是闭浸入，从而普遍闭（使用《态射》之引理
\ref{morphisms-lemma-integral-universally-closed}）。因此环映射
$$
\mathcal{O}(U) \otimes_\mathbf{Z} \mathcal{O}(V)
\longrightarrow
\mathcal{O}(U \times_X V)
$$
是有限型的整满射态射，因而有限，进而满
（使用《态射》之引理 \ref{morphisms-lemma-finite-integral}
和《代数》之引理 \ref{algebra-lemma-finite-epimorphism-surjective}）。
\end{proof}

\begin{definition}
\label{formal-spaces-definition-separated}
设 $S$ 是概形，$X$ 是 $S$ 上的形式代数空间。我们称
\begin{enumerate}
\item 若引理 \ref{formal-spaces-lemma-characterize-quasi-separated} 的等价条件成立，
则 $X$ {\it 拟分离}；
\item 若引理 \ref{formal-spaces-lemma-characterize-separated} 的等价条件成立，
则 $X$ {\it 分离}。
\end{enumerate}
\end{definition}

\noindent
下列引理特别蕴含：弱可容许拓扑环的完备张量积仍是弱可容许拓扑环。

\begin{lemma}
\label{formal-spaces-lemma-fibre-product-affines-over-separated}
设 $S$ 是概形，$X \to Z$ 和 $Y \to Z$ 是 $S$ 上形式代数空间的态射。
假设 $Z$ 分离。
\begin{enumerate}
\item 若 $X$ 和 $Y$ 是仿射形式代数空间，则 $X \times_Z Y$ 也是。
\item 若 $X$ 和 $Y$ 是 McQuillan 仿射形式代数空间，则 $X \times_Z Y$ 也是。
\item 若 $X,Y,Z$ 是分别对应于弱可容许拓扑 $S$-代数
$A,B,C$ 的 McQuillan 仿射形式代数空间，则 $X \times_Z Y$
对应于 $A \widehat{\otimes}_C B$。
\end{enumerate}
\end{lemma}

\begin{proof}
按定义 \ref{formal-spaces-definition-affine-formal-algebraic-space} 写
$X = \colim X_\lambda$ 和 $Y = \colim Y_\mu$。则
$X \times_Z Y = \colim X_\lambda \times_Z Y_\mu$。
由于 $Z$ 分离，这些纤维积仿射，所以 (1) 成立。假设
% P10-E0316
$X$ 和 $Y$ 对应于弱可容许拓扑 $S$-代数 $A$ 和 $B$，并且
$X_\lambda = \Spec(A/I_\lambda)$、
$Y_\mu = \Spec(B/J_\mu)$。则
$$
X_\lambda \times_Z Y_\mu \to
X_\lambda \times Y_\mu \to \Spec(A \otimes B)
$$
是闭浸入。因此，引理
\ref{formal-spaces-lemma-mcquillan-affine-formal-algebraic-space} 的一个条件成立，
从而 $X \times_Z Y$ 是 McQuillan 的。
若 $Z$ 也为对应于 $C$ 的 McQuillan 空间，则
$$
X_\lambda \times_Z Y_\mu = \Spec(A/I_\lambda \otimes_C B/J_\mu)
$$
所以可见与 $X \times_Z Y$ 对应的弱可容许拓扑环是完备张量积
（见定义 \ref{formal-spaces-definition-toplogy-tensor-product}）。
\end{proof}

\begin{lemma}
\label{formal-spaces-lemma-separated-from-separated}
设 $S$ 是概形，$X$ 是 $S$ 上的形式代数空间。
设 $U \to X$ 是态射，其中 $U$ 是 $S$ 上的分离代数空间。
则 $U \to X$ 分离。
\end{lemma}

\begin{proof}
该陈述有意义，因为 $U \to X$ 可由代数空间表示
（引理 \ref{formal-spaces-lemma-space-to-formal-space}）。
设 $T$ 是概形并给定态射 $T \to X$。须证明
$U \times_X T \to T$ 分离。由于
$U \times_X T \to U \times_S T$ 是单态射，只须证明
$U \times_S T \to T$ 分离。后者是 $U \to S$ 的基变换，所以结论成立。
上述论证使用了《空间的态射》之引理
\ref{spaces-morphisms-lemma-base-change-separated}、
\ref{spaces-morphisms-lemma-composition-separated}、
\ref{spaces-morphisms-lemma-monomorphism-separated} 和
\ref{spaces-morphisms-lemma-separated-implies-morphism-separated}。
\end{proof}




\section{拟紧形式代数空间}
\label{formal-spaces-section-quasi-compact}

\noindent
下面刻画拟紧形式代数空间。

\begin{lemma}
\label{formal-spaces-lemma-characterize-quasi-compact}
设 $S$ 是概形，$X$ 是 $S$ 上的形式代数空间。下列条件等价：
\begin{enumerate}
\item $X$ 的约化（引理 \ref{formal-spaces-lemma-reduction-formal-algebraic-space}）
是拟紧代数空间；
\item 可以找到定义 \ref{formal-spaces-definition-formal-algebraic-space}
中的 $\{X_i \to X\}_{i \in I}$，其中 $I$ 有限；
\item 存在可由代数空间表示、étale 且满的态射 $Y \to X$，
其中 $Y$ 是仿射形式代数空间。
\end{enumerate}
\end{lemma}

\begin{proof}
略。
\end{proof}

\begin{definition}
\label{formal-spaces-definition-quasi-compact}
设 $S$ 是概形，$X$ 是 $S$ 上的形式代数空间。
若引理 \ref{formal-spaces-lemma-characterize-quasi-compact} 的等价条件成立，
则称 $X$ {\it 拟紧}。
\end{definition}

\begin{lemma}
\label{formal-spaces-lemma-characterize-quasi-compact-morphism}
设 $S$ 是概形，$f : X \to Y$ 是 $S$ 上形式代数空间的态射。
下列条件等价：
\begin{enumerate}
\item 约化之间的诱导映射
$f_{red} : X_{red} \to Y_{red}$
（引理 \ref{formal-spaces-lemma-reduction-formal-algebraic-space}）是代数空间的拟紧态射；
\item 对每个拟紧概形 $T$ 和态射 $T \to Y$，纤维积
$X \times_Y T$ 是拟紧形式代数空间；
\item 对每个仿射概形 $T$ 和态射 $T \to Y$，纤维积
$X \times_Y T$ 是拟紧形式代数空间；
\item 存在定义 \ref{formal-spaces-definition-formal-algebraic-space} 中的覆盖
$\{Y_j \to Y\}$，使每个 $X \times_Y Y_j$ 都是拟紧形式代数空间。
\end{enumerate}
\end{lemma}

\begin{proof}
略。
\end{proof}

\begin{definition}
\label{formal-spaces-definition-quasi-compact-morphism}
设 $S$ 是概形，$f : X \to Y$ 是 $S$ 上形式代数空间的态射。
若引理 \ref{formal-spaces-lemma-characterize-quasi-compact-morphism} 的等价条件成立，
则称 $f$ {\it 拟紧}。
\end{definition}

\noindent
当态射可由代数空间表示（特别地，当它可表示）时，
这与已有的概念一致。

\begin{lemma}
\label{formal-spaces-lemma-quasi-compact-representable}
设 $S$ 是概形，$f : X \to Y$ 是 $S$ 上形式代数空间的态射，
并且可由代数空间表示。则 $f$ 在定义
\ref{formal-spaces-definition-quasi-compact-morphism} 的意义下拟紧，当且仅当
$f$ 在《自举》之定义 \ref{bootstrap-definition-property-transformation}
的意义下拟紧。
\end{lemma}

\begin{proof}
这立即来自定义和引理 \ref{formal-spaces-lemma-characterize-quasi-compact-morphism}。
\end{proof}





\section{拟紧且拟分离的形式代数空间}
\label{formal-spaces-section-quasi-compact-quasi-separated}

\noindent
下列结果由 Yasuda 给出，见 \cite[Proposition 3.32]{Yasuda}。

\begin{lemma}
\label{formal-spaces-lemma-structure-quasi-compact-quasi-separated}
\begin{reference}
\cite[Proposition 3.32]{Yasuda}
\end{reference}
设 $S$ 是概形，$X$ 是 $S$ 上拟紧且拟分离的形式代数空间。
则对有向集 $\Lambda$ 上的某个代数空间系统
$(X_\lambda,f_{\lambda\mu})$，有 $X = \colim X_\lambda$，
其中每个 $f_{\lambda\mu} : X_\lambda \to X_\mu$ 都是增厚。
\end{lemma}

\begin{proof}
由引理 \ref{formal-spaces-lemma-characterize-quasi-compact}，可以选取仿射形式代数空间
$Y$ 以及可表示、满且 étale 的态射 $Y \to X$。
按定义 \ref{formal-spaces-definition-affine-formal-algebraic-space} 写
$Y = \colim Y_\lambda$。

\medskip\noindent
取 $\lambda \in \Lambda$。由引理 \ref{formal-spaces-lemma-presentation-representable}，
$Y_\lambda \times_X Y$ 是概形。$Y_\lambda \times_X Y$ 的约化
（引理 \ref{formal-spaces-lemma-reduction-formal-algebraic-space}）等于
$Y_{red} \times_{X_{red}}Y_{red}$ 的约化；后者拟紧，因为 $X$ 拟分离且
$Y_{red}$ 仿射。所以 $Y_\lambda \times_X Y$ 是拟紧概形。
因此存在 $\mu \geq \lambda$，使
$\text{pr}_2 : Y_\lambda \times_X Y \to Y$ 分解经过 $Y_\mu$，
见引理 \ref{formal-spaces-lemma-factor-through-thickening}。
令 $Z_\lambda$ 为态射
$\text{pr}_2 : Y_\lambda \times_X Y \to Y_\mu$ 的概形论像。
它与 $\mu$ 的选择无关，我们可以而且将把
$Z_\lambda \subset Y$ 看作态射
$\text{pr}_2 : Y_\lambda \times_X Y \to Y$ 的概形论像。
注意，$Z_\lambda$ 也等于态射
$\text{pr}_1 : Y \times_X Y_\lambda \to Y$ 的概形论像，
因为它同构于定义 $Z_\lambda$ 所用的态射。
我们断言，作为 $Y \times_X Y$ 的子函子，有
$Z_\lambda \times_X Y = Y \times_X Z_\lambda$。
事实上，因为 $Y \to X$ 是 étale 的，由《空间的态射》之引理
\ref{spaces-morphisms-lemma-quasi-compact-scheme-theoretic-image}，
$Z_\lambda \times_X Y$ 是态射
$$
\text{pr}_{13} = \text{pr}_1 \times \text{id}_Y :
Y \times_X Y_\lambda \times_X Y \longrightarrow Y \times_X Y
$$
的概形论像。同理，$Y \times_X Z_\lambda$ 是态射
$$
\text{pr}_{13} = \text{id}_Y \times \text{pr}_2 :
Y \times_X Y_\lambda \times_X Y \longrightarrow Y \times_X Y
$$
的概形论像。断言得证。于是
$R_\lambda = Z_\lambda \times_X Y = Y \times_X Z_\lambda$，
连同态射 $R_\lambda \to Z_\lambda \times_S Z_\lambda$，
定义一个 étale 等价关系。由此得到代数空间
$X_\lambda = Z_\lambda/R_\lambda$。按构造，图表
$$
\xymatrix{
Z_\lambda \ar[r] \ar[d] & Y \ar[d] \\
X_\lambda \ar[r] & X
}
$$
是笛卡尔的（因为 $X$ 是两个投影 $R=Y\times_XY\to Y$ 的余等化子，
$Z_\lambda\subset Y$ 在 $R$ 下不变，且 $R_\lambda$ 是 $R$ 到
$Z_\lambda$ 上的限制）。所以由《空间》之引理
\ref{spaces-lemma-morphism-sheaves-with-P-effective-descent-etale}，
$X_\lambda \to X$ 可表示且为闭浸入。另一方面，由于
$Y_\lambda \subset Z_\lambda$，有
$(X_\lambda)_{red} = X_{red}$；换言之，$X_\lambda \to X$ 是增厚。
最后，我们断言
$$
X = \colim X_\lambda
$$
我们有 $Y \times_X X_\lambda = Z_\lambda \supset Y_\lambda$。
每个态射 $T \to X$（其中 $T$ 是 $S$ 上的概形）都能 étale 局部地提升为到
$Y$ 的态射，而后者又能 étale 局部地提升为到某个 $Y_\lambda$ 的态射。
所以 $T \to X$ 在 $T$ 上 étale 局部地提升为到 $X_\lambda$ 的态射。
证明完成。
\end{proof}

\begin{remark}
\label{formal-spaces-remark-structure-quasi-compact-quasi-separated}
本注把引理 \ref{formal-spaces-lemma-structure-quasi-compact-quasi-separated}
的陈述和证明翻译为 EGA 意义下形式概形的语言。
考察注 \ref{formal-spaces-remark-ideals-of-definition}，可把引理翻译如下：
\begin{itemize}
\item[(*)] 每个拟紧且拟分离的形式概形都有一族定义理想基本系。
\end{itemize}
为证明这一点，先使用《概形的上同调》之引理
\ref{coherent-lemma-induction-principle} 的归纳原理
（为拟紧且拟分离的形式概形重新表述），把问题归约到如下情形：
$\mathfrak X = \mathfrak U \cup \mathfrak V$，其中
$\mathfrak U,\mathfrak V$ 是开形式子概形，$\mathfrak V$ 仿射，
且结果对 $\mathfrak U$、$\mathfrak V$ 和
$\mathfrak U\cap\mathfrak V$ 都成立。任取定义理想
$\mathcal{I}\subset\mathcal{O}_\mathfrak U$ 和
$\mathcal{J}\subset\mathcal{O}_\mathfrak V$。
根据 $\mathfrak U$ 和 $\mathfrak V$ 上都有定义理想基本系这一假设，
并利用 $\mathfrak U\cap\mathfrak V$ 拟紧，可以找到定义理想
$\mathcal{I}'\subset\mathcal{I}$ 和
$\mathcal{J}'\subset\mathcal{J}$，使得
$$
\mathcal{I}'|_{\mathfrak U \cap \mathfrak V} \subset
\mathcal{J}|_{\mathfrak U \cap \mathfrak V}
\quad\text{且}\quad
\mathcal{J}'|_{\mathfrak U \cap \mathfrak V} \subset
\mathcal{I}|_{\mathfrak U \cap \mathfrak V}
$$
设 $U \to U' \to \mathfrak U$ 和 $V \to V' \to \mathfrak V$
分别为定义理想
$\mathcal{I}'\subset\mathcal{I}\subset\mathcal{O}_\mathfrak U$
和 $\mathcal{J}'\subset\mathcal{J}\subset\mathcal{O}_\mathfrak V$
所确定的闭浸入。以 $\mathfrak U\cap V$ 表示 $V$ 的开子概形，
其底层拓扑空间为 $\mathfrak U\cap\mathfrak V$ 的底层拓扑空间。
由 $\mathcal{I}'$ 的选择，存在分解
$\mathfrak U\cap V\to U'$。类似地定义
$U\cap\mathfrak V$，它分解经过 $V'$。然后考虑
$$
Z_U = \text{下列态射的概形论像：}
U \amalg (\mathfrak U \cap V) \longrightarrow U'
$$
以及
$$
Z_V = \text{下列态射的概形论像：}
(U \cap \mathfrak V) \amalg V \longrightarrow V'
$$
由于取拟紧态射的概形论像与限制到开集交换
（《态射》之引理
\ref{morphisms-lemma-quasi-compact-scheme-theoretic-image}），可见
$Z_U\cap\mathfrak V = \mathfrak U\cap Z_V$。
所以 $Z_U$ 和 $Z_V$ 黏合成概形 $Z$，并带有态射 $Z\to\mathfrak X$。
与注 \ref{formal-spaces-remark-weak-ideals-of-definition} 中的讨论类似，
可见 $Z$ 对应于弱定义理想
$\mathcal{I}_Z\subset\mathcal{O}_\mathfrak X$。
注意 $Z_U\subset U'$ 且 $Z_V\subset V'$。因此，以这种方式构造的所有
$\mathcal{I}_Z$ 组成一族弱定义理想基本系。于是其中一个子族组成一族定义理想基本系，
见注 \ref{formal-spaces-remark-ideals-of-definition}。
\end{remark}

\begin{lemma}
\label{formal-spaces-lemma-characterize-affine}
设 $S$ 是概形，$X$ 是 $S$ 上的形式代数空间。
则 $X$ 是仿射形式代数空间，当且仅当其约化 $X_{red}$
（引理 \ref{formal-spaces-lemma-reduction-formal-algebraic-space}）仿射。
\end{lemma}

\begin{proof}
由引理 \ref{formal-spaces-lemma-characterize-quasi-separated}、
\ref{formal-spaces-lemma-characterize-quasi-compact} 以及定义
\ref{formal-spaces-definition-separated}、\ref{formal-spaces-definition-quasi-compact}，
可见 $X$ 拟紧且拟分离。由 Yasuda 引理
（引理 \ref{formal-spaces-lemma-structure-quasi-compact-quasi-separated}），
可以把 $X = \colim X_\lambda$ 写成沿代数空间增厚的滤过余极限。
然而，由《空间的极限》之引理
\ref{spaces-limits-lemma-reduction-scheme}，每个 $X_\lambda$ 都仿射，
因为 $(X_\lambda)_{red}=X_{red}$。所以按定义，$X$ 是仿射形式代数空间。
\end{proof}






\section{可由代数空间表示的态射}
\label{formal-spaces-section-representable}

\noindent
设 $f : X \to Y$ 是可由代数空间表示的形式代数空间态射。
借助代数空间各章中已经完成的工作，对于这类态射，我们有大量现成理论可用。

\begin{lemma}
\label{formal-spaces-lemma-composition-representable}
可由代数空间表示的态射之复合仍可由代数空间表示。
对可表示（由概形表示）的态射同样成立。
\end{lemma}

\begin{proof}
见《自举》之引理 \ref{bootstrap-lemma-composition-transformation}。
\end{proof}

\begin{lemma}
\label{formal-spaces-lemma-base-change-representable}
可由代数空间表示的态射的基变换仍可由代数空间表示。
对可表示（由概形表示）的态射同样成立。
\end{lemma}

\begin{proof}
见《自举》之引理 \ref{bootstrap-lemma-base-change-transformation}。
\end{proof}

\begin{lemma}
\label{formal-spaces-lemma-permanence-representable}
设 $S$ 是概形，$f : X \to Y$ 和 $g : Y \to Z$ 是 $S$ 上形式代数空间的态射。
若 $g \circ f : X \to Z$ 可由代数空间表示，则
$f : X \to Y$ 可由代数空间表示。
\end{lemma}

\begin{proof}
由引理 \ref{formal-spaces-lemma-diagonal-morphism-formal-algebraic-spaces}，
$Y \to Z$ 的对角态射可表示。因此，由《自举》之引理
\ref{bootstrap-lemma-representable-by-spaces-permanence}，
$X \to Y$ 可由代数空间表示。
\end{proof}

\noindent
可由代数空间表示这一性质在源和靶上都是局部的。

\begin{lemma}
\label{formal-spaces-lemma-representable-by-algebraic-spaces-local}
设 $S$ 是概形，$f : X \to Y$ 是 $S$ 上形式代数空间的态射。
下列条件等价：
\begin{enumerate}
\item 态射 $f$ 可由代数空间表示；
\item 存在交换图表
$$
\xymatrix{
U \ar[d] \ar[r] & V \ar[d] \\
X \ar[r] & Y
}
$$
其中 $U,V$ 是形式代数空间，竖直箭头可由代数空间表示，
$U \to X$ 是满 étale 态射，且 $U \to V$ 可由代数空间表示；
\item 对任意交换图表
$$
\xymatrix{
U \ar[d] \ar[r] & V \ar[d] \\
X \ar[r] & Y
}
$$
若 $U,V$ 是形式代数空间且竖直箭头可由代数空间表示，
则态射 $U \to V$ 可由代数空间表示；
\item 存在定义 \ref{formal-spaces-definition-formal-algebraic-space} 意义下的覆盖
$\{Y_j \to Y\}$，并且对每个 $j$，存在定义
\ref{formal-spaces-definition-formal-algebraic-space} 意义下的覆盖
$\{X_{ji} \to Y_j \times_Y X\}$，使得对每个 $j$ 和 $i$，
$X_{ji} \to Y_j$ 可由代数空间表示；
\item 存在定义 \ref{formal-spaces-definition-formal-algebraic-space} 意义下的覆盖
$\{X_i \to X\}$，并且对每个 $i$，存在分解
$X_i \to Y_i \to Y$，其中 $Y_i$ 是仿射形式代数空间，
$Y_i \to Y$ 可由代数空间表示，且 $X_i \to Y_i$ 可由代数空间表示；以及
% P10-E0318
\item 此处还可补充更多条件。
\end{enumerate}
\end{lemma}

\begin{proof}
显然 (1) 蕴含 (2)，因为可取 $U=X$ 和 $V=Y$。
反之，把《自举》之引理 \ref{bootstrap-lemma-representable-by-spaces-cover}
用于 $U \to X \to Y$，可知 (2) 蕴含 (1)。

\medskip\noindent
假设 (1) 成立，并考虑 (3) 中的图表。
则 $U \to Y$ 可由代数空间表示
（它是 $U \to X \to Y$ 的复合；见《自举》之引理
\ref{bootstrap-lemma-composition-transformation}），且分解经过 $V$。
所以由引理 \ref{formal-spaces-lemma-permanence-representable}，
$U \to V$ 可由代数空间表示。

\medskip\noindent
显然 (3) 蕴含 (2)。因此 (1)--(3) 现已等价。

\medskip\noindent
注意，(4) 中的条件是有意义的，因为由引理
\ref{formal-spaces-lemma-fibre-products}，纤维积 $Y_j \times_Y X$ 是形式代数空间。
显然 (4) 蕴含 (5)。

\medskip\noindent
假设有 (5) 中的 $X_i \to Y_i \to Y$。令
$V=\coprod Y_i$ 且 $U=\coprod X_i$，即可看出 (5) 蕴含 (2)。

\medskip\noindent
最后，假设 (1)--(3) 成立。于是可以任取定义
\ref{formal-spaces-definition-formal-algebraic-space} 意义下的覆盖
$\{Y_j \to Y\}$，并对每个 $j$ 任取定义
\ref{formal-spaces-definition-formal-algebraic-space} 意义下的覆盖
$\{X_{ji} \to Y_j \times_Y X\}$。
% P10-E0317
由 (3)，$X_{ij} \to Y_j$ 可由代数空间表示，故 (4) 成立。
证明完成。
\end{proof}

\begin{lemma}
\label{formal-spaces-lemma-algebraic-space-over-affine-formal}
设 $S$ 是概形，$Y$ 是 $S$ 上的仿射形式代数空间。
设 $f : X \to Y$ 是 $(\Sch/S)_{fppf}$ 上层的映射，并且可由代数空间表示。
则 $X$ 是形式代数空间。
\end{lemma}

\begin{proof}
按定义 \ref{formal-spaces-definition-affine-formal-algebraic-space} 写成
$Y = \colim Y_\lambda$。对每个 $\lambda$，纤维积
$X \times_Y Y_\lambda$ 都是代数空间。因此，由引理
\ref{formal-spaces-lemma-colimit-is-formal}，
$X = \colim X \times_Y Y_\lambda$ 是形式代数空间。
\end{proof}

\begin{lemma}
\label{formal-spaces-lemma-representable-by-algebraic-spaces}
设 $S$ 是概形，$Y$ 是 $S$ 上的形式代数空间。
设 $f : X \to Y$ 是 $(\Sch/S)_{fppf}$ 上层的映射，并且可由代数空间表示。
则 $X$ 是形式代数空间。
\end{lemma}

\begin{proof}
设 $\{Y_i \to Y\}$ 如定义
\ref{formal-spaces-definition-formal-algebraic-space} 所述。
则 $X \times_Y Y_i \to X$ 是一族可由代数空间表示、étale 且联合满的态射。
所以只需证明 $X \times_Y Y_i$ 是形式代数空间，见引理
\ref{formal-spaces-lemma-etale-covering-by-formal-algebraic-spaces}。
这由引理 \ref{formal-spaces-lemma-algebraic-space-over-affine-formal} 得出。
\end{proof}

\begin{lemma}
\label{formal-spaces-lemma-affine-representable-by-algebraic-spaces}
设 $S$ 是概形，$f : X \to Y$ 是仿射形式代数空间之间可由代数空间表示的态射。
则 $f$ 可表示（由概形表示）且为仿射态射。
\end{lemma}

\begin{proof}
我们将证明 $f$ 为仿射态射；于是由《空间的态射》之引理
\ref{spaces-morphisms-lemma-affine-local}，
$f$ 可表示且为仿射态射。按定义
\ref{formal-spaces-definition-affine-formal-algebraic-space} 写成
$Y = \colim Y_\mu$ 和 $X = \colim X_\lambda$。
设 $T \to Y$ 是态射，其中 $T$ 是 $S$ 上的概形。
由《自举》之定义 \ref{bootstrap-definition-property-transformation}，
必须证明 $X \times_Y T \to T$ 为仿射态射。
为此可以假设 $T$ 仿射，而需证明 $X \times_Y T$ 仿射。
此时，由引理 \ref{formal-spaces-lemma-factor-through-thickening}，
$T \to Y$ 对某个 $\mu$ 分解经过 $Y_\mu \to Y$。
由于 $f$ 拟紧，由引理 \ref{formal-spaces-lemma-characterize-quasi-compact-morphism}，
$X \times_Y T$ 拟紧。因此，$X \times_Y T \to X$ 对某个
$\lambda$ 分解经过 $X_\lambda$。类似地，增大 $\mu$ 后，
$X_\lambda \to Y$ 分解经过 $Y_\mu$。于是
$X \times_Y T = X_\lambda \times_{Y_\mu} T$。
因仿射概形的纤维积仿射，结论成立。
\end{proof}

\begin{lemma}
\label{formal-spaces-lemma-representable-affine}
设 $S$ 是概形，$\varphi : A \to B$ 是 $S$ 上弱可容许拓扑环的连续映射。
下列条件等价：
\begin{enumerate}
\item $\text{Spf}(\varphi) : \text{Spf}(B) \to \text{Spf}(A)$
可由代数空间表示；
\item $\text{Spf}(\varphi) : \text{Spf}(B) \to \text{Spf}(A)$
可表示（由概形表示）；
\item $\varphi$ 是 taut 的，见定义 \ref{formal-spaces-definition-taut}。
\end{enumerate}
\end{lemma}

\begin{proof}
由引理 \ref{formal-spaces-lemma-affine-representable-by-algebraic-spaces}，
(1) 和 (2) 等价。

\medskip\noindent
假设等价条件 (1) 和 (2) 成立。若 $I \subset A$ 是弱定义理想，则
$\Spec(A/I) \to \text{Spf}(A)$ 可表示且为增厚
（这从形式谱的构造显然可见，也可由引理
\ref{formal-spaces-lemma-mcquillan-affine-formal-algebraic-space} 得出）。
于是其基变换
$\Spec(A/I) \times_{\text{Spf}(A)} \text{Spf}(B) \to \text{Spf}(B)$
可表示且为增厚。所以由引理
\ref{formal-spaces-lemma-mcquillan-affine-formal-algebraic-space}，
存在弱定义理想 $J(I) \subset B$，使得作为
$\text{Spf}(B)$ 的子函子有
$\Spec(A/I) \times_{\text{Spf}(A)} \text{Spf}(B) = \Spec(B/J(I))$。
我们得到笛卡尔图表
$$
\xymatrix{
\Spec(B/J(I)) \ar[d] \ar[r] & \Spec(A/I) \ar[d] \\
\text{Spf}(B) \ar[r] & \text{Spf}(A)
}
$$
由引理 \ref{formal-spaces-lemma-fibre-product-affines-over-separated}，
$B/J(I) = B \widehat{\otimes}_A A/I$。
于是由引理 \ref{formal-spaces-lemma-closure-image-ideal}，
$J(I)$ 是理想 $\varphi(I)B$ 的闭包。
由于 $\text{Spf}(A) = \colim \Spec(A/I)$（其中 $I$ 如上），
可得 $\text{Spf}(B) = \colim \Spec(B/J(I))$。
因此，理想 $J(I)$ 组成一族弱定义理想基本系
（见引理 \ref{formal-spaces-lemma-mcquillan-affine-formal-algebraic-space}）。
所以 (3) 成立。

\medskip\noindent
假设 (3) 成立。实质上，我们只需反向进行上一段的论证。
设 $I \subset A$ 是弱定义理想。由引理
\ref{formal-spaces-lemma-fibre-product-affines-over-separated}，得到笛卡尔图表
$$
\xymatrix{
\text{Spf}(B \widehat{\otimes}_A A/I) \ar[d] \ar[r] & \Spec(A/I) \ar[d] \\
\text{Spf}(B) \ar[r] & \text{Spf}(A)
}
$$
若 $J(I)$ 是 $IB$ 的闭包，则由 $\varphi$ 的 taut 性，
$J(I)$ 在 $B$ 中是开的。因此，若 $J$ 在 $B$ 中是开的且
% P10-E0319
$J \subset J(B)$，则
$B/J \otimes_A A/I = B/(IB + J) = B/J(I)$，因为由引理
\ref{formal-spaces-lemma-closed}，
$J(I) = \bigcap_{J \subset B\text{ 开}} (IB + J)$。
所以定义完备张量积的极限稳定下来，给出
$B \widehat{\otimes}_A A/I = B/J(I)$。
因此
$\text{Spf}(B \widehat{\otimes}_A A/I) = \Spec(B/J(I))$。
这证明，对每个弱定义理想 $I \subset A$，
$\text{Spf}(B) \times_{\text{Spf}(A)} \Spec(A/I)$ 都可表示。
由于从拟紧概形 $T$ 出发的每个态射
$T \to \text{Spf}(A)$ 都对某个弱定义理想 $I$ 分解经过
$\Spec(A/I)$（引理 \ref{formal-spaces-lemma-factor-through-thickening}），
可知 $\text{Spf}(\varphi)$ 可表示，即 (2) 成立。证明完成。
\end{proof}

\begin{lemma}
\label{formal-spaces-lemma-property-goes-up-affine-morphism}
设 $S$ 是概形，$Y$ 是仿射形式代数空间。
设 $f : X \to Y$ 是 $(\Sch/S)_{fppf}$ 上层的映射，且可表示并为仿射态射。
则
\begin{enumerate}
\item $X$ 是仿射形式代数空间；
\item 若 $Y$ 可数指标化，则 $X$ 可数指标化；
\item 若 $Y$ 可数指标化且经典，则 $X$ 可数指标化且经典；
\item 若 $Y$ 弱 adic，则 $X$ 弱 adic；
\item 若 $Y$ 为 adic*，则 $X$ 为 adic*；以及
\item 若 $Y$ 是 Noether 的且 $f$（局部）有限型，则 $X$ 是 Noether 的。
\end{enumerate}
\end{lemma}

\begin{proof}
(1) 的证明。按定义 \ref{formal-spaces-definition-affine-formal-algebraic-space} 写成
$Y = \colim_{\lambda \in \Lambda} Y_\lambda$。
由于 $f$ 可表示且为仿射态射，纤维积
$X_\lambda = Y_\lambda \times_Y X$ 仿射，并且
% P10-E0320
$X = \colim Y_\lambda \times_Y X$。
所以 $X$ 是仿射形式代数空间。

\medskip\noindent
(2) 的证明。若 $Y$ 可数指标化，则在上述论证中可以假设
$\Lambda$ 可数。于是立即可见 $X$ 也可数指标化。

\medskip\noindent
(3)、(4) 和 (5) 的证明。在每种情形中，由假设和引理
\ref{formal-spaces-lemma-implications-between-types}，$Y$ 都是可数指标化的仿射形式代数空间；
再由 (2)，$X$ 亦然。所以由引理 \ref{formal-spaces-lemma-countably-indexed}，
可以对某些弱可容许拓扑 $S$-代数 $A,B$ 写成
$X = \text{Spf}(A)$ 和 $Y = \text{Spf}(B)$。
由引理 \ref{formal-spaces-lemma-morphism-between-formal-spectra}，态射 $f$ 对应于连续
$S$-代数同态 $\varphi : B \to A$。由引理
\ref{formal-spaces-lemma-representable-affine}，$\varphi$ 是 taut 的。
于是，(3) 由引理 \ref{formal-spaces-lemma-taut-ascent-admissible} 得出，
(4) 由引理 \ref{formal-spaces-lemma-taut-ascent-weakly-adic} 得出，
(5) 由引理 \ref{formal-spaces-lemma-taut-is-adic} 得出。

\medskip\noindent
(6) 的证明。结合 (3) 与引理 \ref{formal-spaces-lemma-implications-between-types}，
可见 $X$ 为 adic*。所以可以使用引理
\ref{formal-spaces-lemma-characterize-noetherian-affine} 的判据。
首先，它告诉我们仿射概形 $Y_\lambda$ 是 Noether 的。
于是 $X_\lambda \to Y_\lambda$ 有限型，故 $X_\lambda$ 也是 Noether 的
（《态射》之引理 \ref{morphisms-lemma-finite-type-noetherian}）。
再由该判据，$X$ 是 Noether 的，证明完成。
\end{proof}

\begin{lemma}
\label{formal-spaces-lemma-property-goes-up-affine}
设 $S$ 是概形，$f : X \to Y$ 是仿射形式代数空间之间可由代数空间表示的态射。
则
\begin{enumerate}
\item 若 $Y$ 可数指标化，则 $X$ 可数指标化；
\item 若 $Y$ 可数指标化且经典，则 $X$ 可数指标化且经典；
\item 若 $Y$ 弱 adic，则 $X$ 弱 adic；
\item 若 $Y$ 为 adic*，则 $X$ 为 adic*；以及
\item 若 $Y$ 是 Noether 的且 $f$（局部）有限型，则 $X$ 是 Noether 的。
\end{enumerate}
\end{lemma}

\begin{proof}
结合引理 \ref{formal-spaces-lemma-affine-representable-by-algebraic-spaces} 与
\ref{formal-spaces-lemma-property-goes-up-affine-morphism}。
\end{proof}

\begin{example}
\label{formal-spaces-example-representable-morphism-from-completion}
设 $B$ 是弱可容许拓扑环，$B \to A$ 是环映射（$A$ 上不带拓扑）。
于是可以考虑
$$
A^\wedge = \lim A/JA
$$
其中极限遍历 $B$ 的所有弱定义理想 $J$。
则 $A^\wedge$（赋予极限拓扑）是完备线性拓扑化环。
满射 $A^\wedge \to A/JA$ 的（开）核 $I$ 是 $JA^\wedge$ 的闭包，
见引理 \ref{formal-spaces-lemma-closed}。由引理
\ref{formal-spaces-lemma-topologically-nilpotent}，可见 $I$ 由拓扑幂零元组成。
所以 $I$ 是 $A^\wedge$ 的弱定义理想，进而 $A^\wedge$ 是弱可容许拓扑环。
% P10-E0321
因此，$\varphi : B \to A^\wedge$ 是弱可容许拓扑环的 taut 映射，且
$$
\text{Spf}(A^\wedge) \longrightarrow \text{Spf}(B)
$$
是引理 \ref{formal-spaces-lemma-representable-affine} 所研究现象的一个特例。
\end{example}

\begin{remark}[警告]
\label{formal-spaces-remark-warning}
引理 \ref{formal-spaces-lemma-representable-affine}、
\ref{formal-spaces-lemma-property-goes-up-affine-morphism} 和
\ref{formal-spaces-lemma-property-goes-up-affine} 中的讨论在以下两方面都是最精确的：
\begin{enumerate}
\item 若 $A,B$ 是弱可容许环且 $\varphi : A \to B$ 是连续映射，则一般而言，
$\text{Spf}(\varphi) : \text{Spf}(B) \to \text{Spf}(A)$ 不可表示。
\item 若 $f : Y \to X$ 是仿射形式代数空间之间的可表示态射，
且 $X = \text{Spf}(A)$ 是 McQuillan 的，则不能由此推出 $Y$ 是 McQuillan 的。
\end{enumerate}
(1) 的一个例子是取域 $A=k$（带离散拓扑）以及带 $t$-进拓扑的
$B=k[[t]]$。(2) 的一个例子见《例子》之第
\ref{examples-section-affine-formal-algebraic-space} 节。
\end{remark}

\noindent
尽管有上述警告，我们仍有如下结果。

\begin{lemma}
\label{formal-spaces-lemma-etale}
设 $S$ 是概形，$Y$ 是 $S$ 上的 McQuillan 仿射形式代数空间；
即对某个弱可容许拓扑 $S$-代数 $B$，有
$Y = \text{Spf}(B)$。则下列两个范畴等价：
\begin{enumerate}
\item 仿射形式代数空间之间可由代数空间表示且为 étale 的态射
$f : X \to Y$ 所成的范畴；以及
\item 形如 $A^\wedge$ 的拓扑 $B$-代数所成的范畴，其中
$A$ 是 étale $B$-代数，并且
$A^\wedge = \lim A/JA$，这里 $J \subset B$ 遍历 $B$ 的弱定义理想。
\end{enumerate}
该等价把 $A^\wedge$ 送到 $X = \text{Spf}(A^\wedge)$。
特别地，(1) 中的任意 $X$ 都是 McQuillan 的。
\end{lemma}

\begin{proof}
设 $A$ 是 étale $B$-代数。则对每个开理想 $J \subset B$，
$B/J \to A/JA$ 都是 étale 的。因此，态射
$\text{Spf}(A^\wedge) \to Y$ 可表示且为 étale。
由引理 \ref{formal-spaces-lemma-morphism-between-formal-spectra}，
函子 $\text{Spf}$ 全忠实。为完成证明，下一段将说明 (1) 中任意
$X \to Y$ 都属于本质像。

\medskip\noindent
选取弱定义理想 $J_0 \subset B$。令
$Y_0 = \Spec(B/J_0)$ 和 $X_0 = Y_0 \times_Y X$。
则 $X_0 \to Y_0$ 是仿射概形之间的 étale 态射
（见引理 \ref{formal-spaces-lemma-affine-representable-by-algebraic-spaces}）。
设 $X_0 = \Spec(A_0)$。由《代数》之引理
\ref{algebra-lemma-lift-etale}，可以找到 étale 代数映射
$B \to A$，使得 $A_0 \cong A/J_0A$。
% P10-E0322
考虑定义理想 $J \subset J_0$。如上，可以对某个 étale 环映射
$B/J \to \bar A$ 写成
$\Spec(B/J) \times_Y X = \Spec(\bar A)$。
于是 $B/J \to \bar A$ 和 $B/J \to A/JA$ 都是 étale 环映射，
并且都提升 étale 环映射 $B/J_0 \to A_0$。
由《更多代数》之引理
\ref{more-algebra-lemma-locally-nilpotent-henselian}，
存在唯一的 $B/J$-代数同构
$\varphi_J : A/JA \to \bar A$，它提升模 $J_0$ 的同一化。
由于映射 $\varphi_J$ 唯一，当 $J$ 变化时它们彼此相容。因此
% P10-E0323
$$
X = \colim \Spec(B/J) \times_Y X = \colim \Spec(A/JA) = \text{Spf}(A)
$$
故引理成立。
\end{proof}

\begin{lemma}
\label{formal-spaces-lemma-etale-surjective}
采用引理 \ref{formal-spaces-lemma-etale} 的记号和假设，设
$f : X \to Y$ 对应于 $B \to A^\wedge$。下列条件等价：
\begin{enumerate}
\item $f : X \to Y$ 满；
\item $B \to A$ 忠实平坦；
\item 对每个弱定义理想 $J \subset B$，环映射
$B/J \to A/JA$ 忠实平坦；以及
\item 对某个弱定义理想 $J \subset B$，环映射
$B/J \to A/JA$ 忠实平坦。
\end{enumerate}
\end{lemma}

\begin{proof}
设 $J \subset B$ 是弱定义理想。由于 $J$ 的每个元素都拓扑幂零，
$1+J$ 的每个元素都是单位。因此，$J$ 包含在 $B$ 的 Jacobson 根中
（《代数》之引理 \ref{algebra-lemma-contained-in-radical}）。
所以，平坦环映射 $B \to A$ 忠实平坦，当且仅当
$B/J \to A/JA$ 忠实平坦
（《代数》之引理 \ref{algebra-lemma-ff-rings}）。
由此可见 (2)--(4) 等价。
若 (1) 成立，则对每个弱定义理想 $J \subset B$，态射
$\Spec(A/JA) = \Spec(B/J) \times_Y X \to \Spec(B/J)$ 满，
这蕴含 (3)。反之，假设 (3)。从拟紧概形 $T$ 出发的态射
$T \to Y$ 对 $B$ 的某个定义理想 $J$ 分解经过 $\Spec(B/J)$
% P10-E0324
（引理 \ref{formal-spaces-lemma-factor-through-thickening}）。因此，
$X \times_Y T = \Spec(A/JA) \times_{\Spec(B/J)} T \to T$
是满态射 $\Spec(A/JA) \to \Spec(B/J)$ 的基变换，因而满。
所以 (1) 成立。
\end{proof}





\section{形式代数空间的类型}
\label{formal-spaces-section-types}

\noindent
本节定义“局部 Noether”、“局部 adic*”、“局部弱 adic”、
“局部可数指标化且经典”以及“局部可数指标化”的形式代数空间。
“局部 adic”、“局部经典”和“局部 McQuillan”这几种类型尚未定义，
因为我们不知道如何在这些情形中证明下列引理的类似结果
（只须对仿射形式代数空间之间的 étale 覆盖证明这些引理的类似结果就足够了）。

\begin{lemma}
\label{formal-spaces-lemma-iff-countably-indexed}
设 $S$ 是概形，$X \to Y$ 是仿射形式代数空间之间可由代数空间表示、
满且平坦的态射。则 $X$ 可数指标化，当且仅当 $Y$ 可数指标化。
\end{lemma}

\begin{proof}
假设 $X$ 可数指标化。按引理
\ref{formal-spaces-lemma-countable-affine-formal-algebraic-space} 写成
$X = \colim X_n$。按定义
\ref{formal-spaces-definition-affine-formal-algebraic-space} 写成
$Y = \colim Y_\lambda$。对每个 $n$，可以选取 $\lambda_n$，使得
$X_n \to Y$ 分解经过 $Y_{\lambda_n}$，见引理
\ref{formal-spaces-lemma-factor-through-thickening}。
另一方面，对每个 $\lambda$，概形 $Y_\lambda \times_Y X$ 仿射
（引理 \ref{formal-spaces-lemma-affine-representable-by-algebraic-spaces}），
所以 $Y_\lambda \times_Y X \to X$ 对某个 $n$ 分解经过 $X_n$
（引理 \ref{formal-spaces-lemma-factor-through-thickening}）。考虑图表
$$
\xymatrix{
Y_\lambda \times_Y X \ar[r] \ar[d] & X_n \ar[r] \ar[d] & X \ar[d] \\
Y_\lambda \ar@{..>}[r] \ar@/_1pc/[rr] & Y_{\lambda_n} \ar[r] & Y
}
$$
若能证明虚线箭头存在，则可推出
$Y = \colim Y_{\lambda_n}$，从而 $Y$ 可数指标化。
为此，选取 $\mu \geq \lambda$ 且 $\mu \geq \lambda_n$。
于是 $Y_\lambda \to Y$ 和 $Y_{\lambda_n} \to Y$ 都分解经过
$Y_\mu \to Y$。设 $Y_\mu = \Spec(B_\mu)$，闭子概形
$Y_\lambda$ 对应于 $J \subset B_\mu$，闭子概形
$Y_{\lambda_n}$ 对应于 $J' \subset B_\mu$。我们要证明 $J' \subset J$。
由上图可知 $J'A_\mu \subset JA_\mu$，其中
$Y_\mu \times_Y X = \Spec(A_\mu)$。
由于 $X \to Y$ 满且平坦，态射
$Y_\lambda \times_Y X \to Y_\lambda$ 是仿射概形之间的忠实平坦态射，
故 $B_\mu \to A_\mu$ 忠实平坦。因此如所需，$J' \subset J$。

\medskip\noindent
假设 $Y$ 可数指标化。则由引理
\ref{formal-spaces-lemma-property-goes-up-affine}，$X$ 可数指标化。
\end{proof}

\begin{lemma}
\label{formal-spaces-lemma-iff-cic}
设 $S$ 是概形，$X \to Y$ 是仿射形式代数空间之间可由代数空间表示、
满且平坦的态射。则 $X$ 可数指标化且经典，当且仅当
$Y$ 可数指标化且经典。
\end{lemma}

\begin{proof}
一个方向的蕴含已在引理 \ref{formal-spaces-lemma-property-goes-up-affine} 中证明。
对另一个方向，注意由引理 \ref{formal-spaces-lemma-iff-countably-indexed}，
可以假设 $X,Y$ 都可数指标化。因此，由引理
\ref{formal-spaces-lemma-countably-indexed}，对某些弱可容许拓扑 $S$-代数
$A,B$，有 $X = \text{Spf}(A)$ 和 $Y = \text{Spf}(B)$。
由引理 \ref{formal-spaces-lemma-morphism-between-formal-spectra}，
态射 $X \to Y$ 对应于连续 $S$-代数同态
$\varphi : B \to A$。由引理 \ref{formal-spaces-lemma-representable-affine}，
$\varphi$ 是 taut 的。设 $J \subset B$ 是开理想，
$I \subset A$ 是 $JA$ 的闭包。由引理
\ref{formal-spaces-lemma-fibre-product-affines-over-separated} 和
\ref{formal-spaces-lemma-closure-image-ideal}，可见
$\Spec(B/J) \times_Y X = \Spec(A/I)$。
因此，$B/J \to A/I$ 忠实平坦（因为 $X \to Y$ 满且平坦）。
这意味着 $\varphi : B \to A$ 满足第
\ref{formal-spaces-section-taut-descent} 节的条件（交换 $A$ 与 $B$ 的角色）。
所以由引理 \ref{formal-spaces-lemma-taut-descent-weakly-admissible}，结论成立。
\end{proof}

\begin{lemma}
\label{formal-spaces-lemma-iff-weakly-adic}
设 $S$ 是概形，$X \to Y$ 是仿射形式代数空间之间可由代数空间表示、
满且平坦的态射。则 $X$ 弱 adic，当且仅当 $Y$ 弱 adic。
\end{lemma}

\begin{proof}
证明与引理 \ref{formal-spaces-lemma-iff-cic} 的证明完全相同，
但最后使用引理 \ref{formal-spaces-lemma-taut-descent-weakly-adic}。
\end{proof}

\begin{lemma}
\label{formal-spaces-lemma-iff-adic-star}
设 $S$ 是概形，$X \to Y$ 是仿射形式代数空间之间可由代数空间表示、
满且平坦的态射。则 $X$ 为 adic*，当且仅当 $Y$ 为 adic*。
\end{lemma}

\begin{proof}
证明与引理 \ref{formal-spaces-lemma-iff-cic} 的证明完全相同，
但最后使用引理 \ref{formal-spaces-lemma-taut-descent-adic-star}。
\end{proof}

\begin{lemma}
\label{formal-spaces-lemma-iff-noetherian}
设 $S$ 是概形，$X \to Y$ 是仿射形式代数空间之间可由代数空间表示、
满、平坦且（局部）有限型的态射。则 $X$ 是 Noether 的，
当且仅当 $Y$ 是 Noether 的。
\end{lemma}

\begin{proof}
注意，Noether 仿射形式代数空间为 adic*，见引理
\ref{formal-spaces-lemma-implications-between-types}。所以由引理
\ref{formal-spaces-lemma-iff-adic-star}，可以假设 $X,Y$ 都为 adic*。
我们将使用引理 \ref{formal-spaces-lemma-characterize-noetherian-affine} 的判据证明结论。
具体地，按引理 \ref{formal-spaces-lemma-countable-affine-formal-algebraic-space} 写成
$Y = \colim Y_n$。对每个 $n$，令
$X_n = Y_n \times_Y X$。则 $X_n$ 是仿射概形
（引理 \ref{formal-spaces-lemma-affine-representable-by-algebraic-spaces}），且
$X = \colim X_n$。每个态射 $X_n \to Y_n$ 都忠实平坦且有限型。因此，引理由
$X_n$ 在这种情形中是 Noether 的当且仅当 $Y_n$ 是 Noether 的这一事实得出，见
《代数》之引理 \ref{algebra-lemma-descent-Noetherian}（下降）和
《代数》之引理 \ref{algebra-lemma-Noetherian-permanence}（上升）。
\end{proof}

\begin{lemma}
\label{formal-spaces-lemma-type-local}
设 $S$ 是概形。令
$$
P \in
\left\{
\begin{matrix}
\text{可数指标化},\\
\text{可数指标化且经典},\\
\text{弱 adic},\ adic*,\ \text{Noether}
\end{matrix}
\right\}
$$
设 $X$ 是 $S$ 上的形式代数空间。下列条件等价：
\begin{enumerate}
\item 若 $Y$ 是仿射形式代数空间，且
$f : Y \to X$ 可由代数空间表示并为 étale，则 $Y$ 具有性质 $P$；
\item 对定义 \ref{formal-spaces-definition-formal-algebraic-space} 中的某个
$\{X_i \to X\}_{i \in I}$，每个 $X_i$ 都具有性质 $P$。
\end{enumerate}
\end{lemma}

\begin{proof}
显然 (1) 蕴含 (2)。假设 (2)，并设 $Y \to X$ 如 (1) 所述。
由于纤维积 $X_i \times_X Y$ 是形式代数空间
（引理 \ref{formal-spaces-lemma-fibre-products-general}），可以选取定义
\ref{formal-spaces-definition-formal-algebraic-space} 意义下的覆盖
$\{X_{ij} \to X_i \times_X Y\}$。
由于 $Y$ 拟紧，存在 $(i_1,j_1),\ldots,(i_n,j_n)$，使得
$$
X_{i_1 j_1} \amalg \ldots \amalg X_{i_n j_n} \longrightarrow Y
$$
满且为 étale。于是 $X_{i_kj_k} \to X_{i_k}$ 可由代数空间表示且为 étale，
故由引理 \ref{formal-spaces-lemma-property-goes-up-affine}，
$X_{i_kj_k}$ 具有性质 $P$。因此
$X_{i_1 j_1} \amalg \ldots \amalg X_{i_n j_n}$ 是具有性质 $P$ 的
仿射形式代数空间（这里略去关于仿射形式代数空间有限不交并的一点细节）。
最后应用引理 \ref{formal-spaces-lemma-iff-countably-indexed}、
\ref{formal-spaces-lemma-iff-cic}、\ref{formal-spaces-lemma-iff-weakly-adic}、
\ref{formal-spaces-lemma-iff-adic-star} 和 \ref{formal-spaces-lemma-iff-noetherian} 之一即可得出结论。
\end{proof}

\noindent
上一引理为如下定义铺平了道路。

\begin{definition}
\label{formal-spaces-definition-types-formal-algebraic-spaces}
设 $S$ 是概形，$X$ 是 $S$ 上的形式代数空间。
如果引理 \ref{formal-spaces-lemma-type-local} 的等价条件对相应性质成立，
则称 $X$ \textit{局部可数指标化}、
\textit{局部可数指标化且经典}、
\textit{局部弱 adic}、\textit{局部 adic*} 或
\textit{局部 Noether}。
\end{definition}

\noindent
局部 Noether 代数空间沿闭子集的形式完备化是局部 Noether 形式代数空间。

\begin{lemma}
\label{formal-spaces-lemma-formal-completion-types}
设 $S$ 是概形，$X$ 是 $S$ 上的代数空间，$T \subset |X|$ 是闭子集。
设 $X_{/T}$ 是 $X$ 沿 $T$ 的形式完备化。
\begin{enumerate}
\item 若 $X \setminus T \to X$ 拟紧，则 $X_{/T}$ 局部 adic*。
\item 若 $X$ 局部 Noether，则 $X_{/T}$ 局部 Noether。
\end{enumerate}
\end{lemma}

\begin{proof}
选取满 étale 态射 $U \to X$，其中
$U=\coprod U_i$ 是仿射概形的不交并，见《空间的性质》之引理
\ref{spaces-properties-lemma-cover-by-union-affines}。
设 $T_i \subset U_i$ 是 $T$ 的逆像。由引理
\ref{formal-spaces-lemma-map-completions-representable}，有
$X_{/T} \times_X U_i = (U_i)_{/T_i}$。
所以 $\{(U_i)_{/T_i} \to X_{/T}\}$ 是定义
\ref{formal-spaces-definition-formal-algebraic-space} 意义下的覆盖。
此外，若 $X \setminus T \to X$ 拟紧，则
$U_i \setminus T_i \to U_i$ 亦然；若 $X$ 局部 Noether，
则 $U_i$ 亦然。因此，引理由仿射情形即引理
\ref{formal-spaces-lemma-affine-formal-completion-types} 得出。
\end{proof}

\begin{remark}[警告]
\label{formal-spaces-remark-warning-completion}
设 $X = \Spec(A)$，且 $T \subset X$ 是有限生成理想
$I \subset A$ 的零点集。令 $J=\sqrt I$ 为 $I$ 的根。
由定义可见 $X_{/T}=\text{Spf}(A^\wedge)$，其中
$A^\wedge=\lim A/I^n$ 是 $A$ 的 $I$-进完备化。
另一方面，从 $I$-进完备化到 $J$-进完备化的映射
$A^\wedge \to \lim A/J^n$ 可能不是环同构。例如，令
$$
A = \bigcup\nolimits_{n \geq 1} \mathbf{C}[t^{1/n}]
$$
且 $I=(t)$。则 $J=\mathfrak m$ 是赋值环 $A$ 的极大理想，
且 $J^2=J$。所以 $A$ 的 $J$-进完备化是 $\mathbf C$，
而 $I$-进完备化是例 \ref{formal-spaces-example-david-hansen} 所述的赋值环
（不过特别容易看出 $A \subset A^\wedge$）。
\end{remark}

\begin{lemma}
\label{formal-spaces-lemma-types-fibre-products}
% P10-E0325
设 $S$ 是概形，$X \to Y$ 和 $Z \to Y$ 是 $S$ 上形式代数空间的态射。
则
\begin{enumerate}
\item 若 $X,Z$ 局部可数指标化，则 $X \times_Y Z$ 局部可数指标化。
\item 若 $X,Z$ 局部可数指标化且经典，则
$X \times_Y Z$ 局部可数指标化且经典。
% P10-E0326
\item 若 $X,Z$ 弱 adic，则 $X \times_Y Z$ 弱 adic。
\item 若 $X,Z$ 局部 adic*，则 $X \times_Y Z$ 局部 adic*。
\item 若 $X,Z$ 局部 Noether，且 $X_{red} \to Y_{red}$ 局部有限型，
则 $X \times_Y Z$ 局部 Noether。
\end{enumerate}
\end{lemma}

\begin{proof}
选取定义 \ref{formal-spaces-definition-formal-algebraic-space} 意义下的覆盖
$\{Y_j \to Y\}$。对每个 $j$，选取定义
\ref{formal-spaces-definition-formal-algebraic-space} 意义下的覆盖
$\{X_{ji} \to Y_j \times_Y X\}$，并选取定义
\ref{formal-spaces-definition-formal-algebraic-space} 意义下的覆盖
$\{Z_{jk} \to Y_j \times_Y Z\}$。
由引理 \ref{formal-spaces-lemma-fibre-product-affines-over-separated}，
$X_{ji} \times_{Y_j} Z_{jk}$ 是仿射形式代数空间。因此
$$
\{X_{ji} \times_{Y_j} Z_{jk} \to X \times_Y Z\}
$$
是定义 \ref{formal-spaces-definition-formal-algebraic-space} 意义下的覆盖。
% P10-E0330
所以，只需在 $X,Y,Z$ 都是仿射形式代数空间的情形中证明 (1)、(2)、(3) 和 (4)。

\medskip\noindent
假设 $X,Z$ 可数指标化。按引理
\ref{formal-spaces-lemma-countable-affine-formal-algebraic-space} 写成
$X=\colim X_n$ 和 $Z=\colim Z_m$。
按定义 \ref{formal-spaces-definition-affine-formal-algebraic-space} 写成
$Y=\colim_{\lambda\in\Lambda}Y_\lambda$。
对每个 $n,m$，可以找到 $\lambda_{n,m}\in\Lambda$，使得
$X_n\to Y$ 和 $Z_m\to Y$ 都分解经过 $Y_{\lambda_{n,m}}$
（例如见引理 \ref{formal-spaces-lemma-factor-through-thickening}）。
选取 $\lambda_0\in\Lambda$。对 $t\geq1$ 归纳选取
$\lambda_t\in\Lambda$，使它对所有 $1\leq n,m\leq t$ 都满足
$\lambda_t\geq\lambda_{n,m}$，并且 $\lambda_t\geq\lambda_{t-1}$。
令 $Y'=\colim Y_{\lambda_t}$。则 $Y'\to Y$ 是单态射，且
$X\to Y$ 和 $Z\to Y$ 都分解经过 $Y'$。所以可以用 $Y'$ 代替 $Y$；
换言之，可以假设 $Y$ 可数指标化。

\medskip\noindent
假设 $X,Y,Z$ 可数指标化。由引理 \ref{formal-spaces-lemma-countably-indexed}，
可对某些弱可容许拓扑环 $A,B,C$ 写成
$X=\text{Spf}(A)$、$Y=\text{Spf}(B)$、$Z=\text{Spf}(C)$。
% P10-E0327
由引理 \ref{formal-spaces-lemma-morphism-between-formal-spectra}，
态射 $X\to Y$ 和 $Z\to Y$ 由连续环映射
$B\to A$ 和 $B\to C$ 给出。
由引理 \ref{formal-spaces-lemma-fibre-product-affines-over-separated}，可见
$X\times_Y Z=\text{Spf}(A\widehat{\otimes}_B C)$，并且
$A\widehat{\otimes}_B C$ 是弱可容许拓扑环。
特别地，由引理 \ref{formal-spaces-lemma-completed-tensor-product} 的第 (3) 部分，
$X\times_Y Z$ 可数指标化。这证明了 (1)。

\medskip\noindent
(2) 的证明。在这种情形中，$X,Z$ 可数指标化，故上述论证表明
$X\times_Y Z$ 是 $A\widehat{\otimes}_B C$ 的形式谱，其中 $A,C$ 可容许。
于是由引理 \ref{formal-spaces-lemma-completed-tensor-product} 的第 (2) 部分，
$A\widehat{\otimes}_B C$ 可容许。

\medskip\noindent
(3) 的证明。如前可知，$X\times_Y Z$ 是
$A\widehat{\otimes}_B C$ 的形式谱，其中 $A,C$ 弱 adic。
于是由引理 \ref{formal-spaces-lemma-completed-tensor-product-weakly-adic}，
$A\widehat{\otimes}_B C$ 弱 adic。

\medskip\noindent
(4) 的证明。与上面一样论证，结论由引理
\ref{formal-spaces-lemma-completed-tensor-product} 的第 (4) 部分得出。

\medskip\noindent
(5) 的证明。为了从引理 \ref{formal-spaces-lemma-completed-tensor-product} 的第 (5) 部分
推出情形 (5)，需要证明假设彼此吻合。
% P10-E0328
具体地，采用证明第一段的记号，若
$X_{red}\to Y_{red}$ 局部有限型，则
$(X_{ji})_{red}\to(Y_j)_{red}$ 局部有限型。
这由《空间的态射》之引理
\ref{spaces-morphisms-lemma-finite-type-local} 以及下列交换图表中竖直态射为
étale 这一事实得出：
$$
\xymatrix{
(X_{ji})_{red} \ar[d] \ar[r] & (Y_j)_{red} \ar[d] \\
X_{red} \ar[r] & Y_{red}
}
$$
事实上，由引理 \ref{formal-spaces-lemma-reduction-smooth}，
% P10-E0329
有 $(X_{ji})_{red}=X_{ij}\times_X X_{red}$ 和
$(Y_j)_{red}=Y_j\times_Y Y_{red}$。
所以如上可归约到如下情形：$X,Y,Z$ 是仿射形式代数空间，
$X,Z$ 是 Noether 的，且 $X_{red}\to Y_{red}$ 有限型。
其次，在证明第二段中我们用 $Y'$ 替换了 $Y$，但按构造
$Y_{red}=Y'_{red}$，故有限型假设在该替换下保持不变。
于是可见 $X,Y,Z$ 分别对应于 $A,B,C$，而
$X\times_Y Z$ 对应于 $A\widehat{\otimes}_B C$，其中 $A,C$ 是
Noether adic 环。最后，取约化对应于除去拓扑幂零元理想
（例 \ref{formal-spaces-example-reduction-affine-formal-spectrum}），
所以 $X_{red}\to Y_{red}$ 有限型确实意味着
$B/\mathfrak b\to A/\mathfrak a$ 有限型。证明完成。
\end{proof}

\begin{lemma}
\label{formal-spaces-lemma-structure-locally-noetherian}
设 $S$ 是概形，$X$ 是 $S$ 上的局部 Noether 形式代数空间。
则 $X=\colim X_n$，其中
$X_1\to X_2\to X_3\to\ldots$ 是 $S$ 上局部 Noether 代数空间的
有限阶增厚系统，$X_1=X_{red}$，并且对所有 $m\geq n$，
$X_n$ 是 $X_1$ 在 $X_m$ 中的第 $n$ 个无穷小邻域。
\end{lemma}

\begin{proof}
我们只概述证明，并略去若干细节。令 $X_1=X_{red}$。
按如下规则定义子函子 $X_n\subset X$：若 $T$ 是概形，
则态射 $f:T\to X$ 分解经过 $X_n$，当且仅当闭浸入
$X_1\times_XT\to T$ 的理想层的第 $n$ 次幂为零。
于是 $X_n\subset X$ 是子层，因为拟凝聚模的零化可以 fppf 局部地检验。
我们断言，$X_n\to X$ 可由概形表示且为闭浸入，并且作为 fppf 层有
$X=\colim X_n$。为检验这一点，可以在 $X$ 上 étale 局部地工作。
所以可假设 $X=\text{Spf}(A)$ 是 Noether 仿射形式代数空间。
于是 $X_1=\Spec(A/\mathfrak a)$，其中 $\mathfrak a\subset A$ 是
Noether adic 拓扑环 $A$ 的拓扑幂零元理想。此时
$X_n=\Spec(A/\mathfrak a^n)$，从而得到所需结论。
\end{proof}









\section{态射与连续环映射}
\label{formal-spaces-section-morphisms-rings}

\noindent
本节以 $\textit{WAdm}$ 表示弱可容许拓扑环与连续环同态所成的范畴。
我们定义全子范畴
$$
\textit{WAdm} \supset
\textit{WAdm}^{count} \supset
\textit{WAdm}^{cic} \supset
\textit{WAdm}^{weakly\ adic} \supset
\textit{WAdm}^{adic*} \supset
\textit{WAdm}^{Noeth}
$$
其对象分别为
\begin{enumerate}
\item $\textit{WAdm}^{count}$：具有可数开理想基本系的弱可容许拓扑环 $A$；
\item $\textit{WAdm}^{cic}$：具有可数开理想基本系的可容许拓扑环 $A$；
\item $\textit{WAdm}^{weakly\ adic}$：弱 adic 拓扑环
（第 \ref{formal-spaces-section-weakly-adic} 节）；
\item $\textit{WAdm}^{adic*}$：具有有限生成定义理想的 adic 拓扑环；以及
\item $\textit{WAdm}^{Noeth}$：Noether adic 拓扑环。
\end{enumerate}
显然，这些类型的环的形式谱分别是局部可数指标化、
局部可数指标化且经典、局部弱 adic、局部 adic* 以及局部 Noether
形式代数空间的基本构件。

\medskip\noindent
我们简要回顾可数指标化仿射形式代数空间的态射与
$\textit{WAdm}^{count}$ 中态射之间的关系。
设 $S$ 是概形，$X,Y$ 是可数指标化仿射形式代数空间。
% P10-E0331
写成 $X=\text{Spf}(A)$ 和 $Y=\text{Spf}(B)$，其中拓扑
$S$-代数 $A,B$ 属于 $\textit{WAdm}^{count}$，见引理
\ref{formal-spaces-lemma-countably-indexed}。由引理
\ref{formal-spaces-lemma-morphism-between-formal-spectra}，态射
$f:X\to Y$ 与拓扑 $S$-代数的连续映射
$$
\varphi : B \longrightarrow A
$$
之间存在一一对应。该对应由
$f\mapsto f^\sharp$ 和 $\varphi\mapsto\text{Spf}(\varphi)$ 给出。

\medskip\noindent
设 $S$ 是概形，$f:X\to Y$ 是局部可数指标化形式代数空间的态射。
考虑交换图表
$$
\xymatrix{
U \ar[d] \ar[r] & V \ar[d] \\
X \ar[r] & Y
}
$$
其中 $U,V$ 是仿射形式代数空间，且 $U\to X$ 和 $V\to Y$
可由代数空间表示并为 étale。由定义
\ref{formal-spaces-definition-types-formal-algebraic-spaces}
（因而也由引理 \ref{formal-spaces-lemma-type-local}），可知 $U,V$ 是可数指标化
仿射形式代数空间。由上一段的讨论，$U\to V$ 同构于
$\text{Spf}(\varphi)$，其中
$$
\varphi : B \longrightarrow A
$$
是 $\textit{WAdm}^{count}$ 中拓扑 $S$-代数的连续映射。

\begin{lemma}
\label{formal-spaces-lemma-completion-in-sub}
设 $A\in\Ob(\textit{WAdm})$，$A\to A'$ 是环映射（无拓扑）。
令
$(A')^\wedge=\lim_{I\subset A\text{ w.i.d}}A'/IA'$ 为例
\ref{formal-spaces-example-representable-morphism-from-completion} 中构造的
$\textit{WAdm}$ 对象。
\begin{enumerate}
\item 若 $A$ 属于 $\textit{WAdm}^{count}$，则 $(A')^\wedge$ 亦然。
\item 若 $A$ 属于 $\textit{WAdm}^{cic}$，则 $(A')^\wedge$ 亦然。
\item 若 $A$ 属于 $\textit{WAdm}^{weakly\ adic}$，则 $(A')^\wedge$ 亦然。
\item 若 $A$ 属于 $\textit{WAdm}^{adic*}$，则 $(A')^\wedge$ 亦然。
\item 若 $A$ 属于 $\textit{WAdm}^{Noeth}$ 且 $A'$ 是 Noether 的，
则 $(A')^\wedge$ 属于 $\textit{WAdm}^{Noeth}$。
\end{enumerate}
\end{lemma}

\begin{proof}
回忆 $A\to(A')^\wedge$ 是 taut 的，见例
\ref{formal-spaces-example-representable-morphism-from-completion} 中的讨论。
所以 (1)、(2)、(3) 和 (4) 分别由引理
\ref{formal-spaces-lemma-taut-ascent-countable}、
\ref{formal-spaces-lemma-taut-ascent-admissible}、
\ref{formal-spaces-lemma-taut-ascent-weakly-adic} 和
\ref{formal-spaces-lemma-taut-is-adic} 得出。
最后，假设 $A$ 是 Noether adic 环。由 (4)，
$(A')^\wedge$ 是 adic 的。由《代数》之引理
\ref{algebra-lemma-completion-Noetherian-Noetherian}，
$(A')^\wedge$ 是 Noether 的。所以 (5) 成立。
\end{proof}

\begin{situation}
\label{formal-spaces-situation-local-property}
设 $P$ 是 $\textit{WAdm}^{count}$ 态射的一个性质。
考虑交换图表
\begin{equation}
\label{formal-spaces-equation-localize}
\vcenter{
\xymatrix{
A \ar[r] & (A')^\wedge \\
B \ar[r] \ar[u]^\varphi & (B')^\wedge \ar[u]_{\varphi'}
}
}
\end{equation}
并假设满足下列条件：
\begin{enumerate}
\item $A,B$ 是 $\textit{WAdm}^{count}$ 的对象；
\item $A\to A'$ 和 $B\to B'$ 是 étale 环映射；
\item $(A')^\wedge=\lim A'/IA'$，而
$(B')^\wedge=\lim B'/JB'$，其中 $I\subset A$ 与 $J\subset B$
分别遍历 $A$ 与 $B$ 的弱可容许定义理想；
% P10-E0332
\item $\varphi:B\to A$ 和
$\varphi':(B')^\wedge\to(A')^\wedge$ 连续。
\end{enumerate}
由引理 \ref{formal-spaces-lemma-completion-in-sub}，拓扑环
$(A')^\wedge$ 和 $(B')^\wedge$ 是 $\textit{WAdm}^{count}$ 的对象。
若下列公理成立，则称 $P$ 是\textit{局部性质}：
\begin{enumerate}
\item
\label{formal-spaces-item-axiom-1}
对任意图表 (\ref{formal-spaces-equation-localize})，有
$P(\varphi)\Rightarrow P(\varphi')$；
\item
\label{formal-spaces-item-axiom-2}
对任意图表 (\ref{formal-spaces-equation-localize})，若 $A\to A'$ 忠实平坦，
则 $P(\varphi')\Rightarrow P(\varphi)$；
\item
\label{formal-spaces-item-axiom-3}
若对 $i=1,\ldots,n$ 都有 $P(B\to A_i)$，则
$P(B\to\prod_{i=1,\ldots,n}A_i)$。
\end{enumerate}
由于 $\textit{WAdm}^{count}$ 有有限积，公理
(\ref{formal-spaces-item-axiom-3}) 是有意义的。
\end{situation}

\begin{lemma}
\label{formal-spaces-lemma-property-defines-property-morphisms}
设 $S$ 是概形，$f:X\to Y$ 是 $S$ 上局部可数指标化形式代数空间的态射。
设 $P$ 是 $\textit{WAdm}^{count}$ 态射的局部性质。下列条件等价：
\begin{enumerate}
\item 对每个交换图表
$$
\xymatrix{
U \ar[d] \ar[r] & V \ar[d] \\
X \ar[r] & Y
}
$$
若 $U,V$ 是仿射形式代数空间，且 $U\to X$ 和 $V\to Y$
可由代数空间表示并为 étale，则态射 $U\to V$ 对应于
具有性质 $P$ 的 $\textit{WAdm}^{count}$ 态射；
\item 存在定义 \ref{formal-spaces-definition-formal-algebraic-space} 意义下的覆盖
$\{Y_j\to Y\}$，并且对每个 $j$ 存在定义
\ref{formal-spaces-definition-formal-algebraic-space} 意义下的覆盖
$\{X_{ji}\to Y_j\times_YX\}$，使每个 $X_{ji}\to Y_j$
都对应于具有性质 $P$ 的 $\textit{WAdm}^{count}$ 态射；以及
\item 存在定义 \ref{formal-spaces-definition-formal-algebraic-space} 意义下的覆盖
$\{X_i\to X\}$，并且对每个 $i$ 存在分解
$X_i\to Y_i\to Y$，其中 $Y_i$ 是仿射形式代数空间，
$Y_i\to Y$ 可由代数空间表示并为 étale，且 $X_i\to Y_i$
对应于具有性质 $P$ 的 $\textit{WAdm}^{count}$ 态射。
\end{enumerate}
\end{lemma}

\begin{proof}
显然 (1) 蕴含 (2)，且 (2) 蕴含 (3)。
假设 $\{X_i\to X\}$ 和 $X_i\to Y_i\to Y$ 如 (3) 所述，
并给定 (1) 中的图表。由于 $Y_i\times_YV$ 是形式代数空间
（引理 \ref{formal-spaces-lemma-fibre-products-general}），可以选取定义
\ref{formal-spaces-definition-formal-algebraic-space} 意义下的覆盖
$\{Y_{ij}\to Y_i\times_YV\}$。对每个 $(i,j)$，类似地选取定义
\ref{formal-spaces-definition-formal-algebraic-space} 意义下的覆盖
$\{X_{ijk}\to Y_{ij}\times_{Y_i}X_i\times_XU\}$。
由于 $U$ 拟紧，可以选择
$(i_1,j_1,k_1),\ldots,(i_n,j_n,k_n)$，使得
$$
X_{i_1 j_1 k_1} \amalg \ldots \amalg X_{i_n j_n k_n} \longrightarrow U
$$
满。对 $s=1,\ldots,n$，考虑交换图表
$$
\xymatrix{
& & & X_{i_s j_s k_s} \ar[ld] \ar[d] \ar[rd] \\
X \ar[d] & X_{i_s} \ar[l] \ar[d] &
X_{i_s} \times_X U \ar[l] \ar[d] & Y_{i_s j_s} \ar[ld] \ar[rd] &
X_{i_s} \times_X U \ar[d] \ar[r] &
U \ar[d] \ar[r] & X \ar[d] \\
Y & Y_{i_s} \ar[l] &
Y_{i_s} \times_Y V \ar[l] & &
Y_{i_s} \times_Y V \ar[r] &
V \ar[r] & Y
}
$$
若可数指标化仿射形式代数空间的一个态射所对应的
$\textit{WAdm}^{count}$ 态射具有 $P$，就称该态射具有 $P$。
注意，映射 $X_{i_sj_sk_s}\to X_{i_s}$ 和
$Y_{i_sj_s}\to Y_{i_s}$ 由 étale 环映射的完备化给出，
见引理 \ref{formal-spaces-lemma-etale}。所以由公理 (\ref{formal-spaces-item-axiom-1})，
$P(X_{i_s}\to Y_{i_s})$ 蕴含
$P(X_{i_sj_sk_s}\to Y_{i_sj_s})$。
注意，映射 $Y_{i_sj_s}\to V$ 由 étale 环映射的完备化给出
% P10-E0333
（仍用上一引理）。把公理 (\ref{formal-spaces-item-axiom-2}) 用于图表
$$
\xymatrix{
X_{i_s j_s k_s} \ar@{=}[r] \ar[d] & X_{i_s j_s k_s} \ar[d] \\
Y_{i_s j_s} \ar[r] & V
}
$$
（这是允许的，因为恒等映射是忠实平坦环映射），可得
$P(X_{i_sj_sk_s}\to V)$ 成立。由公理 (\ref{formal-spaces-item-axiom-3})，
$P(\coprod_{s=1,\ldots,n}X_{i_sj_sk_s}\to V)$ 成立。
按构造，态射 $\coprod X_{i_sj_sk_s}\to U$ 满，
所以由引理 \ref{formal-spaces-lemma-etale-surjective}，相应的
$\textit{WAdm}^{count}$ 态射是忠实平坦 étale 环映射的完备化。
再应用一次公理 (\ref{formal-spaces-item-axiom-2})（取 $B'=B$），
即可如所需得到 $P(U\to V)$ 成立。
\end{proof}

\begin{remark}[adic-star 变体]
\label{formal-spaces-remark-variant-adic-star}
设 $P$ 是 $\textit{WAdm}^{adic*}$ 态射的性质。
若公理 (\ref{formal-spaces-item-axiom-1})、(\ref{formal-spaces-item-axiom-2})、
(\ref{formal-spaces-item-axiom-3})（见情形 \ref{formal-spaces-situation-local-property}）
对 $\textit{WAdm}^{adic*}$ 的态射成立，则称 $P$ 是\textit{局部性质}。
以完全相同的方式，可得到引理
\ref{formal-spaces-lemma-property-defines-property-morphisms} 对 $S$ 上局部 adic*
形式代数空间之间态射的变体。
\end{remark}

\begin{remark}[Noether 变体]
\label{formal-spaces-remark-variant-Noetherian}
设 $P$ 是 $\textit{WAdm}^{Noeth}$ 态射的性质。
% P10-E0334
若公理 (\ref{formal-spaces-item-axiom-1})、(\ref{formal-spaces-item-axiom-2})、
(\ref{formal-spaces-item-axiom-3})（见情形 \ref{formal-spaces-situation-local-property}）
对 $\textit{WAdm}^{Noeth}$ 的态射成立，则称 $P$ 是\textit{局部性质}。
以完全相同的方式，可得到引理
\ref{formal-spaces-lemma-property-defines-property-morphisms} 对 $S$ 上局部 Noether
形式代数空间之间态射的变体。
\end{remark}

\begin{situation}
\label{formal-spaces-situation-base-change-local-property}
设 $P$ 是 $\textit{WAdm}^{count}$ 态射的局部性质，见情形
\ref{formal-spaces-situation-local-property}。若对 $\textit{WAdm}^{count}$ 中给定的
$B\to A$ 和 $B\to C$，有
$P(B\to A)\Rightarrow P(C\to A\widehat{\otimes}_BC)$，
则称 $P$ 在\textit{基变换下稳定}。
这是有意义的，因为由引理 \ref{formal-spaces-lemma-completed-tensor-product}，
$A\widehat{\otimes}_BC$ 是 $\textit{WAdm}^{count}$ 的对象。
\end{situation}

\begin{lemma}
\label{formal-spaces-lemma-base-change-property-morphisms}
设 $S$ 是概形，$P$ 是在基变换下稳定的
$\textit{WAdm}^{count}$ 态射的局部性质。
设 $f:X\to Y$ 和 $g:Z\to Y$ 是 $S$ 上局部可数指标化
形式代数空间的态射。若 $f$ 满足引理
\ref{formal-spaces-lemma-property-defines-property-morphisms} 的等价条件，
则 $\text{pr}_2:X\times_YZ\to Z$ 亦然。
\end{lemma}

\begin{proof}
选取定义 \ref{formal-spaces-definition-formal-algebraic-space} 意义下的覆盖
$\{Y_j\to Y\}$。对每个 $j$，选取定义
\ref{formal-spaces-definition-formal-algebraic-space} 意义下的覆盖
$\{X_{ji}\to Y_j\times_YX\}$，并选取定义
\ref{formal-spaces-definition-formal-algebraic-space} 意义下的覆盖
$\{Z_{jk}\to Y_j\times_YZ\}$。
由引理 \ref{formal-spaces-lemma-types-fibre-products}，
$X_{ji}\times_{Y_j}Z_{jk}$ 是可数指标化仿射形式代数空间。
于是
$$
\{X_{ji} \times_{Y_j} Z_{jk} \to X \times_Y Z\}
$$
是定义 \ref{formal-spaces-definition-formal-algebraic-space} 意义下的覆盖。
此外，态射 $X_{ji}\times_{Y_j}Z_{jk}\to Z$ 分解经过 $Z_{jk}$。
由假设，$X_{ji}\to Y_j$ 对应于具有性质 $P$ 的
$\textit{WAdm}^{count}$ 态射 $B_j\to A_{ji}$。
态射 $Z_{jk}\to Y_j$ 对应于 $\textit{WAdm}^{count}$ 中的态射
$B_j\to C_{jk}$。由引理
\ref{formal-spaces-lemma-fibre-product-affines-over-separated}，
$X_{ji} \times_{Y_j} Z_{jk} =
\text{Spf}(A_{ji} \widehat{\otimes}_{B_j} C_{jk})$，
所以只需证明
$C_{jk}\to A_{ji}\widehat{\otimes}_{B_j}C_{jk}$ 具有性质 $P$；
而这恰由 $P$ 在基变换下稳定这一条件保证。
\end{proof}

\begin{remark}[adic-star 变体]
\label{formal-spaces-remark-base-change-variant-adic-star}
设 $P$ 是 $\textit{WAdm}^{adic*}$ 态射的局部性质，见注
\ref{formal-spaces-remark-variant-adic-star}。若对 $\textit{WAdm}^{adic*}$ 中给定的
$B\to A$ 和 $B\to C$，有
$P(B\to A)\Rightarrow P(C\to A\widehat{\otimes}_BC)$，
则称 $P$ 在\textit{基变换下稳定}。
这是有意义的，因为由引理 \ref{formal-spaces-lemma-completed-tensor-product}，
$A\widehat{\otimes}_BC$ 是 $\textit{WAdm}^{adic*}$ 的对象。
以完全相同的方式，可得到引理
\ref{formal-spaces-lemma-base-change-property-morphisms} 对 $S$ 上局部 adic*
形式代数空间之间态射的变体。
\end{remark}

\begin{remark}[Noether 变体]
\label{formal-spaces-remark-base-change-variant-Noetherian}
设 $P$ 是 $\textit{WAdm}^{Noeth}$ 态射的局部性质，见注
\ref{formal-spaces-remark-variant-Noetherian}。若对 $\textit{WAdm}^{Noeth}$ 中给定的
$B\to A$ 和 $B\to C$，性质 $P(B\to A)$ 同时蕴含
$A\widehat{\otimes}_BC$ 是 Noether adic 的
\footnote{判据见引理 \ref{formal-spaces-lemma-completed-tensor-product}。}
以及 $P(C\to A\widehat{\otimes}_BC)$，则称 $P$ 在\textit{基变换下稳定}。
以完全相同的方式，可得到引理
\ref{formal-spaces-lemma-base-change-property-morphisms} 对 $S$ 上局部 Noether
形式代数空间之间态射的变体。
\end{remark}

\begin{remark}[另一个 Noether 变体]
\label{formal-spaces-remark-base-change-variant-variant-Noetherian}
设 $P,Q$ 是 $\textit{WAdm}^{Noeth}$ 态射的局部性质，见注
\ref{formal-spaces-remark-variant-Noetherian}。若对 $\textit{WAdm}^{Noeth}$ 中给定的
$B\to A$ 和 $B\to C$，在 $P(B\to A)$ 与 $Q(B\to C)$ 成立时，
$A\widehat{\otimes}_BC$ 是 Noether adic 的并且
$P(C\to A\widehat{\otimes}_BC)$ 成立，
则称\textit{$P$ 在按 $Q$ 的基变换下稳定}。
与引理 \ref{formal-spaces-lemma-base-change-property-morphisms} 的证明完全相同地论证，
得到如下陈述：给定 $S$ 上局部 Noether 形式代数空间的态射
% P10-E0335
$f:X\to Y$ 和 $g:Y\to Z$，使得
\begin{enumerate}
\item 引理 \ref{formal-spaces-lemma-property-defines-property-morphisms}
关于 $f$ 和 $P$ 的等价条件成立；
\item 引理 \ref{formal-spaces-lemma-property-defines-property-morphisms}
关于 $g$ 和 $Q$ 的等价条件成立，
\end{enumerate}
则引理 \ref{formal-spaces-lemma-property-defines-property-morphisms}
关于 $\text{pr}_2:X\times_YZ\to Z$ 和 $P$ 的等价条件成立。
\end{remark}

\begin{situation}
\label{formal-spaces-situation-composition-local-property}
设 $P$ 是 $\textit{WAdm}^{count}$ 态射的局部性质，见情形
\ref{formal-spaces-situation-local-property}。若对 $\textit{WAdm}^{count}$ 中给定的
$B\to A$ 和 $C\to B$，有
$P(B\to A)\wedge P(C\to B)\Rightarrow P(C\to A)$，
则称 $P$ 在\textit{复合下稳定}。
\end{situation}

\begin{lemma}
\label{formal-spaces-lemma-composition-property-morphisms}
设 $S$ 是概形，$P$ 是在复合下稳定的
$\textit{WAdm}^{count}$ 态射的局部性质。
设 $f:X\to Y$ 和 $g:Y\to Z$ 是 $S$ 上局部可数指标化形式代数空间的态射。
% P10-E0336
若 $f,g$ 满足引理 \ref{formal-spaces-lemma-property-defines-property-morphisms}
的等价条件，则 $g\circ f:X\to Z$ 亦然。
\end{lemma}

\begin{proof}
选取定义 \ref{formal-spaces-definition-formal-algebraic-space} 意义下的覆盖
$\{Z_k\to Z\}$。对每个 $k$，选取定义
\ref{formal-spaces-definition-formal-algebraic-space} 意义下的覆盖
$\{Y_{kj}\to Z_k\times_ZY\}$；再对每个 $k,j$ 选取定义
\ref{formal-spaces-definition-formal-algebraic-space} 意义下的覆盖
$\{X_{kji}\to Y_{kj}\times_YX\}$。
% P10-E0337
若 $f,g$ 满足引理 \ref{formal-spaces-lemma-property-defines-property-morphisms}
的等价条件，则
% P10-E0338
$X_{kji}\to Y_{jk}$ 和 $Y_{jk}\to Z_k$ 分别对应于
$\textit{WAdm}^{count}$ 中具有性质 $P$ 的箭头
$B_{kj}\to A_{kji}$ 和 $C_k\to B_{kj}$。
因此复合也具有该性质，结论成立。
\end{proof}

\begin{remark}[adic-star 变体]
\label{formal-spaces-remark-composition-variant-adic-star}
设 $P$ 是 $\textit{WAdm}^{adic*}$ 态射的局部性质，见注
\ref{formal-spaces-remark-variant-adic-star}。若对 $\textit{WAdm}^{adic*}$ 中给定的
$B\to A$ 和 $C\to B$，有
$P(B\to A)\wedge P(C\to B)\Rightarrow P(C\to A)$，
则称 $P$ 在\textit{复合下稳定}。
以完全相同的方式，可得到引理
\ref{formal-spaces-lemma-composition-property-morphisms} 对 $S$ 上局部 adic*
形式代数空间之间态射的变体。
\end{remark}

\begin{remark}[Noether 变体]
\label{formal-spaces-remark-composition-variant-Noetherian}
设 $P$ 是 $\textit{WAdm}^{Noeth}$ 态射的局部性质，见注
\ref{formal-spaces-remark-variant-Noetherian}。若对 $\textit{WAdm}^{Noeth}$ 中给定的
$B\to A$ 和 $C\to B$，有
$P(B\to A)\wedge P(C\to B)\Rightarrow P(C\to A)$，
则称 $P$ 在\textit{复合下稳定}。
以完全相同的方式，可得到引理
\ref{formal-spaces-lemma-composition-property-morphisms} 对 $S$ 上局部 Noether
形式代数空间之间态射的变体。
\end{remark}

\begin{situation}
\label{formal-spaces-situation-permanence-local-property}
设 $P$ 是 $\textit{WAdm}^{count}$ 态射的局部性质，见情形
\ref{formal-spaces-situation-local-property}。若对 $\textit{WAdm}^{count}$ 中给定的
$B\to A$ 和 $C\to B$，有
$P(C\to B)\wedge P(C\to A)\Rightarrow P(B\to A)$，
则称 $P$ \textit{具有消去性质}。
\end{situation}

\begin{lemma}
\label{formal-spaces-lemma-permanence-property-morphisms}
设 $S$ 是概形，$P$ 是具有消去性质的
$\textit{WAdm}^{count}$ 态射的局部性质。
设 $f:X\to Y$ 和 $g:Y\to Z$ 是 $S$ 上局部可数指标化形式代数空间的态射。
% P10-E0339
若 $g\circ f$ 和 $g$ 满足引理
\ref{formal-spaces-lemma-property-defines-property-morphisms} 的等价条件，
则 $f:X\to Y$ 亦然。
\end{lemma}

\begin{proof}
选取定义 \ref{formal-spaces-definition-formal-algebraic-space} 意义下的覆盖
$\{Z_k\to Z\}$。对每个 $k$，选取定义
\ref{formal-spaces-definition-formal-algebraic-space} 意义下的覆盖
$\{Y_{kj}\to Z_k\times_ZY\}$；再对每个 $k,j$ 选取定义
\ref{formal-spaces-definition-formal-algebraic-space} 意义下的覆盖
$\{X_{kji}\to Y_{kj}\times_YX\}$。
% P10-E0340
设 $X_{kji}\to Y_{jk}$ 和 $Y_{jk}\to Z_k$ 分别对应于
$\textit{WAdm}^{count}$ 中的箭头
$B_{kj}\to A_{kji}$ 和 $C_k\to B_{kj}$。
% P10-E0341
若 $g\circ f$ 和 $g$ 满足引理
\ref{formal-spaces-lemma-property-defines-property-morphisms} 的等价条件，
则 $C_k\to B_{kj}$ 和 $C_k\to A_{kji}$ 具有性质 $P$。
因此 $B_{kj}\to A_{kji}$ 也具有性质 $P$，结论成立。
\end{proof}

\begin{remark}[adic-star 变体]
\label{formal-spaces-remark-permanence-variant-adic-star}
设 $P$ 是 $\textit{WAdm}^{adic*}$ 态射的局部性质，见注
\ref{formal-spaces-remark-variant-adic-star}。若对 $\textit{WAdm}^{adic*}$ 中给定的
$B\to A$ 和 $C\to B$，有
$P(C\to A)\wedge P(C\to B)\Rightarrow P(B\to A)$，
则称 $P$ \textit{具有消去性质}。
以完全相同的方式，可得到引理
% P10-E0342
\ref{formal-spaces-lemma-composition-property-morphisms} 对 $S$ 上局部 adic*
形式代数空间之间态射的变体。
\end{remark}

\begin{remark}[Noether 变体]
\label{formal-spaces-remark-permanence-variant-Noetherian}
设 $P$ 是 $\textit{WAdm}^{Noeth}$ 态射的局部性质，见注
\ref{formal-spaces-remark-variant-Noetherian}。若对 $\textit{WAdm}^{Noeth}$ 中给定的
$B\to A$ 和 $C\to B$，有
% P10-E0343
$P(C\to B)\wedge P(C\to A)\Rightarrow P(C\to B)$，
则称 $P$ \textit{具有消去性质}。
以完全相同的方式，可得到引理
% P10-E0344
\ref{formal-spaces-lemma-composition-property-morphisms} 对 $S$ 上局部 Noether
形式代数空间之间态射的变体。
\end{remark}













\section{Taut 环映射与由代数空间表示}
\label{formal-spaces-section-taut}

\noindent
% P10-E0345
本节简要证明，局部可数指标化形式代数空间之间的态射，
étale 局部地对应于具有可数开理想基本系的弱可容许拓扑环之间的
taut 连续环同态。事实上，这从引理 \ref{formal-spaces-lemma-representable-affine}
已经相当清楚，我们鼓励读者跳过本节。

\begin{lemma}
\label{formal-spaces-lemma-representable-property-rings}
设 $B\to A$ 是 $\textit{WAdm}^{count}$ 中的箭头。下列条件等价：
\begin{enumerate}
\item[(a)] $B\to A$ 是 taut 的（定义 \ref{formal-spaces-definition-taut}）；
% P10-E0346
\item[(b)] 对弱定义理想基本系
$B\supset J_1\supset J_2\supset J_3\supset\ldots$，存在交换图表
$$
\xymatrix{
A \ar[r] & \ldots \ar[r] & A_3 \ar[r] & A_2 \ar[r] & A_1 \\
B \ar[r] \ar[u] & \ldots \ar[r] & B/J_3 \ar[r] \ar[u] &
B/J_2 \ar[r] \ar[u] & B/J_1 \ar[u]
}
$$
使得 $A_{n+1}/J_nA_{n+1}=A_n$，并且
% P10-E0347
$A=\lim A_n$ 作为拓扑环成立。
\end{enumerate}
此外，这些等价条件定义一个局部性质；即它们满足公理
(\ref{formal-spaces-item-axiom-1})、(\ref{formal-spaces-item-axiom-2})、(\ref{formal-spaces-item-axiom-3})。
\end{lemma}

\begin{proof}
(a) 与 (b) 的等价性是立即的。下面将给出这些公理的代数证明，
不过事实上我们已经证明过它们。具体地，利用引理
\ref{formal-spaces-lemma-representable-affine}，等价条件 (a)、(b) 转化为：
相应的仿射形式代数空间态射可由代数空间表示。
由引理 \ref{formal-spaces-lemma-representable-by-algebraic-spaces-local}，
这个条件“在源和靶上 étale 局部”，所以立即得到公理
(\ref{formal-spaces-item-axiom-1})、(\ref{formal-spaces-item-axiom-2}) 和
(\ref{formal-spaces-item-axiom-3})。

\medskip\noindent
公理 (\ref{formal-spaces-item-axiom-1})、(\ref{formal-spaces-item-axiom-2})、
(\ref{formal-spaces-item-axiom-3}) 的直接代数证明。
给定图表 (\ref{formal-spaces-equation-localize})，如情形
\ref{formal-spaces-situation-local-property} 所述。由例
\ref{formal-spaces-example-representable-morphism-from-completion}，
映射 $A\to(A')^\wedge$ 和 $B\to(B')^\wedge$ 满足 (a)、(b)。

\medskip\noindent
假设 $\varphi$ 满足 (a)、(b)。设 $J\subset B$ 是弱定义理想。
则 $JA$ 与 $J(B')^\wedge$ 的闭包分别是弱定义理想
$I\subset A$ 与 $J'\subset(B')^\wedge$。
于是 $I(A')^\wedge$ 的闭包是弱定义理想
$I'\subset(A')^\wedge$。拓扑论论证表明，$I'$ 也是
$J(A')^\wedge$ 与 $J'(A')^\wedge$ 的闭包。
最后，当 $J$ 遍历 $B$ 的弱定义理想基本系时，
$I$ 与 $I'$ 分别遍历 $A$ 与 $(A')^\wedge$ 中的弱定义理想基本系。
所以 (a) 对 $\varphi'$ 成立。这证明了 (\ref{formal-spaces-item-axiom-1})。

\medskip\noindent
假设 $A\to A'$ 忠实平坦，且 $\varphi'$ 满足 (a)、(b)。
设 $J\subset B$ 是弱定义理想。把 (a)、(b) 用于映射
$B\to(B')^\wedge\to(A')^\wedge$，可知
$J(A')^\wedge$ 的闭包 $I'$ 是弱定义理想。
特别地，$I'$ 是开的，故 $I'$ 在 $A$ 中的逆像是开的。
现在有（解释见下）
\begin{align*}
A \cap I'
& =
A \cap \bigcap (J(A')^\wedge + \Ker((A')^\wedge \to A'/I_0A')) \\
& =
A \cap \bigcap \Ker((A')^\wedge \to A'/JA' + I_0 A') \\
& = \bigcap (JA + I_0)
\end{align*}
由引理 \ref{formal-spaces-lemma-closed}，这就是 $JA$ 的闭包。
各交遍历弱定义理想 $I_0\subset A$。
% P10-E0348
第一个等式是因为 $(A')^\wedge$ 中 $0$ 的一族邻域基本系由映射
$(A')^\wedge\to A'/I_0A'$ 的核组成。第二个等式显然。
% P10-E0349
第三个等式是因为 $A\to A'$ 忠实平坦，见《代数》之引理
\ref{algebra-lemma-faithfully-flat-universally-injective}。
因此 $JA$ 的闭包是开的。由引理
\ref{formal-spaces-lemma-topologically-nilpotent}，$JA$ 的闭包是 $A$ 的弱定义理想。
最后，给定弱定义理想 $I\subset A$，由 $B\to(B')^\wedge$
和 $\varphi'$ 的性质 (a)，可找到 $J$，使
$J(A')^\wedge$ 包含在 $I(A')^\wedge$ 的闭包中。
所以 (a) 对 $\varphi$ 成立。这证明了 (\ref{formal-spaces-item-axiom-2})。

\medskip\noindent
我们略去 (\ref{formal-spaces-item-axiom-3}) 的证明。
\end{proof}

\begin{lemma}
\label{formal-spaces-lemma-representable-local-property}
设 $S$ 是概形，$f:X\to Y$ 是 $S$ 上局部可数指标化形式代数空间的态射。
下列条件等价：
\begin{enumerate}
\item 对每个交换图表
$$
\xymatrix{
U \ar[d] \ar[r] & V \ar[d] \\
X \ar[r] & Y
}
$$
若 $U,V$ 是仿射形式代数空间，且 $U\to X$ 和 $V\to Y$
可由代数空间表示并为 étale，则态射 $U\to V$ 对应于
$\textit{WAdm}^{count}$ 中的 taut 映射 $B\to A$；
\item 存在定义 \ref{formal-spaces-definition-formal-algebraic-space} 意义下的覆盖
$\{Y_j\to Y\}$，并且对每个 $j$，存在定义
\ref{formal-spaces-definition-formal-algebraic-space} 意义下的覆盖
$\{X_{ji}\to Y_j\times_YX\}$，使每个 $X_{ji}\to Y_j$
都对应于 $\textit{WAdm}^{count}$ 中的 taut 环映射；
\item 存在定义 \ref{formal-spaces-definition-formal-algebraic-space} 意义下的覆盖
$\{X_i\to X\}$，并且对每个 $i$ 存在分解
$X_i\to Y_i\to Y$，其中 $Y_i$ 是仿射形式代数空间，
$Y_i\to Y$ 可由代数空间表示并为 étale，且 $X_i\to Y_i$
对应于 $\textit{WAdm}^{count}$ 中的 taut 环映射；以及
\item $f$ 可由代数空间表示。
\end{enumerate}
\end{lemma}

\begin{proof}
由引理 \ref{formal-spaces-lemma-representable-property-rings}，
$\textit{WAdm}^{count}$ 中映射的“taut”性质是局部性质。
所以引理 \ref{formal-spaces-lemma-property-defines-property-morphisms}
恰好告诉我们 (1)、(2)、(3) 等价。
另一方面，由引理 \ref{formal-spaces-lemma-representable-affine}，
$\textit{WAdm}^{count}$ 中映射为“taut”恰好对应于相应的可数指标化
仿射形式代数空间态射“可由代数空间表示”。
% P10-E0350
因此，引理 \ref{formal-spaces-lemma-representable-by-algebraic-spaces-local} 的蕴含
(1) $\Rightarrow$ (2) 表明本引理的 (4) 蕴含 (1)。
类似地，引理 \ref{formal-spaces-lemma-representable-by-algebraic-spaces-local} 的蕴含
(4) $\Rightarrow$ (1) 表明本引理的 (2) 蕴含 (4)。
\end{proof}








\section{Adic 态射}
\label{formal-spaces-section-adic}

\noindent
本节把局部 adic* 形式代数空间 $X,Y$ 之间偶尔使用的“adic 态射”
$f:X\to Y$ 这一概念，一方面与 $f$ 可由代数空间表示相对应，
另一方面与我们的 taut 连续环同态概念相对应。
首先回忆：对具有有限生成定义理想的 adic 环，taut 性等价于 adic 性。

\begin{lemma}
\label{formal-spaces-lemma-adic-homomorphism}
设 $A,B$ 是 pre-adic 拓扑环，
$\varphi:A\to B$ 是连续环同态。
\begin{enumerate}
\item 若 $\varphi$ 是 adic 的，则 $\varphi$ 是 taut 的。
\item 若 $B$ 完备，$A$ 有有限生成定义理想，且 $\varphi$ 是 taut 的，
则 $\varphi$ 是 adic 的。
\end{enumerate}
特别地，在范畴 $\textit{WAdm}^{adic*}$ 上，“$\varphi$ 是 adic 的”
与“$\varphi$ 是 taut 的”这两个条件等价。
\end{lemma}

\begin{proof}
(1) 是引理 \ref{formal-spaces-lemma-adic-taut}，(2) 是引理
\ref{formal-spaces-lemma-taut-is-adic}。最后的陈述由 (1)、(2) 得出。
\end{proof}

\noindent
设 $S$ 是概形，$f:X\to Y$ 是 $S$ 上局部 adic* 形式代数空间的态射。
由引理 \ref{formal-spaces-lemma-representable-local-property}，下列条件等价：
\begin{enumerate}
\item $f$ 可由代数空间表示（换言之，引理
\ref{formal-spaces-lemma-representable-by-algebraic-spaces-local} 的等价条件成立）；
\item 对每个交换图表
$$
\xymatrix{
U \ar[d] \ar[r] & V \ar[d] \\
X \ar[r] & Y
}
$$
若 $U,V$ 是仿射形式代数空间，且 $U\to X$ 和 $V\to Y$
可由代数空间表示并为 étale，则态射 $U\to V$ 对应于
$\textit{WAdm}^{adic*}$ 中的 adic 映射
\footnote{由引理 \ref{formal-spaces-lemma-adic-homomorphism}，这等价于 taut。}。
\end{enumerate}
在这种情形中，我们将称\textit{$f$ 是 adic 态射}
（正式定义见下）。这一概念或术语只对局部 adic* 形式代数空间之间的态射定义和使用，
因为否则我们没有上述 (1)、(2) 的等价性。

\begin{definition}
\label{formal-spaces-definition-adic-morphism}
设 $S$ 是概形，$f:X\to Y$ 是 $S$ 上形式代数空间的态射。
假设 $X,Y$ 局部 adic*。若 $f$ 可由代数空间表示，
则称 $f$ 是\textit{adic 态射}。见上面的讨论。
\end{definition}









\section{有限型态射}
\label{formal-spaces-section-finite-type}

\noindent
% P10-E0351
鉴于 Stacks Project 中各项内容的设置方式，下面的做法才是真正正确的，
更强的概念应当另取名称。

\begin{definition}
\label{formal-spaces-definition-finite-type}
设 $S$ 是概形，$f:Y\to X$ 是 $S$ 上形式代数空间的态射。
\begin{enumerate}
\item 若 $f$ 可由代数空间表示，并且在《自举》之定义
\ref{bootstrap-definition-property-transformation} 的意义下局部有限型，
则称 $f$ \textit{局部有限型}。
\item 若 $f$ 局部有限型且拟紧
（定义 \ref{formal-spaces-definition-quasi-compact-morphism}），
则称 $f$ \textit{有限型}。
\end{enumerate}
\end{definition}

\noindent
我们将在第 \ref{formal-spaces-section-tft} 节讨论某些形式代数空间的有限型态射
与拓扑有限型连续环映射 $A\to B$ 之间的关系。

\begin{lemma}
\label{formal-spaces-lemma-characterize-finite-type}
设 $S$ 是概形，$f:X\to Y$ 是 $S$ 上形式代数空间的态射。
下列条件等价：
\begin{enumerate}
\item $f$ 有限型；
\item $f$ 可由代数空间表示，并且在《自举》之定义
\ref{bootstrap-definition-property-transformation} 的意义下有限型。
\end{enumerate}
\end{lemma}

\begin{proof}
这由《自举》之引理
\ref{bootstrap-lemma-transformations-property-implication}、
代数空间态射的蕴含“拟紧 $+$ 局部有限型 $\Rightarrow$ 有限型”，
以及引理 \ref{formal-spaces-lemma-quasi-compact-representable} 得出。
\end{proof}

\begin{lemma}
\label{formal-spaces-lemma-composition-finite-type}
有限型态射的复合有限型。对局部有限型同样成立。
\end{lemma}

\begin{proof}
见《自举》之引理
\ref{bootstrap-lemma-composition-transformation-property}，
并使用《空间的态射》之引理
\ref{spaces-morphisms-lemma-composition-finite-type}。
\end{proof}

\begin{lemma}
\label{formal-spaces-lemma-base-change-finite-type}
有限型态射的基变换有限型。对局部有限型同样成立。
\end{lemma}

\begin{proof}
见《自举》之引理
\ref{bootstrap-lemma-base-change-transformation-property}，
并使用《空间的态射》之引理
\ref{spaces-morphisms-lemma-base-change-finite-type}。
\end{proof}

\begin{lemma}
\label{formal-spaces-lemma-permanence-finite-type}
设 $S$ 是概形，$f:X\to Y$ 和 $g:Y\to Z$ 是 $S$ 上形式代数空间的态射。
若 $g\circ f:X\to Z$ 局部有限型，则 $f:X\to Y$ 局部有限型。
\end{lemma}

\begin{proof}
由引理 \ref{formal-spaces-lemma-permanence-representable}，$f$ 可由代数空间表示。
设 $T$ 是概形且 $T\to Z$ 是态射。把《空间的态射》之引理
\ref{spaces-morphisms-lemma-permanence-finite-type}
用于代数空间的态射
$T\times_ZX\to T\times_ZY\to T$，即可得出结论。
\end{proof}

\noindent
局部有限型这一性质在源和靶上都是局部的。

\begin{lemma}
\label{formal-spaces-lemma-finite-type-local}
设 $S$ 是概形，$f:X\to Y$ 是 $S$ 上形式代数空间的态射。
下列条件等价：
\begin{enumerate}
\item 态射 $f$ 局部有限型；
\item 存在交换图表
$$
\xymatrix{
U \ar[d] \ar[r] & V \ar[d] \\
X \ar[r] & Y
}
$$
其中 $U,V$ 是形式代数空间，竖直箭头可由代数空间表示并为 étale，
$U\to X$ 满，且 $U\to V$ 局部有限型；
\item 对任意交换图表
$$
\xymatrix{
U \ar[d] \ar[r] & V \ar[d] \\
X \ar[r] & Y
}
$$
若 $U,V$ 是形式代数空间，且竖直箭头可由代数空间表示并为 étale，
则态射 $U\to V$ 局部有限型；
\item 存在定义 \ref{formal-spaces-definition-formal-algebraic-space} 意义下的覆盖
$\{Y_j\to Y\}$，并且对每个 $j$，存在定义
\ref{formal-spaces-definition-formal-algebraic-space} 意义下的覆盖
$\{X_{ji}\to Y_j\times_YX\}$，使得对每个 $j,i$，
$X_{ji}\to Y_j$ 局部有限型；
\item 存在定义 \ref{formal-spaces-definition-formal-algebraic-space} 意义下的覆盖
$\{X_i\to X\}$，并且对每个 $i$ 存在分解
$X_i\to Y_i\to Y$，其中 $Y_i$ 是仿射形式代数空间，
$Y_i\to Y$ 可由代数空间表示并为 étale，且 $X_i\to Y_i$ 局部有限型；以及
% P10-E0352
\item 此处还可补充更多条件。
\end{enumerate}
\end{lemma}

\begin{proof}
% P10-E0353
在这 5 种情形中的每一种，态射 $f:X\to Y$ 都可由代数空间表示，
见引理 \ref{formal-spaces-lemma-representable-by-algebraic-spaces-local}。
以下将不再说明地使用这一点。

\medskip\noindent
显然 (1) 蕴含 (2)，因为可以取 $U=X$ 和 $V=Y$。
反之，假设给定 (2) 中的图表。设 $T$ 是概形且 $T\to Y$ 是态射。
则可以考虑
$$
\xymatrix{
U \times_Y T \ar[d] \ar[r] & V \times_Y T \ar[d] \\
X \times_Y T \ar[r] & T
}
$$
竖直箭头为 étale，顶部水平箭头局部有限型，因为它们是相应态射的基变换。
% P10-E0354
所以由《空间的态射》之引理
\ref{spaces-morphisms-lemma-finite-type-local}，
$X\times_YT\to T$ 局部有限型。换言之，(1) 成立。

\medskip\noindent
假设 (1) 成立，并考虑 (3) 中的图表。则 $U\to Y$ 局部有限型
（它是 $U\to X\to Y$ 的复合；见《自举》之引理
\ref{bootstrap-lemma-composition-transformation-property}）。
设 $T$ 是概形且 $T\to V$ 是态射。则投影
$T\times_VU\to T$ 分解为
$$
T \times_V U = (T \times_Y U) \times_{(V \times_Y V)} V
\to T \times_Y U \to T
$$
第二个箭头局部有限型（它是复合 $U\to X\to Y$ 的基变换），
第一个箭头是对角态射 $V\to V\times_YV$ 的基变换；
由引理 \ref{formal-spaces-lemma-diagonal-morphism-formal-algebraic-spaces}，
该对角态射局部有限型。

\medskip\noindent
显然 (3) 蕴含 (2)。因此 (1)--(3) 现已等价。

\medskip\noindent
注意，(4) 中的条件是有意义的，因为由引理
\ref{formal-spaces-lemma-fibre-products}，纤维积 $Y_j\times_YX$ 是形式代数空间。
显然 (4) 蕴含 (5)。

\medskip\noindent
假设有 (5) 中的 $X_i\to Y_i\to Y$。令
$V=\coprod Y_i$ 且 $U=\coprod X_i$，即可看出 (5) 蕴含 (2)。

\medskip\noindent
最后，假设 (1)--(3) 成立。于是可以任取定义
\ref{formal-spaces-definition-formal-algebraic-space} 意义下的覆盖
$\{Y_j\to Y\}$，并对每个 $j$ 任取定义
\ref{formal-spaces-definition-formal-algebraic-space} 意义下的覆盖
$\{X_{ji}\to Y_j\times_YX\}$。
% P10-E0355
由 (3)，$X_{ij}\to Y_j$ 局部有限型，故 (4) 成立。证明完成。
\end{proof}

\begin{example}
\label{formal-spaces-example-finite-type-from-finite-type-ring-map}
设 $S$ 是概形，$A$ 是 $S$ 上的弱可容许拓扑环，
$A\to A'$ 是有限型环映射。则
$$
(A')^\wedge = \lim_{I \subset A\ w.i.d.} A'/IA'
$$
是弱可容许环，且相应的态射
$\text{Spf}((A')^\wedge)\to\text{Spf}(A)$ 可表示，
见例 \ref{formal-spaces-example-representable-morphism-from-completion}。
若 $T\to\text{Spf}(A)$ 是态射且 $T$ 是拟紧概形，则由引理
\ref{formal-spaces-lemma-factor-through-thickening}，它对某个弱定义理想
$I\subset A$ 分解经过 $\Spec(A/I)$。于是
$T\times_{\text{Spf}(A)}\text{Spf}((A')^\wedge)$ 等于
$T\times_{\Spec(A/I)}\Spec(A'/IA')$，故
$\text{Spf}((A')^\wedge)\to\text{Spf}(A)$ 有限型。
\end{example}

\begin{lemma}
\label{formal-spaces-lemma-locally-finite-type-locally-noetherian}
设 $S$ 是概形，$f:X\to Y$ 是 $S$ 上形式代数空间的态射。
% P10-E0356
若 $Y$ 局部 Noether 且 $f$ 局部有限型，则 $X$ 局部 Noether。
\end{lemma}

\begin{proof}
按引理 \ref{formal-spaces-lemma-finite-type-local} 选取
$\{Y_j\to Y\}$ 和 $\{X_{ij}\to Y_j\times_YX\}$。
由引理 \ref{formal-spaces-lemma-property-goes-up-affine-morphism}，
每个 $X_{ij}$ 都是 Noether 的。引理得证。
\end{proof}

\begin{lemma}
\label{formal-spaces-lemma-fibre-product-Noetherian}
设 $S$ 是概形，$f:X\to Y$ 和 $Z\to Y$ 是 $S$ 上形式代数空间的态射。
若 $Z$ 局部 Noether 且 $f$ 局部有限型，则
$Z\times_YX$ 局部 Noether。
\end{lemma}

\begin{proof}
由引理 \ref{formal-spaces-lemma-base-change-finite-type}，
态射 $Z\times_YX\to Z$ 局部有限型。
所以结论由引理 \ref{formal-spaces-lemma-locally-finite-type-locally-noetherian} 得出。
\end{proof}








\section{满态射}
\label{formal-spaces-section-surjective}

\noindent
由引理 \ref{formal-spaces-lemma-reduction-surjective}，下述定义不会与已有的代数空间态射定义，
或可由代数空间表示的形式代数空间态射定义冲突。

\begin{definition}
\label{formal-spaces-definition-surjective}
设 $S$ 为概形。若 $S$ 上形式代数空间的态射 $f:X\to Y$
在底层约化代数空间上诱导满态射 $X_{red}\to Y_{red}$，
则称 $f$ 为\textit{满态射}。
\end{definition}

\begin{lemma}
\label{formal-spaces-lemma-composition-surjective}
两个满态射的复合是满态射。
\end{lemma}

\begin{proof}
省略。
\end{proof}

\begin{lemma}
\label{formal-spaces-lemma-base-change-surjective}
满态射的基变换是满态射。
\end{lemma}

\begin{proof}
省略。
\end{proof}

\begin{lemma}
\label{formal-spaces-lemma-characterize-surjective}
设 $S$ 为概形，$f:X\to Y$ 为 $S$ 上形式代数空间的态射。
下列条件等价：
\begin{enumerate}
\item $f$ 是满态射；
\item 对每个概形 $T$ 以及态射 $T\to Y$，投影
$X\times_YT\to T$ 是形式代数空间的满态射；
\item 对每个仿射概形 $T$ 以及态射 $T\to Y$，投影
$X\times_YT\to T$ 是形式代数空间的满态射；
\item 存在定义 \ref{formal-spaces-definition-formal-algebraic-space} 意义下的覆盖
$\{Y_j\to Y\}$，使得每个 $X\times_YY_j\to Y_j$ 都是形式代数空间的满态射；
\item 存在形式代数空间的满态射 $Z\to Y$，使得
$X\times_YZ\to Z$ 是满态射，并且
% P10-E0357
\item 此处还可补充更多条件。
\end{enumerate}
\end{lemma}

\begin{proof}
省略。
\end{proof}








\section{单态射}
\label{formal-spaces-section-monomorphisms}

\noindent
定义如下。

\begin{definition}
\label{formal-spaces-definition-monomorphism}
设 $S$ 为概形。若 $S$ 上形式代数空间的态射是层之间的单射映射，
则称其为\textit{单态射}。
\end{definition}

\noindent
下面是一个例子。设 $X$ 为代数空间，且 $T\subset |X|$ 为闭子集。
那么，从沿 $T$ 的 $X$ 的形式完备化到 $X$ 的态射
$X_{/T}\to X$ 是单态射。特别地，形式代数空间的单态射一般并不
可由代数空间表示。

\begin{lemma}
\label{formal-spaces-lemma-composition-monomorphism}
两个单态射的复合是单态射。
\end{lemma}

\begin{proof}
省略。
\end{proof}

\begin{lemma}
\label{formal-spaces-lemma-base-change-monomorphism}
单态射的基变换是单态射。
\end{lemma}

\begin{proof}
省略。
\end{proof}

\begin{lemma}
\label{formal-spaces-lemma-characterize-monomorphisms}
设 $S$ 为概形，$f:X\to Y$ 为 $S$ 上形式代数空间的态射。
下列条件等价：
\begin{enumerate}
\item $f$ 是单态射；
\item 对每个概形 $T$ 以及态射 $T\to Y$，投影
$X\times_YT\to T$ 是形式代数空间的单态射；
\item 对每个仿射概形 $T$ 以及态射 $T\to Y$，投影
$X\times_YT\to T$ 是形式代数空间的单态射；
\item 存在定义 \ref{formal-spaces-definition-formal-algebraic-space} 意义下的覆盖
$\{Y_j\to Y\}$，使得每个 $X\times_YY_j\to Y_j$ 都是形式代数空间的单态射，
并且
\item 存在一族态射 $\{Y_j\to Y\}$，使得
$\coprod Y_j\to Y$ 是 $(\Sch/S)_{fppf}$ 上层的满射，且每个
$X\times_YY_j\to Y_j$ 对所有 $j$ 都是单态射；
\item 存在形式代数空间的态射 $Z\to Y$，该态射可由代数空间表示、满、平坦且局部有限呈示，
使得
% P10-E0359
$X\times_YZ\to X$ 是单态射，并且
% P10-E0358
\item 此处还可补充更多条件。
\end{enumerate}
\end{lemma}

\begin{proof}
省略。
\end{proof}







\section{闭浸入}
\label{formal-spaces-section-closed-immersions}

\noindent
定义如下。

\begin{definition}
\label{formal-spaces-definition-closed-immersion}
设 $S$ 为概形，设 $f:Y\to X$ 为 $S$ 上形式代数空间的态射。
若 $f$ 可由代数空间表示，并且在《自举》之定义
\ref{bootstrap-definition-property-transformation} 的意义下是闭浸入，
则称 $f$ 为\textit{闭浸入}。
\end{definition}

\noindent
% P10-E0360
阅读本节时，请跳过开头那些必需的引理。

\begin{lemma}
\label{formal-spaces-lemma-composition-closed-immersion}
两个闭浸入的复合是闭浸入。
\end{lemma}

\begin{proof}
省略。
\end{proof}

\begin{lemma}
\label{formal-spaces-lemma-base-change-closed-immersion}
闭浸入的基变换是闭浸入。
\end{lemma}

\begin{proof}
省略。
\end{proof}

\begin{lemma}
\label{formal-spaces-lemma-characterize-closed-immersion}
设 $S$ 为概形，$f:X\to Y$ 为 $S$ 上形式代数空间的态射。
下列条件等价：
\begin{enumerate}
\item $f$ 是闭浸入；
\item 对每个概形 $T$ 以及态射 $T\to Y$，投影
$X\times_YT\to T$ 是闭浸入；
\item 对每个仿射概形 $T$ 以及态射 $T\to Y$，投影
$X\times_YT\to T$ 是闭浸入；
\item 存在定义 \ref{formal-spaces-definition-formal-algebraic-space} 意义下的覆盖
$\{Y_j\to Y\}$，使得每个 $X\times_YY_j\to Y_j$ 都是闭浸入，并且
\item 存在形式代数空间的态射 $Z\to Y$，该态射可由代数空间表示、满、平坦且局部有限呈示，
使得
% P10-E0361
$X\times_YZ\to X$ 是闭浸入，并且
% P10-E0362
\item 此处还可补充更多条件。
\end{enumerate}
\end{lemma}

\begin{proof}
省略。
\end{proof}

\begin{lemma}
\label{formal-spaces-lemma-closed-immersion-into-McQuillan}
设 $S$ 为概形，$X$ 为 $S$ 上的 McQuillan 仿射形式代数空间。
设 $f:Y\to X$ 为 $S$ 上形式代数空间的闭浸入。则 $Y$ 是 McQuillan
仿射形式代数空间，并且 $f$ 对应于弱可容许拓扑 $S$-代数的连续同态
$A\to B$；该同态是 taut 的，具有闭核，并且像稠密。
\end{lemma}

\begin{proof}
写作 $X=\text{Spf}(A)$，其中 $A$ 是弱可容许拓扑环。
令 $I_\lambda$ 为 $A$ 中弱可容许定义理想的基本系统。则
$Y\times_X\Spec(A/I_\lambda)$ 是 $\Spec(A/I_\lambda)$ 的闭子概形，
因而是仿射的（定义 \ref{formal-spaces-definition-closed-immersion}）。设
$Y\times_X\Spec(A/I_\lambda)=\Spec(B_\lambda)$。
$A/I_\lambda\to B_\lambda$ 为满射。因此投影
$$
B = \lim B_\lambda \longrightarrow B_\lambda
$$
都是满射，因为复合 $A\to B\to B_\lambda$ 都是满射。
由引理 \ref{formal-spaces-lemma-mcquillan-affine-formal-algebraic-space}，可知 $Y$ 是
McQuillan 的。由引理 \ref{formal-spaces-lemma-representable-affine}，环映射
$A\to B$ 是 taut 的。由于 $B$ 完备且 $A\to B$ 连续，核是闭的。
最后，由于对所有 $\lambda$，$A\to B_\lambda$ 都是满射，故 $A$ 在
$B$ 中的像稠密。
\end{proof}

\noindent
尽管我们已经得到上述结果，但一般而言，当靶为 McQuillan 仿射形式代数空间时，
我们并不知道闭浸入的行为；见备注 \ref{formal-spaces-remark-questions}。

\begin{example}
\label{formal-spaces-example-closed-immersion-from-quotient}
设 $S$ 为概形，设 $A$ 为 $S$ 上的弱可容许拓扑环。
设 $K\subset A$ 为闭理想。令
$$
B = (A/K)^\wedge = \lim_{I \subset A\ w.i.d.} A/(I + K)
$$
则态射 $\text{Spf}(B)\to\text{Spf}(A)$ 可由代数空间表示；见例
\ref{formal-spaces-example-representable-morphism-from-completion}。
若 $T\to\text{Spf}(A)$ 是态射且 $T$ 是拟紧概形，则由引理
\ref{formal-spaces-lemma-factor-through-thickening}，该态射经由某个弱定义理想
$I\subset A$ 因子分解到 $\Spec(A/I)$。于是
$T\times_{\text{Spf}(A)}\text{Spf}(B)$
等于 $T\times_{\Spec(A/I)}\Spec(A/(K + I))$，
由此可见 $\text{Spf}(B)\to\text{Spf}(A)$ 是闭浸入。
由于 $K$ 闭，$A\to B$ 的核是 $K$；但要注意，一般而言环映射
$A\to B=(A/K)^\wedge$ 未必是满射。
\end{example}

\begin{lemma}
\label{formal-spaces-lemma-monomorphism-iso-over-red}
设 $S$ 为概形，$f:X\to Y$ 为形式代数空间的态射。假设
\begin{enumerate}
\item $f$ 可由代数空间表示；
\item $f$ 是单态射；
\item 包含映射 $Y_{red}\to Y$ 经由 $f$ 因子分解；
\item $f$ 局部有限型，或 $Y$ 局部 Noether。
\end{enumerate}
则 $f$ 是闭浸入。
\end{lemma}

\begin{proof}
假设 (2) 和 (3)，则
$X_{red}=X\times_YY_{red}=Y_{red}$。
以下不再特别说明地使用这一点。

\medskip\noindent
若 $Y'\to Y$ 是 $S$ 上形式代数空间的 'etale 态射，
则基变换 $f':X\times_YY'\to Y'$ 满足条件 (1)--(4)。
因此由引理 \ref{formal-spaces-lemma-characterize-closed-immersion}，可以假设 $Y$ 是
仿射形式代数空间。

\medskip\noindent
按定义 \ref{formal-spaces-definition-affine-formal-algebraic-space}，写作
$Y=\colim_{\lambda\in\Lambda}Y_\lambda$。于是
$X_\lambda=X\times_YY_\lambda$ 是带有单态射
$f_\lambda:X_\lambda\to Y_\lambda$ 的代数空间，并且该态射诱导同构
$X_{\lambda, red}\to Y_{\lambda, red}$。
因此由《空间的极限》之命题
\ref{spaces-limits-proposition-affine}，$X_\lambda$ 是仿射概形
（因为 $X_{\lambda, red}\to X_\lambda$ 是满射且整）。
要完成证明，只需说明 $X_\lambda\to Y_\lambda$ 是闭浸入，
这将在下一段中完成。

\medskip\noindent
设 $X\to Y$ 是仿射概形的单态射，并且
$X_{red}=X\times_YY_{red}=Y_{red}$。一般而言，这并不意味着
$X\to Y$ 是闭浸入；见《例》之第
\ref{examples-section-epimorphism-not-surjective} 节。
然而，在我们的条件 (4) 下，我们知道上一段中
% P10-E0363
$X_\lambda\to Y_\lambda$ 要么是有限型，要么 $Y_\lambda$ 是 Noether。
这意味着 $X\to Y$ 对应于环映射 $R\to A$，其中
$I\subset R$ 是幂零根（即 $R$ 的最大局部幂零理想），且
$R/I\to A/IA$ 是同构；并且 $R\to A$ 是有限型，或 $R$ 是 Noether。
在第一种情形，由《代数》之引理
\ref{algebra-lemma-surjective-mod-locally-nilpotent}，$R\to A$ 是满射；
在第二种情形，$I$ 有限生成，因而是幂零的，所以由 Nakayama 引理，
$R\to A$ 是满射；见《代数》之引理 \ref{algebra-lemma-NAK} 第 (11) 部分。
\end{proof}








\section{受限幂级数}
\label{formal-spaces-section-restricted-power-series}

\noindent
设 $A$ 为关于线性拓扑完备的拓扑环（《代数更多内容》之定义
\ref{more-algebra-definition-topological-ring}）。设 $I_\lambda$ 为开理想的基本系统。
设 $r\geq 0$ 为整数。在此情形下，通常记
$$
A\{x_1, \ldots, x_r\} =
\lim_\lambda A/I_\lambda[x_1, \ldots, x_r] =
\lim_\lambda (A[x_1, \ldots, x_r]/I_\lambda A[x_1, \ldots, x_r])
$$
并赋予极限拓扑。换言之，这是关于理想 $I_\lambda$ 的多项式环完备化。
我们可以把 $A\{x_1, \ldots, x_r\}$ 的元素看作关于
$x_1, \ldots, x_r$ 的幂级数
$$
f = \sum\nolimits_{E = (e_1, \ldots, e_r)} a_E x_1^{e_1} \ldots x_r^{e_r}
$$
其中系数 $a_E\in A$ 在 $A$ 的拓扑中趋于零。换言之，对任意 $\lambda$，
除有限多个 $a_E$ 外，其余都属于 $I_\lambda$。
因此，$A\{x_1, \ldots, x_r\}$ 的元素有时称为\textit{受限幂级数}。
有时将此环记为 $A\langle x_1, \ldots, x_r\rangle$；我们将避免使用这一记号。

\begin{remark}[受限幂级数的泛性质]
\label{formal-spaces-remark-universal-property}
\begin{reference}
\cite[Chapter 0, 7.5.3]{EGA}
\end{reference}
设 $A\to C$ 为完备线性拓扑环之间的连续映射。
则任意 $A$-代数映射 $A[x_1, \ldots x_r]\to C$ 都唯一延拓为
受限幂级数上的连续映射 $A\{x_1, \ldots, x_r\}\to C$。
\end{remark}

\begin{remark}
\label{formal-spaces-remark-I-adic-completion-and-restricted-power-series}
设 $A$ 为环，$I\subset A$ 为理想。若 $A$ 是 $I$-进完备的，
则 $A[x_1, \ldots, x_r]$ 的 $I$-进完备化 $A[x_1, \ldots, x_r]^\wedge$
作为环是 $A$ 上的受限幂级数环。然而，目前尚不清楚
$A[x_1, \ldots, x_r]^\wedge$ 是否 $I$-进完备。
我们把 $A\{x_1, \ldots, x_r\}$ 上的拓扑看作极限拓扑（它总是完备的），
而常常把 $A[x_1, \ldots, x_r]^\wedge$ 上的拓扑看作 $I$-进拓扑
（不总是完备的）。若 $I$ 有限生成，则
$A\{x_1, \ldots, x_r\}=A[x_1, \ldots, x_r]^\wedge$ 作为拓扑环成立；
见《代数》之引理 \ref{algebra-lemma-hathat-finitely-generated}。
\end{remark}






\section{拓扑有限型代数}
\label{formal-spaces-section-tft}

\noindent
这是我们的定义。这个定义并未得到普遍认可。许多作者会附加进一步的条件，
通常是因为他们只关注特定类型的环，而不是最一般的情形。

\begin{definition}
\label{formal-spaces-definition-topologically-finite-type}
设 $A\to B$ 为拓扑环之间的连续映射（《代数更多内容》之定义
\ref{more-algebra-definition-topological-ring}）。若存在像在 $B$ 中稠密的
$A$-代数映射 $A[x_1,\ldots,x_n]\to B$，则称 $B$
\textit{在 $A$ 上拓扑有限型}。
\end{definition}

\noindent
若 $A$ 是完备线性拓扑环，则受限幂级数环
$A\{x_1,\ldots,x_r\}$ 在 $A$ 上拓扑有限型。若 $k$ 是域，则幂级数环
$k[[x_1,\ldots,x_r]]$ 在 $k$ 上拓扑有限型。

\medskip\noindent
对于弱可容许拓扑环之间的连续 taut 映射，拓扑有限型这一性质，恰好对应于
相应仿射形式代数空间之间的有限型态射。

\begin{lemma}
\label{formal-spaces-lemma-topologically-finite-type-finite-type}
设 $S$ 为概形，$\varphi:A\to B$ 为 $S$ 上弱可容许拓扑环之间的连续映射。
下列条件等价：
\begin{enumerate}
\item $\text{Spf}(\varphi):Y=\text{Spf}(B)\to\text{Spf}(A)=X$
是有限型；
\item $\varphi$ 是 taut 的，且 $B$ 在 $A$ 上拓扑有限型。
\end{enumerate}
\end{lemma}

\begin{proof}
可以使用引理 \ref{formal-spaces-lemma-representable-affine} 将 $\varphi$ 的 taut 性与
$\text{Spf}(\varphi)$ 的可表示性联系起来。以下不再特别说明地使用这一点。
于是 $X=\colim\Spec(A/I)$，而
$Y=\colim\Spec(B/J(I))$，其中 $I\subset A$ 遍历 $A$ 的弱定义理想，
$J(I)$ 是 $IB$ 在 $B$ 中的闭包。

\medskip\noindent
假设 (2)。选取像在 $B$ 中稠密的环映射
$A[x_1,\ldots,x_r]\to B$。于是
$A[x_1,\ldots,x_r]\to B\to B/J(I)$ 的像也稠密，
这意味着它是满射。因此 $B/J(I)$ 是 $A/I$ 上的有限型代数。
设 $T\to X$ 是态射且 $T$ 是拟紧概形。则由引理
\ref{formal-spaces-lemma-factor-through-thickening}，对某个 $I$，$T\to X$ 经由
$\Spec(A/I)$ 因子分解。于是
$T\times_XY=T\times_{\Spec(A/I)}\Spec(B/J(I))$；见引理
\ref{formal-spaces-lemma-representable-affine} 的证明。
% P10-E0364
因此 $T\times_YX\to T$ 是有限型，因为它是有限型态射
$\Spec(B/J(I))\to\Spec(A/I)$ 的基变换。所以 (1) 成立。

\medskip\noindent
假设 (1)。任取如上的 $I\subset A$。由于
$\Spec(A/I)\times_XY=\Spec(B/J(I))$，可知
$A/I\to B/J(I)$ 是有限型的。选取 $b_1,\ldots,b_r\in B$，
它们在 $A/I$ 上映到 $B/J(I)$ 的生成元。我们断言，将 $x_i$ 映到 $b_i$ 的
环映射 $A[x_1,\ldots,x_r]\to B$ 的像是稠密的。为证明这一点，
令 $I'\subset I$ 为另一个弱定义理想。则
$$
B/(J(I') + IB) = B/J(I)
$$
因为 $J(I)$ 是 $IB$ 的闭包，而 $J(I')$ 是开集。于是可以应用《代数》之引理
\ref{algebra-lemma-surjective-mod-locally-nilpotent}，得到
$(A/I')[x_1,\ldots,x_r]\to B/J(I')$ 是满射。
故 (2) 成立，证明完成。
\end{proof}

\noindent
设 $A$ 为关于线性拓扑完备的拓扑环，设 $(I_\lambda)$ 为开理想的基本系统。
令 $\mathcal{C}$ 为逆系统 $(B_\lambda)$ 所成的范畴，其中
\begin{enumerate}
\item $B_\lambda$ 是有限型 $A/I_\lambda$-代数，并且
\item $B_\mu\to B_\lambda$ 是 $A/I_\mu$-代数同态，且诱导同构
$B_\mu/I_\lambda B_\mu\to B_\lambda$。
\end{enumerate}
$\mathcal{C}$ 中的态射由相容的同态系统给出。

\begin{lemma}
\label{formal-spaces-lemma-category-affine-over}
设 $S$ 为概形，$X$ 为 $S$ 上的仿射形式代数空间。假设 $X$ 是 McQuillan，
并令 $A$ 为与 $X$ 关联的弱可容许拓扑环。则在以下两个范畴之间存在反等价：
\begin{enumerate}
\item 上面引入的范畴 $\mathcal{C}$；
\item $X$ 上有限型仿射形式代数空间的态射 $Y\to X$ 所成的范畴。
\end{enumerate}
\end{lemma}

\begin{proof}
令 $(I_\lambda)$ 为 $A$ 中弱可容许定义理想的基本系统。考虑 (2) 中的 $Y$。
则 $Y\times_X\Spec(A/I_\lambda)$ 是仿射的（定义
\ref{formal-spaces-definition-finite-type} 及引理 \ref{formal-spaces-lemma-affine-representable-by-algebraic-spaces}）。
设 $Y\times_X\Spec(A/I_\lambda)=\Spec(B_\lambda)$。
由于 $\Spec(B_\lambda)\to\Spec(A/I_\lambda)$ 是有限型（由定义
\ref{formal-spaces-definition-finite-type}），环映射 $A/I_\lambda\to B_\lambda$ 是有限型的。
于是 $(B_\lambda)$ 是 $\mathcal{C}$ 的对象。

\medskip\noindent
反之，给定 $\mathcal{C}$ 的对象 $(B_\lambda)$，令
$Y=\colim\Spec(B_\lambda)$。这是仿射形式代数空间。我们断言
$$
Y \times_X \Spec(A/I_\lambda) =
\left(\colim_\mu \Spec(B_\mu)\right) \times_X \Spec(A/I_\lambda) =
\Spec(B_\lambda)
$$
为证明这一点，只需在拟紧概形 $U$ 上取值时得到相同结果。
态射
$U\to\left(\colim_\mu\Spec(B_\mu)\right)\times_X\Spec(A/I_\lambda)$
来自某个 $\mu\geq\lambda$ 的态射
$U\to\Spec(B_\mu)\times_{\Spec(A/I_\mu)}\Spec(A/I_\lambda)$
（两次使用引理 \ref{formal-spaces-lemma-factor-through-thickening}）。由于
$\mathcal{C}$ 的对象所满足的第二个条件，
$\Spec(B_\mu)\times_{\Spec(A/I_\mu)}\Spec(A/I_\lambda)=\Spec(B_\lambda)$，
这就证明了所需结论。由此可以说明 $Y\to X$ 是有限型。确切地说，对于任意
$U\to X$ 且 $U$ 为拟紧概形，由上面引用的引理，对某个 $\lambda$ 有分解
$U\to\Spec(A/I_\lambda)\to X$。于是
$$
Y \times_X U =
Y \times_X \Spec(A/I_\lambda)) \times_{\Spec(A/I_\lambda)} U =
\Spec(B_\lambda) \times_{\Spec(A/I_\lambda)} U
$$
% P10-E0365
如所需，这是 $U$ 上的有限型概形。因此构造
$(B_\lambda)\mapsto\colim\Spec(B_\lambda)$ 确实给出从范畴 (1) 到范畴 (2) 的函子。

\medskip\noindent
为完成证明，我们说明上述构造在范畴 (1) 和 (2) 之间定义了准逆函子。
一个方向需要证明，对于范畴 $\mathcal{C}$ 中任意对象 $(B_\lambda)$，有
$$
\left(\colim_\mu \Spec(B_\mu)\right) \times_X \Spec(A/I_\lambda) =
\Spec(B_\lambda)
$$
上面已经证明了这一点。另一个方向需要证明，对范畴 (2) 中给定的 $Y$，有
$$
Y = \colim (Y \times_X \Spec(A/I_\lambda))
$$
这同样可通过在拟紧测试对象上取值，并利用
$X=\colim\Spec(A/I_\lambda)$ 得到。
\end{proof}

\begin{remark}
\label{formal-spaces-remark-questions}
设 $A$ 为弱可容许拓扑环，$(I_\lambda)$ 为弱定义理想的基本系统。
令 $X=\text{Spf}(A)$，换言之，$X$ 是 McQuillan 仿射形式代数空间。
设 $f:Y\to X$ 为仿射形式代数空间的态射。一般而言，$Y$ 未必是 McQuillan。
更具体地，我们可以提出下列问题：
\begin{enumerate}
\item 假设 $f:Y\to X$ 是闭浸入。则 $Y$ 是 McQuillan，且 $f$ 对应于弱可容许拓扑环的连续映射
$\varphi:A\to B$，它是 taut 的，核 $K\subset A$ 是闭理想，并且像
$\varphi(A)$ 在 $B$ 中稠密；见引理 \ref{formal-spaces-lemma-closed-immersion-into-McQuillan}。
什么条件能保证 $B=(A/K)^\wedge$，如例 \ref{formal-spaces-example-closed-immersion-from-quotient} 中那样？
\item 什么条件能保证闭浸入 $f:Y\to X$ 对应于 $A$ 按闭理想作商所得的商环 $A/K$；换言之，
相应的连续映射 $\varphi$ 是满射且开映射？
\item 假设 $f:Y\to X$ 是有限型。由引理 \ref{formal-spaces-lemma-category-affine-over}，有
$Y=\colim\Spec(B_\lambda)$，其中 $(B_\lambda)$ 是 $\mathcal{C}$ 的对象。
此时确实存在固定整数 $r$，使得对所有 $\lambda$，$B_\lambda$ 都由
$A/I_\lambda$ 上的 $r$ 个元素生成（论证实质上已在引理
\ref{formal-spaces-lemma-topologically-finite-type-finite-type} 的 (1) $\Rightarrow$ (2) 证明中给出）。
然而，目前尚不清楚投影 $\lim B_\lambda\to B_\lambda$ 是否为满射，即尚不清楚 $Y$ 是否为 McQuillan。
是否存在 $Y$ 不是 McQuillan 的例子？
\item 假设 $f:Y\to X$ 是有限型且 $Y$ 是 McQuillan。则 $f$ 对应于弱可容许拓扑环的连续映射
$\varphi:A\to B$。事实上，$\varphi$ 是 taut 的，且 $B$ 在 $A$ 上拓扑有限型；见引理
\ref{formal-spaces-lemma-topologically-finite-type-finite-type}。换言之，$f$ 可分解为
$$
Y \longrightarrow \mathbf{A}^r_X \longrightarrow X
$$
其中第一个箭头是 McQuillan 仿射形式代数空间的闭浸入。然而，此时对
$Y\to\mathbf{A}^r_X$，问题 (1) 和 (2) 仍然有效。
\end{enumerate}
下面我们将在 $X$ 可数指标化时回答这些问题，即 $A$ 具有可数的开理想基本系统时。
如果你在更一般的情形下有这些问题的答案，或者有反例，请发送电子邮件至
\href{mailto:stacks.project@gmail.com}{stacks.project@gmail.com}。
\end{remark}

\begin{lemma}
\label{formal-spaces-lemma-closed-immersion-into-countably-indexed}
设 $S$ 为概形，$X$ 为 $S$ 上可数指标化的仿射形式代数空间。
设 $f:Y\to X$ 为 $S$ 上形式代数空间的闭浸入。则 $Y$ 是可数指标化的仿射形式代数空间，
且 $f$ 对应于 $A\to A/K$，其中 $A$ 是
$\textit{WAdm}^{count}$ 中的对象（第 \ref{formal-spaces-section-morphisms-rings} 节），
而 $K\subset A$ 是闭理想。
\end{lemma}

\begin{proof}
由引理 \ref{formal-spaces-lemma-countably-indexed}，$X=\text{Spf}(A)$，其中 $A$ 是
$\textit{WAdm}^{count}$ 中的对象。由于闭浸入是可表示且仿射的，
由引理 \ref{formal-spaces-lemma-property-goes-up-affine-morphism} 可知 $Y$ 是仿射形式代数空间且可数
% P10-E0366
指标化。因此再次应用引理 \ref{formal-spaces-lemma-countably-indexed}，得到
$Y=\text{Spf}(B)$，其中 $B$ 是 $\textit{WAdm}^{count}$ 中的对象。
由引理 \ref{formal-spaces-lemma-closed-immersion-into-McQuillan}，可知 $f$ 由
$\textit{WAdm}^{count}$ 中一个 taut 且像稠密的态射 $A\to B$ 给出。
最后应用引理 \ref{formal-spaces-lemma-dense-image-surjective} 即可完成证明。
\end{proof}

\begin{lemma}
\label{formal-spaces-lemma-quotient-restricted-power-series}
设 $B\to A$ 为 $\textit{WAdm}^{count}$ 中的箭头，见第
\ref{formal-spaces-section-morphisms-rings} 节。下列条件等价：
\begin{enumerate}
\item[(a)] $B\to A$ 是 taut 的，并且对每个弱定义理想 $J\subset B$，
其中 $I\subset A$ 是 $JA$ 的闭包，$B/J\to A/I$ 是有限型；
\item[(b)] $B\to A$ 是 taut 的，并且对 $B$ 的弱定义理想的一个共尾系统
$(J_\lambda)$，$B/J_\lambda\to A/I_\lambda$ 是有限型，其中
$I_\lambda\subset A$ 是 $J_\lambda A$ 的闭包；
\item[(c)] $B\to A$ 是 taut 的，并且 $A$ 在 $B$ 上拓扑有限型；
\item[(d)] 作为拓扑 $B$-代数，$A$ 同构于
$B\{x_1, \ldots, x_n\}$ 对某个闭理想取商所得的代数。
\end{enumerate}
此外，这些等价条件定义了一个局部性质，即它们满足公理
\ref{formal-spaces-item-axiom-1}、\ref{formal-spaces-item-axiom-2}、\ref{formal-spaces-item-axiom-3}。
\end{lemma}

\begin{proof}
由定义，蕴含 (a) $\Rightarrow$ (b)、(c) $\Rightarrow$ (a) 以及
(d) $\Rightarrow$ (c) 都是直接的。假设 (b) 成立，令 $J\subset B$ 和
$I\subset A$ 如 (a) 中所述。取交换图表
$$
\xymatrix{
A \ar[r] & \ldots \ar[r] & A_3 \ar[r] & A_2 \ar[r] & A_1 \\
B \ar[r] \ar[u] & \ldots \ar[r] & B/J_3 \ar[r] \ar[u] &
B/J_2 \ar[r] \ar[u] & B/J_1 \ar[u]
}
$$
使 $A_{n + 1}/J_nA_{n + 1}=A_n$，并且如引理
\ref{formal-spaces-lemma-representable-property-rings} 中那样 $A=\lim A_n$。
对每个 $m$，存在 $\lambda$ 使 $J_\lambda\subset J_m$。
由于 $B/J_\lambda\to A/I_\lambda$ 是有限型，这说明
$B/J_m\to A/I_m$ 是有限型。令 $\alpha_1, \ldots, \alpha_n\in A_1$
为 $A_1$ 在 $B/J_1$ 上的生成元。由于 $A$ 是一个具有满过渡映射的系统的
可数极限，可以找到 $a_1, \ldots, a_n\in A$，它们在 $A_1$ 中分别映到
$\alpha_1, \ldots, \alpha_n$。由备注 \ref{formal-spaces-remark-universal-property}，
存在连续映射 $B\{x_1, \ldots, x_n\}\to A$，将 $x_i$ 映到 $a_i$。
由《代数》之引理 \ref{algebra-lemma-surjective-mod-locally-nilpotent}，
此映射诱导满射 $(B/J_m)[x_1, \ldots, x_n]\to A_m$。
对 $m\geq 1$，得到短正合序列
$$
0 \to K_m \to (B/J_m)[x_1, \ldots, x_n] \to A_m \to 0
$$
诱导的过渡映射 $K_{m + 1}\to K_m$ 是满射，因为
$A_{m + 1}/J_mA_{m + 1}=A_m$。因此这些短正合序列的逆极限是正合的；见《代数》之引理
\ref{algebra-lemma-ML-exact-sequence}。由于
$B\{x_1, \ldots, x_n\}=\lim(B/J_m)[x_1, \ldots, x_n]$
且 $A=\lim A_m$，可知 $B\{x_1, \ldots, x_n\}\to A$ 是满射且开映射。
由于 $A$ 完备，核是闭理想。由此得到 (a)、(b)、(c)、(d) 等价。

\medskip\noindent
给定图表 \ref{formal-spaces-equation-localize}，其情形如
\ref{formal-spaces-situation-local-property} 所述。由例 \ref{formal-spaces-example-finite-type-from-finite-type-ring-map}，
$A\to(A')^\wedge$ 和 $B\to(B')^\wedge$ 满足 (a)、(b)、(c)、(d)。此外，
由引理 \ref{formal-spaces-lemma-representable-property-rings}，为证明公理
\ref{formal-spaces-item-axiom-1} 和 \ref{formal-spaces-item-axiom-2}，可以假设 $B\to A$ 与
$(B')^\wedge\to(A')^\wedge$ 都是 taut 的。现在取 $B$ 中的弱定义理想
$J$。令 $J'\subset(B')^\wedge$、$I\subset A$、$I'\subset(A')^\wedge$
分别为 $J(B')^\wedge$、$JA$、$J(A')^\wedge$ 的闭包。由上面的结论，
只需考虑交换图表
$$
\xymatrix{
A/I \ar[r] & (A')^\wedge/I' \\
B/J \ar[r] \ar[u]^{\overline{\varphi}} &
(B')^\wedge/J' \ar[u]_{\overline{\varphi}'}
}
$$
并证明：(1) $\overline{\varphi}$ 有限型蕴含 $\overline{\varphi}'$
有限型；(2) 若 $A\to A'$ 忠实平坦，则 $\overline{\varphi}'$ 有限型蕴含
$\overline{\varphi}$ 有限型。注意，由 $(B')^\wedge$ 和 $(A')^\wedge$ 的拓扑构造，
$(B')^\wedge/J'=B'/JB'$ 且 $(A')^\wedge/I'=A'/IA'$。
特别地，图表中的水平映射是 \'etale 的。(1) 现在由《代数》之引理
\ref{algebra-lemma-compose-finite-type} 得出；(2) 则由《下降》之引理
\ref{descent-lemma-finite-type-local-source-fppf-algebra} 得出，因为环映射
$A/I\to(A')^\wedge/I'=A'/IA'$ 是忠实平坦且 \'etale 的。

\medskip\noindent
公理 \ref{formal-spaces-item-axiom-3} 的证明略去。
\end{proof}

\begin{lemma}
\label{formal-spaces-lemma-quotient-restricted-power-series-admissible}
在引理 \ref{formal-spaces-lemma-quotient-restricted-power-series} 中，若 $B$ 是可容许的
（例如 adic），则等价条件 (a)--(d) 还等价于
\begin{enumerate}
\item[(e)] $B\to A$ 是 taut 的，并且对某个定义理想 $J\subset B$，
其中 $I\subset A$ 是 $JA$ 的闭包，$B/J\to A/I$ 是有限型。
\end{enumerate}
\end{lemma}

\begin{proof}
只需证明 (e) 蕴含 (a)。设 $J'\subset B$ 为弱定义理想，令
$I'\subset A$ 为 $J'A$ 的闭包。我们需要证明 $B/J'\to A/I'$ 是有限型。
若相应结论对较小的弱定义理想 $J''=J'\cap J$ 成立，则对 $J'$ 也成立。
因此可假设 $J'\subset J$。由于 $J$ 是定义理想（而不仅是弱定义理想），
存在 $n\geq 1$ 使 $J^n\subset J'$。于是可以考虑图表
$$
\xymatrix{
0 \ar[r] & I/I' \ar[r] & A/I' \ar[r] & A/I \ar[r] & 0 \\
0 \ar[r] & J/J' \ar[r] \ar[u] & B/J' \ar[r] \ar[u] & B/J \ar[r] \ar[u] & 0
}
$$
其两行都是正合的。由于 $I'\subset A$ 是开集，并且 $I$ 是 $JA$ 的闭包，
可知 $I/I'=(J/J')\cdot A/I'$。由于 $J/J'$ 是幂零理想，且
$B/J\to A/I$ 是有限型，由《代数》之引理
\ref{algebra-lemma-finite-type-mod-nilpotent}，可知 $A/I'$ 在 $B/J'$ 上有限型，
这正是所需结论。
\end{proof}

\begin{lemma}
\label{formal-spaces-lemma-representable-affine-finite-type}
设 $S$ 为概形，$f:X\to Y$ 为仿射形式代数空间的态射。假设 $Y$ 可数指标化。
下列条件等价：
\begin{enumerate}
\item $f$ 局部有限型；
\item $f$ 有限型；
\item $f$ 对应于 $\textit{WAdm}^{count}$ 中的态射 $B\to A$
（第 \ref{formal-spaces-section-morphisms-rings} 节），并满足引理
\ref{formal-spaces-lemma-quotient-restricted-power-series} 的等价条件。
\end{enumerate}
\end{lemma}

\begin{proof}
由于 $X$ 和 $Y$ 是仿射的，条件 (1) 和 (2) 等价是显然的。在 (1) 或 (2) 的情形，
按定义 $f$ 可由代数空间表示，因而由引理
\ref{formal-spaces-lemma-affine-representable-by-algebraic-spaces} 可知它是仿射的。
因此若 (1) 或 (2) 成立，由引理 \ref{formal-spaces-lemma-property-goes-up-affine-morphism}，
可知 $X$ 是可数指标化的。按引理 \ref{formal-spaces-lemma-countably-indexed}，写作
$X=\text{Spf}(A)$、$Y=\text{Spf}(B)$，其中 $A$ 和 $B$ 是
$\textit{WAdm}^{count}$ 中的拓扑 $S$-代数。由引理
\ref{formal-spaces-lemma-morphism-between-formal-spectra}，可知 $f$ 对应于连续映射
$B\to A$。于是结论由引理
\ref{formal-spaces-lemma-topologically-finite-type-finite-type} 得出。
\end{proof}

\begin{lemma}
\label{formal-spaces-lemma-finite-type-local-property}
设 $S$ 为概形，$f:X\to Y$ 为 $S$ 上局部可数指标化形式代数空间的态射。
下列条件等价：
\begin{enumerate}
\item 对每个交换图表
$$
\xymatrix{
U \ar[d] \ar[r] & V \ar[d] \\
X \ar[r] & Y
}
$$
若 $U,V$ 是仿射形式代数空间，$U\to X$ 和 $V\to Y$ 可由代数空间表示且为
\'etale，则态射 $U\to V$ 对应于 $\textit{WAdm}^{count}$ 中的态射，
该态射是 taut 且拓扑有限型；
\item 存在定义 \ref{formal-spaces-definition-formal-algebraic-space} 意义下的覆盖
$\{Y_j\to Y\}$，且对每个 $j$，存在定义
\ref{formal-spaces-definition-formal-algebraic-space} 意义下的覆盖
$\{X_{ji}\to Y_j\times_Y X\}$，使每个 $X_{ji}\to Y_j$ 对应于
$\textit{WAdm}^{count}$ 中的态射，该态射是 taut 且拓扑有限型；
\item 存在定义 \ref{formal-spaces-definition-formal-algebraic-space} 意义下的覆盖
$\{X_i\to X\}$，并且对每个 $i$ 存在分解 $X_i\to Y_i\to Y$，其中 $Y_i$
是仿射形式代数空间，$Y_i\to Y$ 可由代数空间表示且为 \'etale，且
$X_i\to Y_i$ 对应于 $\textit{WAdm}^{count}$ 中的态射，
% P10-E0368
该态射是 taut 且拓扑有限型，并且
\item $f$ 局部有限型。
\end{enumerate}
\end{lemma}

\begin{proof}
由引理 \ref{formal-spaces-lemma-quotient-restricted-power-series}，性质
$P(\varphi)=$“$\varphi$ 是 taut 且拓扑有限型”在
$\text{WAdm}^{count}$ 上是局部的。因此由引理
\ref{formal-spaces-lemma-property-defines-property-morphisms}，条件 (1)、(2)、(3) 等价。
另一方面，由引理 \ref{formal-spaces-lemma-representable-affine-finite-type}，
$\textit{WAdm}^{count}$ 中态射的性质 $P$ 恰好对应于可数指标化仿射形式代数空间
态射的局部有限型性质。因此引理 \ref{formal-spaces-lemma-finite-type-local} 的蕴含
(1) $\Rightarrow$ (3) 表明，本引理的 (4) 蕴含 (1)。类似地，
引理 \ref{formal-spaces-lemma-finite-type-local} 的蕴含 (4) $\Rightarrow$ (1) 表明，
本引理的 (2) 蕴含 (4)。
\end{proof}






\section{态射的分离公理}
\label{formal-spaces-section-separation-axioms}

\noindent
本节是《空间的态射》第 \ref{spaces-morphisms-section-separation-axioms} 节
关于形式代数空间态射的对应版本。

\begin{definition}
\label{formal-spaces-definition-separated-morphism}
设 $S$ 为概形。设 $f:X\to Y$ 为 $S$ 上形式代数空间的态射。
设 $\Delta_{X/Y}:X\to X\times_YX$ 为对角态射。
\begin{enumerate}
\item 若 $\Delta_{X/Y}$ 是闭浸入，则称 $f$ 为\textit{分离的}；
\item 若 $\Delta_{X/Y}$ 是拟紧的，则称 $f$ 为\textit{拟分离的}。
\end{enumerate}
\end{definition}

\noindent
由于由引理 \ref{formal-spaces-lemma-diagonal-morphism-formal-algebraic-spaces}，
$\Delta_{X/Y}$ 可表示（由概形表示），所以可以考虑从仿射概形 $T$ 到
$X\times_YX$ 的态射，并检查
$$
E = T \times_{X \times_Y X} X \longrightarrow T
$$
是否拟紧或为闭浸入；见引理 \ref{formal-spaces-lemma-quasi-compact-representable} 或定义
\ref{formal-spaces-definition-closed-immersion}。
注意，概形 $E$ 是两个态射 $a,b:T\to X$ 的等化子，其中它们作为到
$Y$ 的态射相同；并且 $E\to T$ 是单态射且局部有限型。

\begin{lemma}
\label{formal-spaces-lemma-base-change-separated}
定义 \ref{formal-spaces-definition-separated-morphism} 中列出的所有分离公理在基变换下稳定。
\end{lemma}

\begin{proof}
设 $f:X\to Y$ 和 $Y'\to Y$ 为形式代数空间的态射。设
$f':X'\to Y'$ 是 $f$ 沿 $Y'\to Y$ 的基变换。则
$\Delta_{X'/Y'}$ 是 $\Delta_{X/Y}$ 沿态射
$X'\times_{Y'}X'\to X\times_YX$ 的基变换。
定义 \ref{formal-spaces-definition-separated-morphism} 中使用的对角态射的每个性质
都在基变换下稳定。因此引理成立。
\end{proof}

\begin{lemma}
\label{formal-spaces-lemma-fibre-product-after-map}
设 $S$ 为概形，$f:X\to Z$、$g:Y\to Z$ 以及 $Z\to T$ 为 $S$ 上形式代数空间的态射。
考虑诱导态射 $i:X\times_ZY\to X\times_TY$。则
\begin{enumerate}
\item $i$ 可表示（由概形表示），局部有限型、局部拟有限、分离且为单态射；
\item 若 $Z\to T$ 分离，则 $i$ 是闭浸入，并且
\item 若 $Z\to T$ 拟分离，则 $i$ 拟紧。
\end{enumerate}
\end{lemma}

\begin{proof}
由一般范畴论，下面的图表
% P10-E0369
$$
\xymatrix{
X \times_Z Y \ar[r]_i \ar[d] & X \times_T Y \ar[d] \\
Z \ar[r]^-{\Delta_{Z/T}} \ar[r] & Z \times_T Z
}
$$
是纤维积图表。因此 $i$ 是对角态射 $\Delta_{Z/T}$ 的基变换。
于是引理由引理 \ref{formal-spaces-lemma-diagonal-morphism-formal-algebraic-spaces} 得出。
\end{proof}

\begin{lemma}
\label{formal-spaces-lemma-composition-separated}
定义 \ref{formal-spaces-definition-separated-morphism} 中列出的所有分离公理在态射复合下稳定。
\end{lemma}

\begin{proof}
设 $f:X\to Y$ 和 $g:Y\to Z$ 为适用所讨论公理的形式代数空间态射。
对角态射 $\Delta_{X/Z}$ 是复合
$$
X \longrightarrow X \times_Y X \longrightarrow X \times_Z X
$$
我们的分离公理通过要求对角态射具有某个性质 $\mathcal{P}$ 来定义。
由上面的引理 \ref{formal-spaces-lemma-fibre-product-after-map}，第二个箭头也具有这一性质。
因此，由于具有性质 $\mathcal{P}$ 的（可表示）态射的复合仍是具有性质
$\mathcal{P}$ 的态射，引理得证。
\end{proof}

\begin{lemma}
\label{formal-spaces-lemma-separated-local}
设 $S$ 为概形，$f:X\to Y$ 为 $S$ 上形式代数空间的态射。令 $\mathcal{P}$ 为定义
\ref{formal-spaces-definition-separated-morphism} 中的任一分离公理。下列条件等价：
\begin{enumerate}
\item $f$ 具有性质 $\mathcal{P}$；
\item 对每个概形 $Z$ 以及态射 $Z\to Y$，$f$ 的基变换
$Z\times_YX\to Z$ 具有性质 $\mathcal{P}$；
\item 对每个仿射概形 $Z$ 以及态射 $Z\to Y$，$f$ 的基变换
$Z\times_YX\to Z$ 具有性质 $\mathcal{P}$；
\item 对每个仿射概形 $Z$ 以及态射 $Z\to Y$，形式代数空间
$Z\times_YX$ 具有性质 $\mathcal{P}$（见定义 \ref{formal-spaces-definition-separated}）；
\item 存在定义 \ref{formal-spaces-definition-formal-algebraic-space} 意义下的覆盖
$\{Y_j\to Y\}$，使得对所有 $j$，基变换 $Y_j\times_YX\to Y_j$ 具有性质 $\mathcal{P}$。
\end{enumerate}
\end{lemma}

\begin{proof}
以下将反复使用引理 \ref{formal-spaces-lemma-base-change-separated}，不再逐一说明。
特别地，(1) 蕴含 (2)，(2) 蕴含 (3) 是显然的。

\medskip\noindent
假设 (3)，设 $Z\to Y$ 为态射，其中 $Z$ 是仿射概形。令 $U,V$ 为仿射概形，
并设 $a:U\to Z\times_YX$ 和 $b:V\to Z\times_YX$ 为态射。则
$$
U \times_{Z \times_Y X} V =
(Z \times_Y X) \times_{\Delta, (Z \times_Y X) \times_Z (Z \times_Y X)}
(U \times_Z V)
$$
当 $\mathcal{P}=$“拟分离”时，这说明该空间是拟紧的；当
$\mathcal{P}=$“分离”时，这说明它是带有到 $U\times_ZV$ 的闭浸入的仿射概形。
因此 (4) 成立。

\medskip\noindent
假设 (4)，设 $Z\to Y$ 为态射，其中 $Z$ 是仿射概形。令 $U,V$ 为仿射概形，
并设 $a:U\to Z\times_YX$ 和 $b:V\to Z\times_YX$ 为态射。
将上面的论证反向进行，可知 $U\times_{Z\times_YX}V\to U\times_ZV$
在 $\mathcal{P}=$“拟分离”时拟紧，在 $\mathcal{P}=$“分离”时为闭浸入。
由于可以选择上述 $U,V$，使 $U$ 遍历 $Z\times_YX$ 的一个
\'etale 覆盖，我们得到相应态射
$$
U \times_Z V \to (Z \times_Y X) \times_Z (Z \times_Y X)
$$
构成由仿射概形组成的 \'etale 覆盖。因此
$\Delta:(Z\times_YX)\to(Z\times_YX)\times_Z(Z\times_YX)$
拟紧，或为闭浸入（分别对应两种 $\mathcal{P}$）。于是 (3) 成立。

\medskip\noindent
证明 (3) 蕴含 (5)。假设 (3)，令 $\{Y_j\to Y\}$ 如定义
\ref{formal-spaces-definition-formal-algebraic-space} 所述。我们需要证明态射
$$
\Delta_j :
Y_j \times_Y X
\longrightarrow
(Y_j \times_Y X) \times_{Y_j} (Y_j \times_Y X) =
Y_j \times_Y X \times_Y X
$$
具有相应性质（即拟紧或为闭浸入）。按定义
\ref{formal-spaces-definition-affine-formal-algebraic-space}，写作
$Y_j=\colim Y_{j,\lambda}$。在上式中以 $Y_{j,\lambda}$ 替代 $Y_j$，
由假设 (3) 可得所需性质。由于显示的箭头是箭头 $\Delta_{j,\lambda}$ 的余极限，
并且可以通过对映入 $Y_j\times_YX\times_YX$ 的仿射概形作基变换来检验
$\Delta_j$ 是否具有相应性质，故由引理
\ref{formal-spaces-lemma-factor-through-thickening} 得出结论。

\medskip\noindent
证明 (5) 蕴含 (1)。令 $\{Y_j\to Y\}$ 如 (5) 所述。则有纤维积图表
$$
\xymatrix{
\coprod Y_j \times_Y X \ar[r] \ar[d] &
X \ar[d] \\
\coprod Y_j \times_Y X \times_Y X \ar[r] &
X \times_Y X
}
$$
由假设，左侧竖直箭头是拟紧或闭浸入。由《空间》之引理
\ref{spaces-lemma-descent-representable-transformations-property}，
右侧竖直箭头也拟紧或为闭浸入。
\end{proof}






\section{固有态射}
\label{formal-spaces-section-proper}

\noindent
下面是我们将使用的定义。

\begin{definition}
\label{formal-spaces-definition-proper}
设 $S$ 为概形，设 $f:Y\to X$ 为 $S$ 上形式代数空间的态射。
若 $f$ 可由代数空间表示，并且在《自举》之定义
\ref{bootstrap-definition-property-transformation} 的意义下是固有的，
则称 $f$ 为\textit{固有态射}。
\end{definition}

\noindent
由定义可知，固有态射是有限型的。

\begin{lemma}
\label{formal-spaces-lemma-proper-local}
设 $S$ 为概形，$f:X\to Y$ 为 $S$ 上形式代数空间的态射。
下列条件等价：
\begin{enumerate}
\item $f$ 是固有态射；
\item 对每个概形 $Z$ 以及态射 $Z\to Y$，$f$ 的基变换
$Z\times_YX\to Z$ 是固有态射；
\item 对每个仿射概形 $Z$ 以及每个态射 $Z\to Y$，$f$ 的基变换
$Z\times_YX\to Z$ 是固有态射；
\item 对每个仿射概形 $Z$ 以及每个态射 $Z\to Y$，形式代数空间
$Z\times_YX$ 是在 $Z$ 上固有的代数空间；
\item 存在定义 \ref{formal-spaces-definition-formal-algebraic-space} 意义下的覆盖
$\{Y_j\to Y\}$，使得对所有 $j$，基变换 $Y_j\times_YX\to Y_j$ 是固有的。
\end{enumerate}
\end{lemma}

\begin{proof}
省略。
\end{proof}

\begin{lemma}
\label{formal-spaces-lemma-base-change-proper}
形式代数空间的固有态射在基变换下保持固有性。
\end{lemma}

\begin{proof}
这是引理 \ref{formal-spaces-lemma-proper-local} 及基变换的传递性的直接推论。
\end{proof}






\section{形式代数空间与 fpqc 覆盖}
\label{formal-spaces-section-fpqc}

\noindent
本节是《空间的性质》第 \ref{spaces-properties-section-fpqc} 节关于形式代数空间的对应版本。
请先阅读该节。

\begin{lemma}
\label{formal-spaces-lemma-sheaf-fpqc}
\begin{slogan}
形式代数空间是 fpqc 层
\end{slogan}
设 $S$ 为概形，设 $X$ 为 $S$ 上的形式代数空间。则 $X$ 满足 fpqc 拓扑的层性质。
\end{lemma}

\begin{proof}
证明与《空间的性质》之命题
\ref{spaces-properties-proposition-sheaf-fpqc} 的证明\textbf{完全相同}。
由于 $X$ 是 Zariski 拓扑下的层，只需证明如下事实。给定仿射概形之间的满平坦态射
$f:T'\to T$，有：$X(T)$ 是两个映射
$X(T')\to X(T'\times_TT')$ 的等化子。见《拓扑》之引理
\ref{topologies-lemma-sheaf-property-fpqc}。

\medskip\noindent
设 $a,b:T\to X$ 为两个满足 $a\circ f=b\circ f$ 的态射。我们需要证明 $a=b$。
考虑纤维积
$$
E = X \times_{\Delta_{X/S}, X \times_S X, (a, b)} T.
$$
由引理 \ref{formal-spaces-lemma-diagonal-formal-algebraic-space}，对角态射
$\Delta_{X/S}$ 是可表示的单态射。因此 $E\to T$ 是概形的单态射。
由假设 $a\circ f=b\circ f$，$T'\to T$ 唯一地经由 $E$ 因子分解。
考虑交换图表
$$
\xymatrix{
T' \times_T E \ar[r] \ar[d] & E \ar[d] \\
T' \ar[r] \ar@/^5ex/[u] \ar[ru] & T
}
$$
由于投影 $T'\times_TE\to T'$ 是带截面的单态射，可知它是同构。
于是由《下降》之引理
\ref{descent-lemma-descending-property-isomorphism}，$E\to T$ 是同构。
这就得到所需的 $a=b$。

\medskip\noindent
现在设 $c:T'\to X$ 为态射，且两个复合
$T'\times_TT'\to T'\to X$ 相同。我们需要找到态射 $a:T\to X$，
使其与 $T'\to T$ 的复合为 $c$。选取形式仿射概形 $U$ 及一个 \'etale 态射
$U\to X$，使 $|U|\to|X_{red}|$ 的像包含
$|c|:|T'|\to|X_{red}|$ 的像。这由定义
\ref{formal-spaces-definition-formal-algebraic-space}、《空间的性质》之引理
\ref{spaces-properties-lemma-topology-points}、有限个形式仿射代数空间的并仍为形式仿射代数空间
这一事实以及 $|T'|$ 拟紧这一事实得到（略去小论证）。
$U\to X$ 可由概形表示（引理 \ref{formal-spaces-lemma-presentation-representable}）且分离
（引理 \ref{formal-spaces-lemma-separated-from-separated}）。因此
$$
V = U \times_{X, c} T' \longrightarrow T'
$$
是概形的 \'etale 且分离的态射。由我们对 $U\to X$ 的选择，它也是满射
（若不想论证这一点，也可将 $U$ 换成形式仿射代数空间的互不相交并，使
$U\to X$ 满射；其余论证仍然成立）。条件
$c\circ\text{pr}_0=c\circ\text{pr}_1$ 意味着在 $V/T'/T$ 上得到下降数据
（《下降》之定义 \ref{descent-definition-descent-datum}），因为
\begin{align*}
V \times_{T'} (T' \times_T T')
& = U \times_{X, c \circ \text{pr}_0} (T' \times_T T') \\
& = (T' \times_T T') \times_{c \circ \text{pr}_1, X} U \\
& = (T' \times_T T') \times_{T'} V
\end{align*}
态射 $V\to T'$ 是 ind-拟仿射的，由《态射更多内容》之引理
\ref{more-morphisms-lemma-etale-separated-ind-quasi-affine} 得出
（因为 \'etale 态射局部拟有限；见《态射》之引理
\ref{morphisms-lemma-etale-locally-quasi-finite}）。由《群胚更多内容》之引理
\ref{more-groupoids-lemma-ind-quasi-affine}，下降数据是有效的。
设 $W\to T$ 为一个态射，使存在同构
$\alpha:T'\times_TW\to V$，该同构与 $V$ 上给定的下降数据以及
$T'\times_TW$ 上的典型下降数据相容。则 $W\to T$ 是满且 \'etale 的
（《下降》之引理 \ref{descent-lemma-descending-property-surjective} 和
\ref{descent-lemma-descending-property-etale}）。考虑复合
$$
b' : T' \times_T W \longrightarrow V = U \times_{X, c} T' \longrightarrow U
$$
由我们对 $\alpha$ 的选择以及 $c$ 的相应性质，两个复合
$b' \circ (\text{pr}_0, 1),
b' \circ (\text{pr}_1, 1) :
(T' \times_T T') \times_T W \to T' \times_T W \to U$
相同（计算略）。因此由《下降》之引理
\ref{descent-lemma-fpqc-universal-effective-epimorphisms}，$b'$ 下降为态射
$b:W\to U$。图表
$$
\xymatrix{
T' \times_T W \ar[r] \ar[d] & W \ar[r]_b & U \ar[d] \\
T' \ar[rr]^c &  & X
}
$$
交换。这意味着我们已经在 $T$ 的 \'etale 局部找到了 $a$，即得到
$a':W\to X$。然而，由于第一段已经证明唯一性，这个 \'etale 局部解满足粘合条件，
即 $\text{pr}_0^*a'=\text{pr}_1^*a'$ 作为 $X(W\times_TW)$ 的元素成立。
由于 $X$ 是 \'etale 层，得到唯一的 $a\in X(T)$，其在 $W$ 上限制为 $a'$。
\end{proof}





\section{从仿射形式概形出发的映射}
\label{formal-spaces-section-map-out-of}

\noindent
我们证明一些稍后会有用的结果。在论文 \cite{Bhatt-Algebraize} 中，读者可以找到
同类性质的非常一般的结果。

\begin{lemma}
\label{formal-spaces-lemma-map-into-affine}
设 $S$ 为概形，$A$ 为弱可容许拓扑 $S$-代数，$X$ 为 $S$ 上的仿射概形。
则自然映射
$$
\Mor_S(\Spec(A), X)
\longrightarrow
\Mor_S(\text{Spf}(A), X)
$$
是双射。
\end{lemma}

\begin{proof}
若 $X$ 是仿射的，写作 $X=\Spec(B)$，则由引理
\ref{formal-spaces-lemma-morphism-between-formal-spectra}，态射
$\text{Spf}(A)\to\Spec(B)$ 对应于连续的 $S$-代数映射 $B\to A$，
其中 $B$ 具有离散拓扑。这些正是 $S$-代数映射，对应于
$\Spec(A)\to\Spec(B)$ 的态射。
\end{proof}

\begin{lemma}
\label{formal-spaces-lemma-map-into-scheme}
设 $S$ 为概形，$A$ 为弱可容许拓扑 $S$-代数，并且对某个弱定义理想
$I\subset A$，$A/I$ 是局部环。设 $X$ 为 $S$ 上的概形。则自然映射
$$
\Mor_S(\Spec(A), X)
\longrightarrow
\Mor_S(\text{Spf}(A), X)
$$
是双射。
\end{lemma}

\begin{proof}
设 $\varphi:\text{Spf}(A)\to X$ 为态射。由于 $\Spec(A/I)$ 是局部的，
可知 $\varphi$ 将 $\Spec(A/I)$ 映入 $X$ 的某个仿射开子集 $U$。
然而，这又意味着对每个定义理想 $J$，$\Spec(A/J)$ 都映入 $U$。
因此可以应用引理 \ref{formal-spaces-lemma-map-into-affine}，可知 $\varphi$ 来自一个
$\Spec(A)\to X$ 的态射。这证明了该映射的满射性。
单射性的证明略去。
\end{proof}

\begin{lemma}
\label{formal-spaces-lemma-map-into-algebraic-space}
设 $S$ 为概形，$R$ 为完备局部 Noether $S$-代数，$X$ 为 $S$ 上的代数空间。
则自然映射
$$
\Mor_S(\Spec(R), X)
\longrightarrow
\Mor_S(\text{Spf}(R), X)
$$
是双射。
\end{lemma}

\begin{proof}
设 $\mathfrak m$ 为 $R$ 的极大理想。我们需要证明，对这样的 $R$，映射
$$
\Mor_S(\Spec(R), X) \longrightarrow
\lim \Mor_S(\Spec(R/\mathfrak m^n), X)
$$
是双射。

\medskip\noindent
单射性：设 $x,x':\Spec(R)\to X$ 为两个在右侧映到同一元素的态射。
考虑纤维积
$$
T = \Spec(R) \times_{(x, x'), X \times_S X, \Delta} X
$$
则 $T$ 是概形，且 $T\to\Spec(R)$ 局部有限型、为单态射、分离且局部拟有限；见《空间的态射》之引理
\ref{spaces-morphisms-lemma-properties-diagonal}。特别地，$T$ 局部 Noether；见《态射》之引理
\ref{morphisms-lemma-finite-type-noetherian}。令 $t\in T$ 为映到
$\Spec(R)$ 闭点的唯一点；由于 $x$ 与 $x'$ 在 $R/\mathfrak m$ 上相同，这一点存在。
则 $R\to\mathcal{O}_{T,t}$ 是 Noether 环的局部环映射，并且对所有 $n$，
$R/\mathfrak m^n\to\mathcal{O}_{T,t}/\mathfrak m^n\mathcal{O}_{T,t}$ 是同构
（因为 $x$ 和 $x'$ 在所有 $\Spec(R/\mathfrak m^n)$ 上相同）。
由于 $\mathcal{O}_{T,t}$ 单射地映入其完备化（见《代数》之引理
\ref{algebra-lemma-intersect-powers-ideal-module-zero}），可知
$R=\mathcal{O}_{T,t}$。因此 $x$ 和 $x'$ 在 $R$ 上相同。

\medskip\noindent
满射性：设 $(x_n)$ 是右侧的一个元素。选取概形 $U$ 及满的 \'etale 态射
$U\to X$。令 $x_0:\Spec(k)\to X$ 为在剩余域
$k=R/\mathfrak m$ 上诱导的态射。概形态射
$U\times_{X,x_0}\Spec(k)\to\Spec(k)$ 是满的 \'etale 态射。
因此 $U\times_{X,x_0}\Spec(k)$ 是 $k$ 的有限可分扩张的谱所成的非空互不相交并；见《态射》之引理
\ref{morphisms-lemma-etale-over-field}。于是可以找到有限可分扩张 $k'/k$ 和
$k'$-点 $u_0:\Spec(k')\to U$，使图表
$$
\xymatrix{
\Spec(k') \ar[d] \ar[r]_-{u_0} & U \ar[d] \\
\Spec(k) \ar[r]^-{x_0} & X
}
$$
交换。令 $R\subset R'$ 为 Noether 完备局部环的有限 \'etale 扩张，
其在剩余域上诱导 $k'/k$（见《代数》之引理
\ref{algebra-lemma-henselian-cat-finite-etale} 和
\ref{algebra-lemma-complete-henselian}）。令 $x'_n$ 为 $x_n$ 在
$\Spec(R'/\mathfrak m^nR')$ 上的限制。由《空间的态射更多内容》之引理
\ref{spaces-more-morphisms-lemma-etale-formally-etale}，可以找到
$(u'_n)\in\lim\Mor_S(\Spec(R'/\mathfrak m^nR'),U)$，其映到 $(x'_n)$。
由引理 \ref{formal-spaces-lemma-map-into-scheme}，族 $(u'_n)$ 来自唯一的态射
$u':\Spec(R')\to U$。令 $x':\Spec(R')\to X$ 为其复合。
注意，$R'\otimes_RR'$ 是完备局部 Noether 环的谱的有限积，仍适用当前讨论。
因此图表
$$
\xymatrix{
\Spec(R' \otimes_R R') \ar[r] \ar[d] & \Spec(R') \ar[d]^{x'} \\
\Spec(R') \ar[r]^{x'} & X
}
$$
由上面证明的单射性以及 $x'_n$ 是 $x_n$ 的限制（后者在
$R/\mathfrak m^n$ 上定义）而交换。由于 $\{\Spec(R')\to\Spec(R)\}$
是 fppf 覆盖，故 $x'$ 下降为态射 $x:\Spec(R)\to X$。
证明 $x_n$ 是 $x$ 在 $\Spec(R/\mathfrak m^n)$ 上的限制这一点略去。
\end{proof}

\begin{lemma}
\label{formal-spaces-lemma-adic-into-completion}
设 $S$ 为概形，$X$ 为 $S$ 上的代数空间。设 $T\subset |X|$ 为闭子集，且
$X\setminus T\to X$ 拟紧。设 $R$ 为完备局部 Noether $S$-代数。
则 adic 态射 $p:\text{Spf}(R)\to X_{/T}$ 对应于唯一的态射
$g:\Spec(R)\to X$，使得 $g^{-1}(T)=\{\mathfrak m_R\}$。
\end{lemma}

\begin{proof}
由于引理 \ref{formal-spaces-lemma-formal-completion-types}，$X_{/T}$ 是 adic* 的，
所以该陈述有意义（并且可以使用 adic 态射这一术语；见定义
\ref{formal-spaces-definition-adic-morphism}）。设 $p$ 已给定。由引理
\ref{formal-spaces-lemma-map-into-algebraic-space}，得到唯一态射 $g:\Spec(R)\to X$，
它对应于复合 $\text{Spf}(R)\to X_{/T}\to X$。
令 $Z\subset X$ 为 $T$ 上的约化诱导闭子空间结构。
% P10-E0371
包含态射 $Z\to X$ 对应于一个态射 $Z\to X_{/T}$。由于 $p$ 是 adic 的，
它可由代数空间表示，于是
$$
\text{Spf}(R) \times_{X_{/T}} Z = \text{Spf}(R) \times_X Z
$$
是带有到 $\text{Spf}(R)$ 的闭浸入的代数空间。（等式成立是因为
$X_{/T}\to X$ 是单态射。）因此这个纤维积等于 $\Spec(R/J)$，其中
理想 $J\subset R$ 包含某个 $n_0\geq1$ 的 $\mathfrak m_R^{n_0}$。
这意味着 $\Spec(R)\times_XZ$ 是 $\Spec(R)$ 的闭子概形，设
$\Spec(R)\times_XZ=\Spec(R/I)$，其与 $n\geq n_0$ 时的
$\Spec(R/\mathfrak m_R^n)$ 的交等于 $\Spec(R/J)$。
用代数语言说，这意味着对所有 $n\geq n_0$，
$I+\mathfrak m_R^n=J+\mathfrak m_R^n=J$。由 Krull 交定理，
这说明 $I=J$，结论得证。
\end{proof}





\section{形式代数空间的小 \'etale 位点}
\label{formal-spaces-section-etale-sites}

\noindent
下述定义的动机来自经典形式概形：形式概形
$(\mathfrak X,\mathcal{O}_\mathfrak X)$ 的底层拓扑空间是约化
$\mathfrak X_{red}$ 的底层拓扑空间。

\medskip\noindent
还有一个重要备注。假设 $X$ 是带有约化 $X_{red}$ 的代数空间
（《空间的性质》之定义 \ref{spaces-properties-definition-reduced-induced-space}）。则由《空间的态射更多内容》之定理
\ref{spaces-more-morphisms-theorem-topological-invariance} 和引理
\ref{spaces-more-morphisms-lemma-topological-invariance}，有
$$
X_{spaces, \etale} = X_{red, spaces, \etale},\quad
X_\etale = X_{red, \etale},\quad
X_{affine, \etale} = X_{red, affine, \etale}
$$
因此，在我们的形式代数空间恰好是代数空间的情形，下述定义不会与已有概念冲突。

\begin{definition}
\label{formal-spaces-definition-etale-sites}
设 $S$ 为概形，设 $X$ 为带有约化 $X_{red}$ 的形式代数空间
（引理 \ref{formal-spaces-lemma-reduction-formal-algebraic-space}）。
\begin{enumerate}
\item $X$ 的\textit{小 \'etale 位点} $X_\etale$ 是《空间的性质》之定义
\ref{spaces-properties-definition-etale-site} 中的位点 $X_{red,\etale}$；
\item 位点 $X_{spaces,\etale}$ 是《空间的性质》之定义
\ref{spaces-properties-definition-spaces-etale-site} 中的位点
$X_{red,spaces,\etale}$；
\item 位点 $X_{affine,\etale}$ 是《空间的性质》之引理
\ref{spaces-properties-lemma-alternative} 中的位点
$X_{red,affine,\etale}$。
\end{enumerate}
\end{definition}

\noindent
在引理 \ref{formal-spaces-lemma-identify-spaces-etale} 中我们将看到，
% P10-E0372
$X_{spaces,\etale}$ 可以用形式代数空间的态射来描述，这些态射可由代数空间表示且为 \'etale。
由《空间的性质》之引理
\ref{spaces-properties-lemma-compare-etale-sites} 和
\ref{spaces-properties-lemma-alternative}，有恒等
\begin{equation}
\label{formal-spaces-equation-etale-topos}
\Sh(X_\etale)=\Sh(X_{spaces,\etale})=\Sh(X_{affine,\etale})
\end{equation}
我们称之为 $X$ 的\textit{（小）\'etale 拓扑斯}。

\begin{lemma}
\label{formal-spaces-lemma-functoriality-etale-site}
设 $S$ 为概形，设 $f:X\to Y$ 为 $S$ 上形式代数空间的态射。
\begin{enumerate}
\item 存在连续函子
$Y_{spaces,\etale}\to X_{spaces,\etale}$，它诱导位点的态射
$$
f_{spaces, \etale}:X_{spaces,\etale}\to Y_{spaces,\etale}.
$$
\item 规则 $f\mapsto f_{spaces,\etale}$ 与复合相容，换言之
$(f\circ g)_{spaces,\etale}=f_{spaces,\etale}\circ g_{spaces,\etale}$；见《位点》之定义
\ref{sites-definition-composition-morphisms-sites}。
\item 与 $f_{spaces,\etale}$ 关联的拓扑斯态射，通过
\ref{formal-spaces-equation-etale-topos} 诱导拓扑斯态射
$f_{small}:\Sh(X_\etale)\to\Sh(Y_\etale)$，且其构造与复合相容。
\end{enumerate}
\end{lemma}

\begin{proof}
这里唯一要注意的是，由引理 \ref{formal-spaces-lemma-reduction-formal-algebraic-space}，
$f$ 诱导约化之间的态射 $X_{red}\to Y_{red}$。因此本引理立即由
代数空间态射的对应引理（《空间的性质》之引理
\ref{spaces-properties-lemma-functoriality-etale-site}）得到。
\end{proof}

\noindent
如果形式代数空间的态射 $X\to Y$ 是 \'etale 的，则拓扑斯态射
$\Sh(X_\etale)\to\Sh(Y_\etale)$ 是局部化。下面给出一个陈述。

\begin{lemma}
\label{formal-spaces-lemma-etale-morphism-topoi}
设 $S$ 为概形，$f:X\to Y$ 为 $S$ 上形式代数空间的态射。
假设 $f$ 可由代数空间表示且为 \'etale。此时存在余连续函子
$j:X_\etale\to Y_\etale$。拓扑斯态射 $f_{small}$ 是与 $j$ 关联的拓扑斯态射；见《位点》之引理
\ref{sites-lemma-cocontinuous-morphism-topoi}。此外，$j$ 也是连续的，
所以《位点》之引理 \ref{sites-lemma-when-shriek} 适用。
\end{lemma}

\begin{proof}
只要证明诱导的态射 $X_{red}\to Y_{red}$ 是 \'etale 的，结论就立即由代数空间的情形
（《空间的性质》之引理 \ref{spaces-properties-lemma-etale-morphism-topoi}）得到。
注意 $X\times_YY_{red}$ 是在约化代数空间 $Y_{red}$ 上 \'etale 的代数空间，
因而本身是约化的（依《空间的性质》第
\ref{spaces-properties-section-types-properties} 节中约化代数空间的定义）。
因此 $X_{red}=X\times_YY_{red}$，如所需。
\end{proof}

\begin{lemma}
\label{formal-spaces-lemma-affine-identify-affine-etale}
设 $S$ 为概形，$X$ 为 $S$ 上的仿射形式代数空间。则
$X_{affine,\etale}$ 等价于下列对象所成的范畴：对象是形式代数空间的态射
$\varphi:U\to X$，满足
\begin{enumerate}
\item $U$ 是仿射形式代数空间；
\item $\varphi$ 可由代数空间表示且为 \'etale。
\end{enumerate}
\end{lemma}

\begin{proof}
记 $\mathcal{C}$ 为引理中引入的范畴。注意，对 $\mathcal{C}$ 中的
$\varphi:U\to X$，$\varphi$ 可由概形表示且为仿射；见引理
\ref{formal-spaces-lemma-affine-representable-by-algebraic-spaces}。回忆
$X_{affine,\etale}=X_{red,affine,\etale}$。因此可以定义函子
$$
\mathcal{C}\longrightarrow X_{affine,\etale},\quad
(U\to X)\longmapsto U\times_XX_{red}
$$
因为 $U\times_XX_{red}$ 是仿射概形。

\medskip\noindent
为完成证明，我们构造准逆。按定义
\ref{formal-spaces-definition-affine-formal-algebraic-space}，写作 $X=\colim X_\lambda$。
对每个 $\lambda$，$X_{red}\subset X_\lambda$ 是一个增厚。因此对每个 $\lambda$ 有等价
$$
X_{red,affine,\etale}=X_{\lambda,affine,\etale}
$$
例如由《代数更多内容》之引理
\ref{more-algebra-lemma-locally-nilpotent-henselian} 得到。因此，若
$U_{red}\to X_{red}$ 是仿射的 \'etale 态射，则得到一族相容的仿射概形 \'etale 态射
$U_\lambda\to X_\lambda$，其与定义 $X$ 的系统中的过渡态射相容。于是令
$$
U=\colim U_\lambda
$$
作为 $X$ 上的仿射形式代数空间。构造给出
$U\times_XX_\lambda=U_\lambda$。这说明 $U\to X$ 可由代数空间表示且为 \'etale。
构造彼此为逆这一点的核验略去。
\end{proof}

\begin{lemma}
\label{formal-spaces-lemma-affine-etale-mcquillan}
设 $S$ 为概形，$X$ 为 $S$ 上的仿射形式代数空间。假设 $X$ 是 McQuillan，即
$X=\text{Spf}(A)$，其中 $A$ 是弱可容许拓扑 $S$-代数。
则 $(X_{affine,\etale})^{opp}$ 等价于下列范畴：
\begin{enumerate}
\item 对象是形如 $B^\wedge=\lim B/JB$ 的 $A$-代数，其中 $A\to B$ 是 \'etale 环映射，
且 $J$ 遍历 $A$ 的弱定义理想；
\item 态射是连续的 $A$-代数同态。
\end{enumerate}
\end{lemma}

\begin{proof}
结合引理 \ref{formal-spaces-lemma-affine-identify-affine-etale} 和引理
\ref{formal-spaces-lemma-etale} 即得。
\end{proof}

\begin{lemma}
\label{formal-spaces-lemma-identify-spaces-etale}
设 $S$ 为概形，$X$ 为 $S$ 上的形式代数空间。则 $X_{spaces,\etale}$ 等价于下列对象所成的范畴：
对象是形式代数空间的态射 $\varphi:U\to X$，且 $\varphi$ 可由代数空间表示并为 \'etale。
\end{lemma}

\begin{proof}
记 $\mathcal{C}$ 为引理中引入的范畴。
回忆 $X_{spaces, \etale} = X_{red, spaces, \etale}$。
因此可以定义函子
$$
\mathcal{C} \longrightarrow X_{spaces, \etale},\quad
(U \to X) \longmapsto U \times_X X_{red}
$$
因为 $U \times_X X_{red}$ 是在 $X_{red}$ 上 \'etale 的代数空间。

\medskip\noindent
为完成证明，我们构造一个准逆函子。
取 $X_{red, spaces, \etale}$ 中的对象
$\psi : V \to X_{red}$。
考虑函子
$U_{V, \psi} : (\Sch/S)_{fppf} \to \textit{Sets}$，其定义为
$$
U_{V, \psi}(T) = \{(a, b) \mid
a : T \to X,
\ b : T \times_{a, X} X_{red} \to V,
\ \psi \circ b = a|_{T \times_{a, X} X_{red}}\}
$$
我们断言，变换 $U_{V, \psi} \to X$，$(a, b) \mapsto a$，
定义了范畴 $\mathcal{C}$ 中的一个对象。
首先证明 $U_{V, \psi}$ 是形式代数空间。
注意，$U_{V, \psi}$ 是 fppf 位点上的层（略去一些细节）。
接下来，设 $X_i \to X$ 是由仿射形式代数空间给出的 \'etale 覆盖，
如定义 \ref{formal-spaces-definition-formal-algebraic-space} 中所述。
令 $V_i = V \times_{X_{red}} X_{i, red}$，并记
$\psi_i : V_i \to X_{i, red}$ 为投影。则
我们有
$$
U_{V, \psi} \times_X X_i = U_{V_i, \psi_i}
$$
这是形式上的论证，因为 $X_{i, red} = X_i \times_X X_{red}$
（这是由于 $X_i \to X$ 可由代数空间表示且为 \'etale）。
因此，只需证明 $U_{V_i, \psi_i}$ 是仿射形式代数空间，
因为这样就有一个覆盖 $U_{V_i, \psi_i} \to U_{V, \psi}$，
满足定义 \ref{formal-spaces-definition-formal-algebraic-space}。
另一方面，我们在引理 \ref{formal-spaces-lemma-etale-morphism-topoi} 的证明中已经看到，
$\psi_i : V_i \to X_i$ 是仿射形式代数空间之间的可表示且 \'etale
态射 $U_i \to X_i$ 的基变换。
于是，不难看出 $U_i = U_{V_i, \psi_i}$，如所需。

\medskip\noindent
我们略去验证 $U_{V, \psi} \to X$ 可由代数空间表示且为 \'etale。
这样便得到另一个方向的函子
$(V, \psi) \mapsto (U_{V, \psi} \to X)$。
构造彼此互逆的验证略去。
\end{proof}

\begin{lemma}
\label{formal-spaces-lemma-identify-affine-etale}
设 $S$ 为概形，$X$ 为 $S$ 上的形式代数空间。
则 $X_{affine, \etale}$ 等价于下列对象所成的范畴：对象是形式代数空间的态射
$\varphi : U \to X$，满足
\begin{enumerate}
\item $U$ 是仿射形式代数空间，
\item $\varphi$ 可由代数空间表示且为 \'etale。
\end{enumerate}
\end{lemma}

\begin{proof}
这由引理 \ref{formal-spaces-lemma-identify-spaces-etale} 和
\ref{formal-spaces-lemma-characterize-affine} 合并得到。
\end{proof}





\section{结构层}
\label{formal-spaces-section-structure-sheaf}


\noindent
如第 \ref{formal-spaces-section-etale-sites} 节所述，本节讨论的动机
来自经典形式概形。经典形式概形定义为二元组
$(\mathfrak X, \mathcal{O}_\mathfrak X)$，其中 $\mathcal{O}_\mathfrak X$
是拓扑环层。

\medskip\noindent
设 $X$ 为形式代数空间。$X$ 的一个{\it 结构层}
是 \'etale 位点 $X_\etale$ 上的拓扑环层 $\mathcal{O}_X$
（该位点在第 \ref{formal-spaces-section-etale-sites} 节中定义），并且对所有满足下列条件的情形，
都有
$$
\mathcal{O}_X(U_{red}) = \lim \Gamma(U_\lambda, \mathcal{O}_{U_\lambda})
$$
作为拓扑环成立：
\begin{enumerate}
\item $\varphi : U \to X$ 是形式代数空间的态射，
\item $U$ 是仿射形式代数空间，
\item $\varphi$ 可由代数空间表示且为 \'etale，
\item $U_{red} \to X_{red}$ 是 $X_\etale$ 中相应的仿射对象，见
引理 \ref{formal-spaces-lemma-identify-affine-etale}，
\item $U = \colim U_\lambda$ 是定义
\ref{formal-spaces-definition-affine-formal-algebraic-space} 中所述的 $U$ 的余极限表示。
\end{enumerate}
结构层确实存在，但其行为可能出人意料。

\begin{lemma}
\label{formal-spaces-lemma-structure-sheaf}
每个形式代数空间都有结构层。
\end{lemma}

\begin{proof}
设 $S$ 为概形，设 $X$ 为 $S$ 上的形式代数空间。
由（\ref{formal-spaces-equation-etale-topos}），只需在 $X_{affine, \etale}$ 上构造
$\mathcal{O}_X$ 为拓扑环层即可。
记 $\mathcal{C}$ 为如下范畴：对象是形式代数空间的态射
$\varphi : U \to X$，其中 $U$ 是仿射形式代数空间，
且 $\varphi$ 可由代数空间表示并为 \'etale。
由引理 \ref{formal-spaces-lemma-identify-affine-etale}，函子
$U \mapsto U_{red}$ 是等价
$\mathcal{C} \to X_{affine, \etale}$。
因此，按引理之前给出的规则，我们已经在
$X_{affine, \etale}$ 上得到拓扑环的预层 $\mathcal{O}_X$。
所以只需验证层条件。

\medskip\noindent
按 $X_{affine, \etale}$ 的定义，一个覆盖对应于有限族
$\{g_i : U_i \to U\}_{i = 1, \ldots, n}$，其中 $g_i$ 是 $\mathcal{C}$ 中的态射，
并且 $\{U_{i, red} \to U_{red}\}$ 是 \'etale 覆盖。
$g_i$ 可由代数空间表示（引理
\ref{formal-spaces-lemma-permanence-representable}），因而是仿射的（引理
\ref{formal-spaces-lemma-affine-representable-by-algebraic-spaces}）。
% P10-E0373
于是 $g_i$ 是 \'etale 的（这形式地由《空间的性质》之引理
\ref{spaces-properties-lemma-etale-permanence} 得到，因为 $U_i$ 与 $U$
在《Bootstrap》第 \ref{bootstrap-section-representable-by-spaces-properties} 节的意义下
均 \'etale 地态射到 $X$）。
最后，按定义 \ref{formal-spaces-definition-affine-formal-algebraic-space} 写
$U = \colim U_\lambda$。

\medskip\noindent
准备工作完成后，层性质如下证明。对每个 $\lambda$，令
$U_{i, \lambda} = U_i \times_U U_\lambda$，并令
$U_{ij, \lambda} = (U_i \times_U U_j) \times_U U_\lambda$。
由上述论证，这些是仿射概形，$\{U_{i, \lambda} \to U_\lambda\}$ 是 \'etale 覆盖，
且
$U_{ij, \lambda} = U_{i, \lambda} \times_{U_\lambda} U_{j, \lambda}$。
此外有 $U_i = \colim U_{i, \lambda}$ 以及
$U_i \times_U U_j = \colim U_{ij, \lambda}$。
对每个 $\lambda$，有正合序列
$$
0 \to
\Gamma(U_\lambda, \mathcal{O}_{U_\lambda}) \to
\prod\nolimits_i \Gamma(U_{i, \lambda}, \mathcal{O}_{U_{i, \lambda}}) \to
\prod\nolimits_{i, j} \Gamma(U_{ij, \lambda}, \mathcal{O}_{U_{ij, \lambda}})
$$
因为在 $U_\lambda$ 上、关于 \'etale 拓扑，结构层满足层条件
（见《\'Etale 上同调》之命题
\ref{etale-cohomology-proposition-quasi-coherent-sheaf-fpqc}）。
由于极限与极限交换，这些正合序列的逆极限仍为正合序列：
$$
0 \to
\lim \Gamma(U_\lambda, \mathcal{O}_{U_\lambda}) \to
\prod\nolimits_i \lim \Gamma(U_{i, \lambda}, \mathcal{O}_{U_{i, \lambda}}) \to
\prod\nolimits_{i, j} \lim
\Gamma(U_{ij, \lambda}, \mathcal{O}_{U_{ij, \lambda}})
$$
这意味着
$$
0 \to
\mathcal{O}_X(U_{red}) \to
\prod\nolimits_i \mathcal{O}_X(U_{i, red}) \to
\prod\nolimits_{i, j} \mathcal{O}_X((U_i \times_U U_j)_{red})
$$
作为阿贝尔群序列是正合的。我们仍需证明
$\mathcal{O}_X(U_{red})$ 上的极限拓扑由
$\prod_{i = 1, \ldots, n} \mathcal{O}_X(U_{i, red})$ 上的极限拓扑诱导。
映射
$\mathcal{O}_X(U_{red}) \to \mathcal{O}_{U_\lambda}(U_\lambda)$
的核构成开子模的基本系。由于这些核是映射
$\prod_i \mathcal{O}_X(U_{i, red}) \to
\prod_i \mathcal{O}_{U_{i, \lambda}}(U_{i, \lambda})$
的核的逆像，而这些核构成
% P10-E0374
$\mathcal{O}_X(U_{red}) \to \mathcal{O}_{U_\lambda}(U_\lambda)$
的开子模基本系，结论得证。
\end{proof}

\begin{remark}
\label{formal-spaces-remark-not-enough-sections}
结构层并不总是有“足够多的截面”。
在例子（第 \ref{examples-section-affine-formal-algebraic-space} 节）中，
我们已经看到存在不是 McQuillan 的仿射形式代数空间，
甚至还有一些例子，其点不能由正则函数分离。
\end{remark}


\noindent
在下一个引理中，我们证明可数指标的仿射形式代数空间上的结构层
具有消失的高阶上同调。
对于非可数指标的情形，推测这一结论一般不成立。

\begin{lemma}
\label{formal-spaces-lemma-higher-vanishing-structure-sheaf}
若 $X$ 是可数指标的仿射形式代数空间，则当 $n > 0$ 时有
$H^n(X_\etale, \mathcal{O}_X) = 0$。
\end{lemma}

\begin{proof}
我们可以在 $X_{affine, \etale}$ 上工作，因为它给出同一个拓扑斯。
我们将应用《位点上的上同调》之引理
\ref{sites-cohomology-lemma-cech-vanish-collection}
来证明消失性。由于 $X_{affine, \etale}$ 有有限不交并，
这将问题化为覆盖的 Čech 复形，其中覆盖由
$X_{affine, \etale} = X_{red, affine, \etale}$ 中的单个箭头
$\{U_{red} \to V_{red}\}$ 给出
（见《\'Etale 上同调》之引理
\ref{etale-cohomology-lemma-cech-complex}）。
因此我们必须证明
$$
0 \to \mathcal{O}_X(V_{red}) \to \mathcal{O}_X(U_{red}) \to
\mathcal{O}_X(U_{red} \times_{V_{red}} U_{red}) \to \ldots
$$
是正合的。下面在 $V_{red} = X_{red}$ 的情形证明这一点；
一般情形的证明完全相同。

\medskip\noindent
回忆 $X = \text{Spf}(A)$，其中 $A$ 是弱可容许拓扑环，
并且具有可数的弱定义理想基本系。
由引理 \ref{formal-spaces-lemma-affine-identify-affine-etale} 和
\ref{formal-spaces-lemma-affine-etale-mcquillan}，$X_{affine, \etale}$ 中的对象
$U_{red}$ 对应于仿射形式代数空间的态射 $U \to X$，
该态射可由代数空间表示且为 \'etale，并且
$U = \text{Spf}(B^\wedge)$，其中 $B$ 是 \'etale 的 $A$-代数。
根据结构层的定义规则，我们看到
$$
\mathcal{O}_X(U_{red}) = B^\wedge
$$
回忆 $B^\wedge = \lim B/JB$，其中极限遍历
$A$ 的弱定义理想 $J \subset A$。
逐一考察定义可得
$$
\mathcal{O}_X(U_{red} \times_{X_{red}} U_{red}) = (B \otimes_A B)^\wedge
$$
依此类推。由于 $U \to X$ 是覆盖，映射 $A \to B$
是忠实平坦的，见引理 \ref{formal-spaces-lemma-etale-surjective}。
因此复形
$$
0 \to A \to B \to B \otimes_A B \to B \otimes_A B \otimes_A B \to \ldots
$$
是{\bf 泛正合}的，见《下降》之引理
\ref{descent-lemma-ff-exact}。
我们的目标是证明
$$
H^n(0 \to A^\wedge \to B^\wedge \to (B \otimes_A B)^\wedge
\to (B \otimes_A B \otimes_A B)^\wedge \to \ldots)
$$
当 $n > 0$ 时为零。为看清其中的机制，将上述（完备化之前的）
正合复形拆成短正合序列
% P10-E0375
$$
0 \to A \to B \to M_1 \to 0,\quad
0 \to M_i \to B^{\otimes_A i + 1} \to M_{i + 1} \to 0
$$
由上面的论述，这些都是泛正合的短正合序列。
因此对 $A$ 的每个理想 $J$，都有
$JM_i = M_i \cap J(B^{\otimes_A i + 1})$。
特别地，$M_i$ 作为 $B^{\otimes_A i + 1}$ 的子模所承继的拓扑，
与 $M_i$ 作为 $B^{\otimes_A i}$ 的商模所承继的拓扑相同。
由于 $A$ 中存在可数的弱定义理想基本系，序列
$$
0 \to A^\wedge \to B^\wedge \to M_1^\wedge \to 0,\quad
0 \to M_i^\wedge \to (B^{\otimes_A i + 1})^\wedge \to M_{i + 1}^\wedge \to 0
$$
由引理 \ref{formal-spaces-lemma-ses} 保持正合。这就证明了引理。
\end{proof}

\begin{remark}
\label{formal-spaces-remark-bad-quasi-coherent}
即使结构层具有良好性质，也不意味着存在良好的拟凝聚模理论。
例如，在例子（第 \ref{examples-section-nonabelian-QCoh} 节）中，
我们看到，对于几乎所有诺特仿射形式代数空间，最自然的拟凝聚模概念
都会导出一个非阿贝尔的模范畴。
\end{remark}





\section{形式代数空间的余极限}
\label{formal-spaces-section-colimits-of-formal}

\noindent
本节将第 \ref{formal-spaces-section-global-colimits} 节的结果推广到
形式代数空间的态射系。
我们指出，在下面的引理中，条件
“$f_{\lambda \mu} : X_\lambda \to X_\mu$ 是闭浸入，
并诱导约化空间 $X_{\lambda, red} \to X_{\mu, red}$ 的同构”
可以改写为
“$f_{\lambda \mu}$ 可由代数空间表示且为增厚”。

\begin{lemma}
\label{formal-spaces-lemma-colimit-affine-formal}
设 $S$ 为概形。设给定有向集 $\Lambda$ 以及其上的仿射形式代数空间系统
$(X_\lambda, f_{\lambda \mu})$，其中每个
$f_{\lambda \mu} : X_\lambda \to X_\mu$ 都是闭浸入，
并诱导约化空间 $X_{\lambda, red} \to X_{\mu, red}$ 的同构。
则 $X = \colim_{\lambda \in \Lambda} X_\lambda$
是 $S$ 上的仿射形式代数空间。
\end{lemma}

\begin{proof}
我们可以写成
$X_\lambda = \colim_{\omega \in \Omega_\lambda} X_{\lambda, \omega}$，
其中这是有向集 $\Omega_\lambda$ 上仿射概形的余极限，
且过渡态射
$X_{\lambda, \omega} \to X_{\lambda, \omega'}$
是增厚。对每个 $\lambda, \mu \in \Lambda$ 以及
$\omega \in \Omega_\lambda$，当 $\mu \geq \lambda$ 时，
存在 $\omega' \in \Omega_\mu$，使得态射
$X_{\lambda, \omega} \to X_\mu$ 分解为经过 $X_{\mu, \omega'}$，
见引理 \ref{formal-spaces-lemma-factor-through-thickening}。于是
态射 $X_{\lambda, \omega} \to X_{\mu, \omega'}$
是闭浸入，并在约化上诱导同构，因而是增厚。令
$\Omega = \coprod_{\lambda \in \Lambda} \Omega_\lambda$，
并规定 $(\lambda, \omega) \leq (\mu, \omega')$ 当且仅当
$\lambda \leq \mu$ 且 $X_{\lambda, \omega} \to X_\mu$ 分解为经过
$X_{\mu, \omega'}$。由上述事实可知，$\Omega$ 是有向集，并且
$X = \colim_{\lambda \in \Lambda} X_\lambda =
\colim_{(\lambda, \omega) \in \Omega} X_{\lambda, \omega}$。
证明完成。
\end{proof}

\begin{lemma}
\label{formal-spaces-lemma-colimit-formal-spaces-is-formal-space}
设 $S$ 为概形。设给定有向集 $\Lambda$ 以及其上的形式代数空间系统
$(X_\lambda, f_{\lambda \mu})$，其中每个
$f_{\lambda \mu} : X_\lambda \to X_\mu$ 都是闭浸入，
并诱导约化空间 $X_{\lambda, red} \to X_{\mu, red}$ 的同构。
则 $X = \colim_{\lambda \in \Lambda} X_\lambda$
是 $S$ 上的形式代数空间。
\end{lemma}

\begin{proof}
由于我们在 fppf 层范畴中取余极限，可知 $X$ 是一个层。
选取并固定 $\lambda \in \Lambda$。按定义
\ref{formal-spaces-definition-formal-algebraic-space} 选取覆盖
$\{X_{i, \lambda} \to X_\lambda\}$。特别地，
$\{X_{i, \lambda, red} \to X_{\lambda, red}\}$ 是由仿射概形给出的
\'etale 覆盖。
对每个 $\mu \geq \lambda$，存在笛卡尔图
$$
\xymatrix{
X_{i, \lambda} \ar[r] \ar[d] & X_{i, \mu} \ar[d] \\
X_\lambda \ar[r] & X_\mu
}
$$
其竖直箭头为 \'etale。事实上，\'etale 态射
$X_{i, \lambda, red} \to X_{\lambda, red} = X_{\mu, red}$
对应于形式代数空间的 \'etale 态射 $X_{i, \mu} \to X_\mu$，
其中 $X_{i, \mu}$ 是仿射形式代数空间，见
引理 \ref{formal-spaces-lemma-affine-identify-affine-etale}。
同一个引理说明 $X_{i, \mu}$ 到 $X_\lambda$ 的基变换
与 $X_{i, \lambda}$ 一致。还可推出，当
$\mu' \geq \mu \geq \lambda$ 时，
$X_{i, \mu} = X_\mu \times_{X_{\mu'}} X_{i, \mu'}$。
令 $X_i = \colim X_{i, \mu}$。则
$X_{i, \mu} = X_i \times_X X_\mu$（作为函子）。
由于从仿射（或拟紧）概形 $T$ 到
$X = \colim X_\mu$ 的任意态射都落入某个 $X_\mu$，
% P10-E0376
因此可得 $\colim X_{i, \mu} \to \colim X_\mu$ 是 \'etale。
于是，只要证明 $\colim X_{i, \mu}$ 是仿射形式代数空间，
引理便得证。
注意，态射 $X_{i, \mu} \to X_{i, \mu'}$
是闭浸入 $X_\mu \to X_{\mu'}$ 的基变换，因而是闭浸入。
最后，由于 $X_{\mu, red} \to X_{\mu', red}$ 是同构，
态射 $X_{i, \mu, red} \to X_{i, \mu', red}$ 也是同构。
因此问题归结为引理 \ref{formal-spaces-lemma-colimit-affine-formal} 中讨论的情形。
\end{proof}





\section{再次完备化}
\label{formal-spaces-section-recompletion}

\noindent
本节定义沿约化空间的闭子集对形式代数空间进行完备化。
这是第 \ref{formal-spaces-section-completion} 节的自然推广。

\begin{lemma}
\label{formal-spaces-lemma-completion-affine-formal-is-affine-formal}
设 $S$ 为概形，设 $X$ 为 $S$ 上的仿射形式代数空间。
设 $T \subset |X_{red}|$ 为闭子集。则函子
$$
X_{/T} : (\Sch/S)_{fppf} \longrightarrow \textit{Sets},\quad
U \longmapsto \{f : U \to X : f(|U|) \subset T\}
$$
是仿射形式代数空间。
\end{lemma}

\begin{proof}
按定义 \ref{formal-spaces-definition-affine-formal-algebraic-space} 写
$X = \colim X_\lambda$。于是 $X_{\lambda, red} = X_{red}$，
我们可以并且将 $T$ 看作 $|X_\lambda| = |X_{\lambda, red}|$ 的闭子集。
由引理 \ref{formal-spaces-lemma-completion-affine-is-affine-formal-algebraic-space}，
对每个 $\lambda$，完备化 $(X_\lambda)_{/T}$ 是仿射形式代数空间。
过渡态射 $(X_\lambda)_{/T} \to (X_\mu)_{/T}$ 是闭浸入，
因为它们是过渡态射 $X_\lambda \to X_\mu$ 的基变换，
见引理 \ref{formal-spaces-lemma-map-completions-representable}。
此外，由引理 \ref{formal-spaces-lemma-reduction-completion}，
态射 $((X_\lambda)_{/T})_{red} \to ((X_\mu)_{/T})_{red}$
是同构。
由于 $X_{/T} = \colim (X_\lambda)_{/T}$，由引理
\ref{formal-spaces-lemma-colimit-affine-formal} 即得结论。
\end{proof}

\begin{lemma}
\label{formal-spaces-lemma-completion-fas-is-fas}
设 $S$ 为概形，设 $X$ 为 $S$ 上的形式代数空间。
设 $T \subset |X_{red}|$ 为闭子集。则函子
$$
X_{/T} : (\Sch/S)_{fppf} \longrightarrow \textit{Sets},\quad
U \longmapsto \{f : U \to X \mid f(|U|) \subset T\}
$$
是形式代数空间。
\end{lemma}

\begin{proof}
函子 $X_{/T}$ 是 fppf 层，因为若 $\{U_i \to U\}$ 是 fppf 覆盖，
则 $\coprod |U_i| \to |U|$ 是满射。

\medskip\noindent
按定义 \ref{formal-spaces-definition-formal-algebraic-space} 选取覆盖
$\{g_i : X_i \to X\}_{i \in I}$。
态射 $X_i \times_X X_{/T} \to X_{/T}$ 是 \'etale 的
（见《空间》之引理
\ref{spaces-lemma-base-change-representable-transformations-property}），
并且映射 $\coprod X_i \times_X X_{/T} \to X_{/T}$ 是层的满射。
因此只需证明 $X_{/T} \times_X X_i$ 是仿射形式代数空间。
$X_i \times_X X_{/T}$ 的一个 $U$-值点是态射 $U \to X_i$，
其像包含于闭子集
$|g_{i, red}|^{-1}(T) \subset |X_{i, red}|$。
于是结论由引理
\ref{formal-spaces-lemma-completion-affine-formal-is-affine-formal} 得到。
\end{proof}

\begin{definition}
\label{formal-spaces-definition-completion-formal-algebraic-space}
设 $S$ 为概形，设 $X$ 为 $S$ 上的形式代数空间。
设 $T \subset |X_{red}|$ 为闭子集。引理
\ref{formal-spaces-lemma-completion-is-formal-algebraic-space} 中的形式代数空间
$X_{/T}$ 称为沿 $T$ 的 $X$ 的{\it 完备化}。
\end{definition}

\noindent
设 $f : X \to X'$ 是概形 $S$ 上形式代数空间的态射。
设 $T \subset |X_{red}|$ 与 $T' \subset |X'_{red}|$ 为闭子集，
且 $|f_{red}|(T) \subset T'$。显然，$f$ 定义了完备化之间的形式代数空间态射
$$
X_{/T} \longrightarrow X'_{/T'}
$$

\begin{lemma}
\label{formal-spaces-lemma-map-recompletions}
设 $S$ 为概形，设 $f : X' \to X$ 是
$S$ 上形式代数空间的态射。设 $T \subset |X_{red}|$
为闭子集，并令 $T' = |f_{red}|^{-1}(T) \subset |X'_{red}|$。
则
$$
\xymatrix{
X'_{/T'} \ar[r] \ar[d] & X' \ar[d]^f \\
X_{/T} \ar[r] & X
}
$$
是 $S$ 上形式代数空间的笛卡尔图。
\end{lemma}

\begin{proof}
事实上，按构造水平箭头是单态射。因此只需证明：
概形 $U$ 到 $X'$ 的态射 $g : U \to X'$ 定义
$X'_{/T}$ 中的点，当且仅当 $f \circ g$ 定义 $X_{/T}$ 中的点。
换言之，我们需证明
% P10-E0377
$g(U)$ 包含于 $T' \subset |X'_{red}|$
当且仅当 $(f \circ g)(U)$ 包含于 $T \subset |X_{red}|$。
由我们对 $T'$ 取为 $T$ 的逆像，立即得到这一点。
\end{proof}

\begin{lemma}
\label{formal-spaces-lemma-reduction-recompletion}
设 $S$ 为概形，设 $X$ 为 $S$ 上的形式代数空间。
设 $T \subset |X_{red}|$ 为闭子集。完备化 $X_{/T}$ 沿 $T$ 的约化
$(X_{/T})_{red}$ 是 $X_{red}$ 的约化诱导闭子空间 $Z$，
它对应于 $T$。
\end{lemma}

\begin{proof}
这由引理 \ref{formal-spaces-lemma-reduction-formal-algebraic-space}、
《空间的性质》之定义
\ref{spaces-properties-definition-reduced-induced-space}
（该定义使用《空间的性质》之引理
\ref{spaces-properties-lemma-reduced-closed-subspace} 构造 $Z$）
以及 $X_{/T}$ 的定义得到：$Z$ 与 $(X_{/T})_{red}$ 都是约化代数空间，
并由同一映射性质刻画：
% P10-E0378
源为约化代数空间的态射 $g : Y \to X$ 可分解通过它们，
当且仅当 $|Y|$ 映入 $T \subset |X|$。
\end{proof}

\begin{lemma}
\label{formal-spaces-lemma-recompletion-affine-types}
设 $S$ 为概形，设 $X$ 为 $S$ 上的仿射形式代数空间。
设 $T \subset X_{red}$ 为闭子集，设 $X_{/T}$ 是沿 $T$ 的 $X$ 的形式完备化。
则
\begin{enumerate}
\item $X_{/T}$ 是仿射形式代数空间，
\item 若 $X$ 是 McQuillan，则 $X_{/T}$ 是 McQuillan，
\item 若 $|X_{red}| \setminus T$ 是拟紧的且 $X$ 是可数指标的，
则 $X_{/T}$ 是可数指标的，
\item 若 $|X_{red}| \setminus T$ 是拟紧的且 $X$ 是 adic* 的，
则 $X_{/T}$ 是 adic* 的，
\item 若 $X$ 是诺特的，则 $X_{/T}$ 是诺特的。
\end{enumerate}
\end{lemma}

\begin{proof}
(1) 是引理 \ref{formal-spaces-lemma-completion-affine-formal-is-affine-formal}。
若 $X$ 是 McQuillan，则对某个弱可容许拓扑环 $A$ 有
$X = \text{Spf}(A)$。于是
$X_{/T} \to X \to \Spec(A)$ 满足引理
\ref{formal-spaces-lemma-mcquillan-affine-formal-algebraic-space} 的性质 (2)，
因此 $X_{/T}$ 是 McQuillan，见定义
\ref{formal-spaces-definition-types-affine-formal-algebraic-space}。

\medskip\noindent
假设 $X$ 与 $T$ 如 (3) 中所述。
则 $X = \text{Spf}(A)$，其中 $A$ 有弱定义理想的基本系
$A \supset I_1 \supset I_2 \supset I_3 \supset \ldots$，
见引理 \ref{formal-spaces-lemma-countably-indexed}。
由《代数》之引理 \ref{algebra-lemma-qc-open}，
可找到有限生成理想
$\overline{J} = (\overline{f}_1, \ldots, \overline{f}_r) \subset A/I_1$，
使得 $T$ 在 $\Spec(A/I_1) = |X_{red}|$ 内由 $\overline{J}$ 定义。
取 $f_i \in A$ 为 $\overline{f}_i$ 的提升。
若 $Z = \Spec(B)$ 是仿射概形，且 $g : Z \to X$ 是
（集合论意义下）像包含于 $T$ 的态射，则
$g^\sharp : A \to B$ 对某个 $n$ 经由 $A/I_n$ 分解，
并且对每个 $i$，$g^\sharp(f_i)$ 在 $B$ 中幂零。
因此，对于某些 $n, m \geq 1$，理想
$J_{m, n} = (f_1, \ldots, f_r)^m + I_n$ 映到 $B$ 中为零。
由此 $X_{/T}$ 是
$\lim_{n, m} A/J_{m, n}$ 的形式谱，因而是可数指标的。
这证明了 (3)。

\medskip\noindent
(4) 的证明。这里论证与 (3) 相同。
不过此时可以取 $I_n = I^n$，其中 $I \subset A$ 是有限生成理想。
于是显然 $X_{/T}$ 是 $\lim A/J^n$ 的形式谱，
其中 $J = (f_1, \ldots, f_r) + I$。细节略去。

\medskip\noindent
(5) 的证明。在此情形，$X_{red}$ 是诺特环的谱，
因而满足 $|X_{red}| \setminus T$ 拟紧的假设。
因此如 (4) 的证明，我们看到 $X_{/T}$ 是 $\lim A/J^n$ 的谱，
这是诺特的 adic 拓扑环，见《代数》之引理
\ref{algebra-lemma-completion-Noetherian-Noetherian}。
\end{proof}

\begin{lemma}
\label{formal-spaces-lemma-recompletion-types}
设 $S$ 为概形，设 $X$ 为 $S$ 上的形式代数空间。
设 $T \subset X_{red}$ 为闭子集，设 $X_{/T}$ 是沿 $T$ 的 $X$ 的形式完备化。
则
\begin{enumerate}
\item 若 $X_{red} \setminus T \to X_{red}$ 是拟紧的且 $X$ 局部可数指标，
则 $X_{/T}$ 局部可数指标，
\item 若 $X_{red} \setminus T \to X_{red}$ 是拟紧的且 $X$ 局部 adic*，
则 $X_{/T}$ 局部 adic*，并且
\item 若 $X$ 局部诺特，则 $X_{/T}$ 局部诺特。
\end{enumerate}
\end{lemma}

\begin{proof}
按定义 \ref{formal-spaces-definition-formal-algebraic-space} 选取覆盖
$\{X_i \to X\}$。令 $T_i \subset X_{i, red}$ 为 $T$ 的逆像。
由引理 \ref{formal-spaces-lemma-map-recompletions}，有
$X_i \times_X X_{/T} = (X_i)_{/T_i}$。
因此 $\{(X_i)_{/T_i} \to X_{/T}\}$ 是
满足定义 \ref{formal-spaces-definition-formal-algebraic-space} 的覆盖。
此外，若 $X_{red} \setminus T \to X_{red}$ 是拟紧的，
则 $X_{i, red} \setminus T_i \to X_{i, red}$ 也是拟紧的；
若 $X$ 局部可数指标、局部 adic* 或局部诺特，则
% P10-E0379
$X_i$ 分别是可数指标、adic* 或诺特的。
因此引理由仿射情形（即引理
\ref{formal-spaces-lemma-recompletion-affine-types}）得到。
\end{proof}







\section{沿闭子空间的完备化}
\label{formal-spaces-section-completion-subspace}

\noindent
% P10-E0380
本节是关于相对于闭子空间的完备化的第
\ref{formal-spaces-section-completion} 节的类似物。

\begin{definition}
\label{formal-spaces-definition-completion-subspace}
设 $S$ 为概形，设 $X$ 为 $S$ 上的代数空间。
设 $Z \subset X$ 为闭子空间，并记 $Z_n \subset X$ 为 $n$ 阶无穷小邻域。
形式代数空间
$$
X^\wedge_Z = \colim Z_n
$$
（见引理 \ref{formal-spaces-lemma-colimit-formal-spaces-is-formal-space}）
称为沿 $Z$ 的 $X$ 的{\it 完备化}。
\end{definition}

\noindent
例如，若 $X = \Spec(A)$ 且 $Z$ 由理想 $I \subset A$ 定义，
则有
$$
X^\wedge_Z = \colim \Spec(A/I^n) = \text{Spf}(B)
$$
其中 $B = \lim A/I^n$，极限取极限拓扑。注意，$B$ 是
弱 adic 拓扑环，但一般不是 adic 的。

\medskip\noindent
回到一般情形，若 $T = |Z|$，则存在典范态射
$X^\wedge_Z \to X_{/T}$，比较沿 $Z$ 与沿 $T$ 的完备化
（第 \ref{formal-spaces-section-completion} 节），它不一定是同构
（见引理 \ref{formal-spaces-lemma-when-isomorphism}）。

\medskip\noindent
设 $f : X \to X'$ 是概形 $S$ 上代数空间的态射。
假设 $Z \subset X$ 与 $Z' \subset X'$ 是闭子空间，
且 $f|_Z$ 将 $Z$ 映入 $Z'$ 并诱导态射 $Z \to Z'$。
显然，$f$ 定义形式代数空间的态射
$$
X^\wedge_Z \longrightarrow (X')^\wedge_{Z'}
$$
介于这些完备化之间。

\begin{lemma}
\label{formal-spaces-lemma-map-completions-subspaces-representable}
设 $S$ 为概形，设 $f : X' \to X$ 是 $S$ 上代数空间的态射。
设 $Z \subset X$ 为闭子空间，并令 $Z' = f^{-1}(Z) = X' \times_X Z$。
则
$$
\xymatrix{
(X')^\wedge_{Z'} \ar[r] \ar[d] & X' \ar[d]^f \\
X^\wedge_Z \ar[r] & X
}
$$
是层的笛卡尔图。特别地，态射
$(X')^\wedge_{Z'} \to X^\wedge_Z$ 可由代数空间表示。
\end{lemma}

\begin{proof}
事实上，设 $Y \to X$ 是从概形到 $X$ 的态射，并且
$Y \to X$ 分解通过 $Z$。则 $Y \times_X X' \to X$
是代数空间的态射，且 $Y \times_X X' \to X'$
分解通过 $Z'$。由于对所有 $n \geq 1$ 都有
$Z'_n = X' \times_X Z_n$，无穷小邻域也具有同样性质。
因此函子的笛卡尔方图由公式
$X^\wedge_Z = \colim Z_n$ 以及
$(X')^\wedge_{Z'} = \colim Z'_n$ 得到。
\end{proof}

\begin{lemma}
\label{formal-spaces-lemma-reduction-completion-subspace}
设 $S$ 为概形，设 $X$ 为 $S$ 上的代数空间。
设 $Z \subset X$ 为闭子空间。完备化 $X^\wedge_Z$ 沿 $Z$ 的约化
$(X^\wedge_Z)_{red}$ 是 $Z_{red}$。
\end{lemma}

\begin{proof}
略。
\end{proof}

\begin{lemma}
\label{formal-spaces-lemma-affine-formal-completion-subspace-types}
设 $S$ 为概形，设 $X = \Spec(A)$ 是 $S$ 上的仿射概形。
设 $Z \subset X$ 是对应于理想 $I \subset A$ 的闭子概形。则
\begin{enumerate}
\item 仿射形式代数空间 $X^\wedge_Z$ 是弱 adic 的，
\item 若 $I$ 有限生成，则 $X^\wedge_Z = \text{Spf}(A^\wedge)$，
其中 $A^\wedge$ 是 $A$ 的 $I$-进完备化，
\item 若 $Z \to X$ 是有限呈示的，即 $I$ 有限生成，
则 $X^\wedge_Z$ 是 adic* 的，
\item 若 $X$ 是诺特的，则 $X^\wedge_Z$ 是诺特的。
\end{enumerate}
\end{lemma}

\begin{proof}
略。
\end{proof}

\begin{lemma}
\label{formal-spaces-lemma-formal-completion-subspace-types}
设 $S$ 为概形，设 $X$ 为 $S$ 上的代数空间。
设 $Z \subset X$ 是闭子空间，设 $X^\wedge_Z$ 是沿 $Z$ 的 $X$ 的形式完备化。
\begin{enumerate}
\item 形式代数空间 $X^\wedge_Z$ 是局部弱 adic 的，
\item 若 $Z \to X$ 是有限呈示的，则 $X^\wedge_Z$ 是局部 adic* 的，
\item 若 $X$ 局部诺特，则
% P10-E0381
$X_Z$ 是局部诺特的。
\end{enumerate}
\end{lemma}

\begin{proof}
略。
\end{proof}

\begin{lemma}
\label{formal-spaces-lemma-when-isomorphism}
设 $S$ 为概形，设 $X$ 为 $S$ 上的代数空间。
设 $Z \subset X$ 是闭子空间，令 $T = |Z|$。
典范态射 $c : X^\wedge_Z \to X_{/T}$ 是同构，当 $Z \to X$
有限呈示时成立，但一般不成立。
\end{lemma}

\begin{proof}
两种构造都与 \'etale 局部化交换，因此只需在 $X$ 仿射时证明。
设 $X = \Spec(A)$，且 $Z$ 对应于理想 $I \subset A$。
若 $I$ 有限生成，则 $X^\wedge_Z$ 与 $X_{/T}$ 都等于
$\text{Spf}(A^\wedge)$，其中 $A^\wedge$ 是 $A$ 的 $I$-进完备化，
见引理 \ref{formal-spaces-lemma-affine-formal-completion-types} 和
\ref{formal-spaces-lemma-affine-formal-completion-subspace-types}。
若 $A = \mathbf{Z}[x_1, x_2, \ldots]$ 且
$I = (x_1, x_2, \ldots)$，则
$\Spec(A/(x_n^n, n \geq 1)) \to X$ 经由 $X_{/T}$ 分解，
但不经由 $X^\wedge_Z$ 分解，因此 $c$ 不是同构。
\end{proof}
